\begin{filecontents*}[overwrite]{\jobname.ist}
headings_flag 1
numhead_positive "数字"
symhead_positive "符号"
heading_prefix "{\\bfseries\\hfil "
heading_suffix "\\hfil}\\nopagebreak\\par\\vskip 5pt "
\end{filecontents*}

\setCJKmainfont[
  ItalicFont     = KaiTi,
  BoldFont       = SimHei,
  BoldItalicFont = FangSong,
]{SimSun}

\usepackage{amsmath,amssymb,amsthm}
\usepackage{mathtools}
\usepackage{tikz}              %
\usepackage{etoolbox}          %
\usepackage[a4paper,margin=2.8cm]{geometry}
\usepackage{enumitem}
\setlist{itemsep=2pt,topsep=4pt}
\emergencystretch=1.5em %
\vfuzz=3pt              %
\hbadness=2500          %
\xeCJKsetup{PunctStyle=quanjiao}%


\newif\ifzhnames
\zhnamestrue
\newcommand{\trn}[2]{\ifzhnames #2\else #1\fi}

\newcommand{\trhAxiom}{公理}%
\newcommand{\trhTheorem}{定理}%
\newcommand{\trhProposition}{命题}%
\newcommand{\trhLemma}{引理}%
\newcommand{\trhCorollary}{推论}%
\newcommand{\trhDefinition}{定义}%
\newcommand{\trhExample}{例}%
\newcommand{\trhRemark}{注记}%
\newcommand{\trhProof}{证明}%

\newcommand{\trHofer}{\trn{Hofer}{霍弗}}%
\newcommand{\trZehnder}{\trn{Zehnder}{曾德尔}}%
\newcommand{\trHamilton}{\trn{Hamilton}{哈密顿}}%
\newcommand{\trPoincare}{\trn{Poincar\'e}{庞加莱}}%
\newcommand{\trArnold}{\trn{Arnold}{阿诺德}}%
\newcommand{\trMoser}{\trn{Moser}{莫泽}}%
\newcommand{\trDarboux}{\trn{Darboux}{达布}}%
\newcommand{\trGromov}{\trn{Gromov}{格罗莫夫}}%
\newcommand{\trRabinowitz}{\trn{Rabinowitz}{拉比诺维茨}}%
\newcommand{\trFloer}{\trn{Floer}{Floer}}%
\newcommand{\trWeinstein}{\trn{Weinstein}{温斯坦}}%
\newcommand{\trLiouville}{\trn{Liouville}{刘维尔}}%
\newcommand{\trBirkhoff}{\trn{Birkhoff}{伯克霍夫}}%
\newcommand{\trLyapunov}{\trn{Lyapunov}{李雅普诺夫}}%
\newcommand{\trLjusternik}{\trn{Ljusternik}{柳斯捷尔尼克}}%
\newcommand{\trSchnirelman}{\trn{Schnirelman}{施尼雷尔曼}}%
\newcommand{\trConley}{\trn{Conley}{康利}}%
\newcommand{\trEkeland}{\trn{Ekeland}{埃克朗}}%
\newcommand{\trLasry}{\trn{Lasry}{拉斯里}}%
\newcommand{\trHerman}{\trn{Herman}{赫尔曼}}%
\newcommand{\trDacorogna}{\trn{Dacorogna}{达科格纳}}%
\newcommand{\trKunzle}{\trn{Künzle}{屈恩茨勒}}%
\newcommand{\trCroke}{\trn{Croke}{克罗克}}%
\newcommand{\trMassey}{\trn{Massey}{马西}}%
\newcommand{\trMa}{\trn{Renyi Ma}{Renyi Ma}}%
\newcommand{\trLutz}{\trn{Lutz}{卢茨}}%
\newcommand{\trMartinet}{\trn{Martinet}{马蒂内}}%
\newcommand{\trBennequin}{\trn{Bennequin}{贝内坎}}%
\newcommand{\trYoccoz}{\trn{Yoccoz}{约科兹}}%
\newcommand{\trHadamard}{\trn{Hadamard}{阿达玛}}%
\newcommand{\trFet}{\trn{Fet}{费特}}%
\newcommand{\trBolotin}{\trn{Bolotin}{博洛京}}%
\newcommand{\trGluck}{\trn{Gluck}{格卢克}}%
\newcommand{\trZiller}{\trn{Ziller}{齐勒}}%
\newcommand{\trKlingenberg}{\trn{Klingenberg}{克林根贝格}}%
\newcommand{\trBangert}{\trn{Bangert}{班格特}}%
\newcommand{\trFranks}{\trn{Franks}{弗兰克斯}}%
\newcommand{\trOffin}{\trn{Offin}{奥芬}}%
\newcommand{\trKoslov}{\trn{Koslov}{科兹洛夫}}%
\newcommand{\trMaupertuis}{\trn{Maupertuis}{莫佩蒂}}%
\newcommand{\trJiang}{\trn{Jiang}{蒋美跃}}%
\newcommand{\trCalabi}{\trn{Calabi}{卡拉比}}%
\newcommand{\trLefschetz}{\trn{Lefschetz}{莱夫谢茨}}%
\newcommand{\trHopf}{\trn{Hopf}{霍普夫}}%
\newcommand{\trNikishin}{\trn{Nikishin}{尼基申}}%
\newcommand{\trSimon}{\trn{Simon}{西蒙}}%
\newcommand{\trLoewner}{\trn{Loewner}{勒夫纳}}%
\newcommand{\trLeCalvez}{\trn{Le Calvez}{勒卡尔韦}}%
\newcommand{\trGambaudo}{\trn{Gambaudo}{冈博多}}%
\newcommand{\trFortune}{\trn{Fortune}{福图内}}%
\newcommand{\trBanyaga}{\trn{Banyaga}{巴尼亚加}}%
\newcommand{\trMaslov}{\trn{Maslov}{马斯洛夫}}%
\newcommand{\trCech}{\trn{Čech}{切赫}}%
\newcommand{\trSchmidt}{\trn{Schmidt}{施密特}}%
\newcommand{\trWarner}{\trn{Warner}{华纳}}%
\newcommand{\trBrezis}{\trn{Brezis}{布雷齐斯}}%
\newcommand{\trNirenberg}{\trn{Nirenberg}{尼伦贝格}}%
\newcommand{\trDold}{\trn{Dold}{多尔德}}%
\newcommand{\trGivental}{\trn{Givental}{吉文塔尔}}%
\newcommand{\trBrunella}{\trn{Brunella}{布鲁内拉}}%
\newcommand{\trChekanov}{\trn{Chekanov}{切卡诺夫}}%
\newcommand{\trHormander}{\trn{Hörmander}{赫尔曼德}}%
\newcommand{\trTraynor}{\trn{Traynor}{特雷诺}}%
\newcommand{\trRybakowski}{\trn{Rybakowski}{雷巴科夫斯基}}%
\newcommand{\trWhitney}{\trn{Whitney}{惠特尼}}%
\newcommand{\trFrechet}{\trn{Fréchet}{弗雷歇}}%
\newcommand{\trHirsch}{\trn{Hirsch}{赫希}}%
\newcommand{\trHeine}{\trn{Heine}{海涅}}%
\newcommand{\trStein}{\trn{Stein}{施泰因}}%
\newcommand{\trEliasson}{\trn{Eliasson}{埃利亚松}}%
\newcommand{\trForster}{\trn{Forster}{福斯特}}%
\newcommand{\trLeibniz}{\trn{Leibniz}{莱布尼茨}}%
\newcommand{\trOno}{\trn{Ono}{小野}}%
\newcommand{\trLeHongVan}{\trn{L\^{e} H\^{\`{o}}ng V\^{a}n}{黎红云}}%
\newcommand{\trHurewicz}{\trn{Hurewicz}{胡雷维茨}}%
\newcommand{\trChern}{\trn{Chern}{陈}}%
\newcommand{\trCasson}{\trn{Casson}{卡森}}%
\newcommand{\trSimons}{\trn{Simons}{西蒙斯}}%
\newcommand{\trNovikov}{\trn{Novikov}{诺维科夫}}%
\newcommand{\trFukaya}{\trn{Fukaya}{深谷}}%
\newcommand{\trBourbaki}{\trn{Bourbaki}{布尔巴基}}%
\newcommand{\trGalerkin}{\trn{Galerkin}{伽辽金}}%
\newcommand{\trSchrodinger}{\trn{Schrödinger}{薛定谔}}%
\newcommand{\trBourgain}{\trn{Bourgain}{布尔甘}}%
\newcommand{\trUstilovsky}{\trn{Ustilovsky}{乌斯季洛夫斯基}}%
\newcommand{\trAudin}{\trn{Audin}{奥丹}}%
\newcommand{\trLafontaine}{\trn{Lafontaine}{拉丰丹}}%
\newcommand{\trAebischer}{\trn{Aebischer}{艾比舍尔}}%
\newcommand{\trBorer}{\trn{Borer}{博勒尔}}%
\newcommand{\trLeuenberger}{\trn{Leuenberger}{洛伊恩贝格尔}}%
\newcommand{\trKalin}{\trn{Kälin}{凯林}}%
\newcommand{\trReimann}{\trn{Reimann}{赖曼}}%
\newcommand{\trDeBoor}{\trn{de Boor}{德布尔}}%
\newcommand{\trHolder}{\trn{Hölder}{赫尔德}}%
\newcommand{\trYoung}{\trn{Young}{杨}}%
\newcommand{\trBeurling}{\trn{Beurling}{贝尔林}}%
\newcommand{\trMarshall}{\trn{Marshall}{马歇尔}}%
\newcommand{\trBeckner}{\trn{Beckner}{贝克纳}}%
\newcommand{\trAbraham}{\trn{Abraham}{亚伯拉罕}}%
\newcommand{\trRobbin}{\trn{Robbin}{罗宾}}%
\newcommand{\trNekhoroshev}{\trn{Nekhoroshev}{涅霍罗舍夫}}%
\newcommand{\trMishchenko}{\trn{Mishchenko}{米申科}}%
\newcommand{\trFomenko}{\trn{Fomenko}{福缅科}}%
\newcommand{\trDuistermaat}{\trn{Duistermaat}{杜斯特马特}}%
\newcommand{\trDouglis}{\trn{Douglis}{道格利斯}}%
\newcommand{\trVekua}{\trn{Vekua}{韦夸}}%
\newcommand{\trVick}{\trn{Vick}{维克}}%
\newcommand{\trSwitzer}{\trn{Switzer}{斯威策}}%
\newcommand{\trJordan}{\trn{Jordan}{若尔当}}%
\newcommand{\trLima}{\trn{Lima}{利马}}%
\newcommand{\trMcDuff}{\trn{McDuff}{麦克达夫}}%
\newcommand{\trSalamon}{\trn{Salamon}{萨拉蒙}}%
\newcommand{\trViterbo}{\trn{Viterbo}{维泰博}}%
\newcommand{\trEliashberg}{\trn{Eliashberg}{埃利亚什伯格}}%
\newcommand{\trWitt}{\trn{Witt}{维特}}%
\newcommand{\trJost}{\trn{Jost}{约斯特}}%
\newcommand{\trSobolev}{\trn{Sobolev}{索伯列夫}}%
\newcommand{\trBanach}{\trn{Banach}{巴拿赫}}%
\newcommand{\trHilbert}{\trn{Hilbert}{希尔伯特}}%
\newcommand{\trRiemann}{\trn{Riemann}{黎曼}}%
\newcommand{\trCauchy}{\trn{Cauchy}{柯西}}%
\newcommand{\trSchwarz}{\trn{Schwarz}{施瓦茨}}%
\newcommand{\trLagrange}{\trn{Lagrange}{拉格朗日}}%
\newcommand{\trJacobi}{\trn{Jacobi}{雅可比}}%
\newcommand{\trFourier}{\trn{Fourier}{傅里叶}}%
\newcommand{\trParseval}{\trn{Parseval}{帕塞瓦尔}}%
\newcommand{\trPlancherel}{\trn{Plancherel}{普朗歇尔}}%
\newcommand{\trNeumann}{\trn{Neumann}{诺伊曼}}%
\newcommand{\trDirichlet}{\trn{Dirichlet}{狄利克雷}}%
\newcommand{\trFredholm}{\trn{Fredholm}{弗雷德霍姆}}%
\newcommand{\trBrouwer}{\trn{Brouwer}{布劳威尔}}%
\newcommand{\trAlexander}{\trn{Alexander}{亚历山大}}%
\newcommand{\trSpanier}{\trn{Spanier}{斯帕尼尔}}%
\newcommand{\trDeRham}{\trn{de Rham}{德拉姆}}%
\newcommand{\trHodge}{\trn{Hodge}{霍奇}}%
\newcommand{\trCartan}{\trn{Cartan}{嘉当}}%
\newcommand{\trLie}{\trn{Lie}{李}}%
\newcommand{\trWeyl}{\trn{Weyl}{外尔}}%
\newcommand{\trNoether}{\trn{Noether}{诺特}}%
\newcommand{\trMilnor}{\trn{Milnor}{米尔诺}}%
\newcommand{\trThom}{\trn{Thom}{托姆}}%
\newcommand{\trSard}{\trn{Sard}{萨德}}%
\newcommand{\trBaire}{\trn{Baire}{贝尔}}%
\newcommand{\trLebesgue}{\trn{Lebesgue}{勒贝格}}%
\newcommand{\trBorel}{\trn{Borel}{博雷尔}}%
\newcommand{\trRellich}{\trn{Rellich}{雷利希}}%
\newcommand{\trArzela}{\trn{Arzel\`a}{阿尔泽拉}}%
\newcommand{\trAscoli}{\trn{Ascoli}{阿斯科利}}%
\newcommand{\trPeano}{\trn{Peano}{佩亚诺}}%
\newcommand{\trMorse}{\trn{Morse}{莫尔斯}}%
\newcommand{\trPalais}{\trn{Palais}{帕莱}}%
\newcommand{\trSmale}{\trn{Smale}{斯梅尔}}%
\newcommand{\trKrasnoselskii}{\trn{Krasnoselskii}{克拉斯诺谢利斯基}}%
\newcommand{\trLaplace}{\trn{Laplace}{拉普拉斯}}%
\newcommand{\trTaubes}{\trn{Taubes}{陶布斯}}%
\newcommand{\trLalonde}{\trn{Lalonde}{拉隆德}}%
\newcommand{\trPolterovich}{\trn{Polterovich}{波尔泰罗维奇}}%
\newcommand{\trSikorav}{\trn{Sikorav}{西科拉夫}}%
\newcommand{\trGinzburg}{\trn{Ginzburg}{金茨堡}}%
\newcommand{\trStruwe}{\trn{Struwe}{施特鲁韦}}%
\newcommand{\trAmbrosetti}{\trn{Ambrosetti}{安布罗塞蒂}}%
\newcommand{\trClarke}{\trn{Clarke}{克拉克}}%
\newcommand{\trBenci}{\trn{Benci}{本奇}}%
\newcommand{\trChaperon}{\trn{Chaperon}{沙普隆}}%
\newcommand{\trKuksin}{\trn{Kuksin}{库克辛}}%
\newcommand{\trWysocki}{\trn{Wysocki}{维索茨基}}%
\newcommand{\trBialy}{\trn{Bialy}{比亚雷}}%
\newcommand{\trCieliebak}{\trn{Cieliebak}{齐利巴克}}%
\newcommand{\trFlucher}{\trn{Flucher}{弗卢歇尔}}%
\newcommand{\trGoing}{\trn{G\"oing}{格因}}%
\newcommand{\trGratza}{\trn{Gratza}{格拉察}}%
\newcommand{\trKaas}{\trn{Kaas}{卡斯}}%
\newcommand{\trKriener}{\trn{Kriener}{克里纳}}%
\newcommand{\trSiburg}{\trn{Siburg}{齐堡}}%
\newcommand{\trMcLachlan}{\trn{McLachlan}{麦克拉克伦}}%
\newcommand{\trFlepp}{\trn{Flepp}{弗莱普}}%
\newcommand{\trCarleman}{\trn{Carleman}{卡莱曼}}%
\newcommand{\trLeray}{\trn{Leray}{勒雷}}%
\newcommand{\trSchauder}{\trn{Schauder}{绍德尔}}%
\newcommand{\trAbbas}{\trn{Abbas}{阿巴斯}}%
\newcommand{\trNash}{\trn{Nash}{纳什}}%
\newcommand{\trRiesz}{\trn{Riesz}{里斯}}%
\newcommand{\trLipschitz}{\trn{Lipschitz}{利普希茨}}%
\newcommand{\trTerng}{\trn{Terng}{滕楚莲}}%
\newcommand{\trBlanchard}{\trn{Blanchard}{布兰查德}}%
\newcommand{\trBruning}{\trn{Br\"uning}{布吕宁}}%
\newcommand{\trMawhin}{\trn{Mawhin}{莫万}}%
\newcommand{\trWillem}{\trn{Willem}{维勒姆}}%
\newcommand{\trChang}{\trn{Chang}{张恭庆}}%
\newcommand{\trChangSA}{\trn{Chang}{张圣容}}%
\newcommand{\trStokes}{\trn{Stokes}{斯托克斯}}%
\newcommand{\trGauss}{\trn{Gauss}{高斯}}%
\newcommand{\trPugh}{\trn{Pugh}{皮尤}}%
\newcommand{\trRobinson}{\trn{Robinson}{罗宾逊}}%
\newcommand{\trSeifert}{\trn{Seifert}{塞弗特}}%
\newcommand{\trSchweitzer}{\trn{Schweitzer}{施魏策尔}}%
\newcommand{\trKuperberg}{\trn{Kuperberg}{库珀伯格}}%
\newcommand{\trWilson}{\trn{Wilson}{威尔逊}}%
\newcommand{\trReeb}{\trn{Reeb}{Reeb}}%
\newcommand{\trEuler}{\trn{Euler}{欧拉}}%
\newcommand{\trLegendre}{\trn{Legendre}{勒让德}}%
\newcommand{\trWeierstrass}{\trn{Weierstrass}{维尔斯特拉斯}}%
\newcommand{\trCourant}{\trn{Courant}{库朗}}%
\newcommand{\trPoisson}{\trn{Poisson}{泊松}}%
\newcommand{\trTaylor}{\trn{Taylor}{泰勒}}%
\newcommand{\trSiegel}{\trn{Siegel}{西格尔}}%
\newcommand{\trLewis}{\trn{Lewis}{刘易斯}}%
\newcommand{\trAvez}{\trn{Avez}{阿韦兹}}%
\newcommand{\trKolmogorov}{\trn{Kolmogorov}{柯尔莫哥洛夫}}%
\newcommand{\trBernoulli}{\trn{Bernoulli}{伯努利}}%
\newcommand{\trKronecker}{\trn{Kronecker}{克罗内克}}%
\newcommand{\trFloquet}{\trn{Floquet}{弗洛凯}}%
\newcommand{\trMather}{\trn{Mather}{马瑟}}%
\newcommand{\trCantor}{\trn{Cantor}{康托尔}}%
\newcommand{\trPoschel}{\trn{P\"oschel}{波歇尔}}%
\newcommand{\trGenecand}{\trn{Genecand}{热内坎}}%
\newcommand{\trChenciner}{\trn{Chenciner}{尚西纳}}%
\newcommand{\trRiemannin}{\trn{Riemann}{黎曼}}%
\newcommand{\trAmann}{\trn{Amann}{阿曼}}%
\newcommand{\trDeimling}{\trn{Deimling}{戴姆林}}%
\newcommand{\trZeidler}{\trn{Zeidler}{蔡德勒}}%

\newcommand{\trSymplectic}{辛}%
\newcommand{\trAntisymplectic}{反辛}%
\newcommand{\trSymplecticGeometry}{辛几何}%
\newcommand{\trSymplecticTopology}{辛拓扑}%
\newcommand{\trSymplecticManifold}{辛流形}%
\newcommand{\trSymplecticVectorSpace}{辛向量空间}%
\newcommand{\trSymplecticForm}{辛形式}%
\newcommand{\trSymplecticStructure}{辛结构}%
\newcommand{\trSymplecticGroup}{辛群}%
\newcommand{\trSymplecticMatrix}{\trSymplectic{}矩阵}%
\newcommand{\trGroupAction}{群作用}%
\newcommand{\trAbelian}{交换}%
\newcommand{\trMapping}{映射}%
\newcommand{\trIdentityMap}{恒同映射}%
\newcommand{\trSymplecticMapping}{辛映射}%
\newcommand{\trSymplecticDiffeomorphism}{辛微分同胚}%
\newcommand{\trSymplectomorphism}{辛微分同胚}%
\newcommand{\trSymplecticHomeomorphism}{辛同胚}%
\newcommand{\trSymplecticEmbedding}{辛嵌入}%
\newcommand{\trSymplecticCapacity}{辛容量}%
\newcommand{\trSymplecticInvariant}{辛不变量}%
\newcommand{\trSymplecticBasis}{辛基}%
\newcommand{\trSymplecticCoordinates}{辛坐标}%
\newcommand{\trSymplecticComplement}{辛补}%
\newcommand{\trSymplecticLinear}{辛线性}%
\newcommand{\trSymplecticIsomorphism}{辛同构}%
\newcommand{\trCapacity}{容量}%
\newcommand{\trInnerCapacity}{内容量}%
\newcommand{\trOuterCapacity}{外容量}%
\newcommand{\trInnerRegularity}{内正则性}%
\newcommand{\trOuterRegularity}{外正则性}%
\newcommand{\trNontriviality}{非平凡性}%
\newcommand{\trInclusionMapping}{包含映射}%
\newcommand{\trUniformization}{单值化}%
\newcommand{\trGromovWidth}{\trGromov{}宽度}%
\newcommand{\trSqueezingTheorem}{挤压定理}%
\newcommand{\trHamiltonian}{哈密顿}%
\newcommand{\trHamiltonianVectorField}{\trHamiltonian{}向量场}%
\newcommand{\trHamiltonianSystem}{\trHamiltonian{}系统}%
\newcommand{\trHamiltonianEquation}{\trHamiltonian{}方程}%
\newcommand{\trHamiltonianFunction}{\trHamiltonian{}函数}%
\newcommand{\trHamiltonianDynamics}{\trHamiltonian{}动力学}%
\newcommand{\trHamiltonianMap}{\trHamiltonian{}映射}%
\newcommand{\trHamiltonianDiffeomorphism}{\trHamiltonian{}微分同胚}%
\newcommand{\trHamiltonianFlow}{\trHamiltonian{}流}%
\newcommand{\trHamiltonianFormalism}{\trHamiltonian{}形式体系}%
\newcommand{\trAlmostComplex}{近复}%
\newcommand{\trAlmostComplexStructure}{近复结构}%
\newcommand{\trComplexStructure}{复结构}%
\newcommand{\trCompatible}{相容}%
\newcommand{\trTamed}{驯服}%
\newcommand{\trContact}{切触}%
\newcommand{\trContactStructure}{切触结构}%
\newcommand{\trIsotropic}{迷向}%
\newcommand{\trFoliation}{叶状结构}%
\newcommand{\trLeaf}{叶}%
\newcommand{\trConjugate}{共轭}%
\newcommand{\trContactForm}{切触形式}%
\newcommand{\trContactType}{切触型}%
\newcommand{\trCharacteristic}{特征}%
\newcommand{\trClosedCharacteristic}{闭特征}%
\newcommand{\trCharacteristicLineBundle}{特征线丛}%
\newcommand{\trEnergy}{能量}%
\newcommand{\trEnergySurface}{能量面}%
\newcommand{\trEnergyLevel}{能级}%
\newcommand{\trCanonical}{正则}%
\newcommand{\trRegular}{正则}%
\newcommand{\trCanonicalTransformation}{正则变换}%
\newcommand{\trCanonicalCoordinates}{正则坐标}%
\newcommand{\trPolarCoordinates}{极坐标}%
\newcommand{\trChart}{坐标卡}%
\newcommand{\trCube}{立方体}%
\newcommand{\trActionAngle}{作用量-角}%
\newcommand{\trGeneratingFunction}{生成函数}%
\newcommand{\trAction}{作用量}%
\newcommand{\trActionFunctional}{作用量泛函}%
\newcommand{\trReducedAction}{约化作用量}%
\newcommand{\trActionSpectrum}{作用量谱}%
\newcommand{\trLeastActionPrinciple}{最小作用量原理}%
\newcommand{\trDisplacementEnergy}{位移能量}%
\newcommand{\trNormalForm}{范式}%
\newcommand{\trBirkhoffInvariant}{\trBirkhoff{}不变量}%
\newcommand{\trNonresonance}{非共振}%
\newcommand{\trResonance}{共振}%
\newcommand{\trQuadraticForm}{二次型}%
\newcommand{\trPositiveDefinite}{正定}%
\newcommand{\trEllipsoid}{椭球}%
\newcommand{\trCylinder}{柱体}%
\newcommand{\trTorus}{环面}%
\newcommand{\trCircle}{圆周}%
\newcommand{\trAnnulus}{环形区域}%
\newcommand{\trBall}{球}%
\newcommand{\trUnitBall}{单位球}%
\newcommand{\trDisc}{圆盘}%
\newcommand{\trVolume}{体积}%
\newcommand{\trVolumeForm}{体积形式}%
\newcommand{\trVolumePreserving}{保体积}%
\newcommand{\trLiouvilleMeasure}{\trLiouville{}测度}%
\newcommand{\trSqueeze}{挤压}%
\newcommand{\trRigidity}{刚性}%

\newcommand{\trDynamicalSystem}{动力系统}%
\newcommand{\trFlow}{流}%
\newcommand{\trPoissonBracket}{\trPoisson{}括号}%
\newcommand{\trInvolution}{对合}%
\newcommand{\trIntegral}{积分}%
\newcommand{\trFirstIntegral}{第一积分}%
\newcommand{\trFrequencyVector}{频率向量}%
\newcommand{\trAccumulationPoint}{聚点}%
\newcommand{\trHyperbolic}{双曲}%
\newcommand{\trDivergent}{发散}%
\newcommand{\trDivergence}{发散性}%
\newcommand{\trOrbit}{轨道}%
\newcommand{\trTrajectory}{轨线}%
\newcommand{\trPeriodic}{周期}%
\newcommand{\trPeriodicOrbit}{周期轨道}%
\newcommand{\trQuasiPeriodic}{拟周期}%
\newcommand{\trPeriodicSolution}{周期解}%
\newcommand{\trFixedPoint}{不动点}%
\newcommand{\trEquilibrium}{平衡点}%
\newcommand{\trRecurrence}{回归}%
\newcommand{\trReturnMap}{返回映射}%
\newcommand{\trInvariant}{不变}%
\newcommand{\trInvariantMeasure}{不变测度}%
\newcommand{\trMeasure}{测度}%
\newcommand{\trMeasurePreserving}{保测度}%
\newcommand{\trMeasureZero}{零测}%
\newcommand{\trStable}{稳定}%
\newcommand{\trUnstable}{不稳定}%
\newcommand{\trStableManifold}{稳定流形}%
\newcommand{\trUnstableManifold}{不稳定流形}%
\newcommand{\trConnectingOrbit}{连接轨道}%
\newcommand{\trHeteroclinic}{异宿}%
\newcommand{\trHomoclinic}{同宿}%
\newcommand{\trClosing}{闭合}%
\newcommand{\trClosingLemma}{闭合引理}%
\newcommand{\trNonClosingLemma}{非闭合引理}%
\newcommand{\trClosingProblem}{闭合问题}%
\newcommand{\trDiophantine}{丢番图}%
\newcommand{\trIrrational}{无理}%
\newcommand{\trTranslation}{平移}%
\newcommand{\trDynamics}{动力学}%
\newcommand{\trCompletion}{完备化}%
\newcommand{\trCommutator}{交换子}%
\newcommand{\trPseudoMetric}{伪度量}%
\newcommand{\trTopologicalGroup}{拓扑群}%
\newcommand{\trOrderStructure}{序结构}%
\newcommand{\trPartialOrder}{偏序}%
\newcommand{\trBijection}{双射}%
\newcommand{\trPrimitive}{原形式}%
\newcommand{\trIntegrationByParts}{分部积分}%
\newcommand{\trBumpFunction}{隆起函数}%
\newcommand{\trMonotone}{单调}%
\newcommand{\trUniversal}{万有}%
\newcommand{\trSpecialCapacity}{特殊容量}%
\newcommand{\trDisjoint}{不相交}%
\newcommand{\trProperlySeparated}{正常分离}%
\newcommand{\trConvexHull}{凸包}%
\newcommand{\trQuasiAutonomous}{拟自治}%
\newcommand{\trMinimal}{极小}%
\newcommand{\trSimple}{简单}%
\newcommand{\trSimplicity}{简单性}%
\newcommand{\trCupLength}{杯长}%
\newcommand{\trMonkeySaddle}{猴鞍}%
\newcommand{\trCenterOfMass}{质心}%
\newcommand{\trCenterOfGravity}{重心}%
\newcommand{\trSectionalCurvature}{截曲率}%
\newcommand{\trCapProduct}{帽积}%
\newcommand{\trProperMap}{恰当映射}%
\newcommand{\trZeroSection}{零截面}%
\newcommand{\trFundamentalClass}{基本类}%
\newcommand{\trReduction}{约化}%
\newcommand{\trPhaseFunction}{相位函数}%
\newcommand{\trFourierIntegralOperator}{傅里叶积分算子}%
\newcommand{\trSplitting}{分裂}%
\newcommand{\trEquicontinuous}{等度连续}%
\newcommand{\trPrecompact}{预紧}%
\newcommand{\trLimitSet}{极限集}%
\newcommand{\trSemiflow}{半流}%
\newcommand{\trContinuum}{连续统}%
\newcommand{\trPeriodGroup}{周期群}%
\newcommand{\trCapDisk}{帽片}%
\newcommand{\trGluingMethod}{粘合方法}%
\newcommand{\trBrokenTrajectory}{破裂轨迹}%
\newcommand{\trConnectedSimpleSystem}{连通单系}%
\newcommand{\trModuliSpace}{模空间}%
\newcommand{\trTrivialization}{平凡化}%
\newcommand{\trInstantonHomology}{瞬子同调}%
\newcommand{\trConnection}{联络}%
\newcommand{\trWindingNumber}{环绕数}%
\newcommand{\trChainMap}{链同态}%
\newcommand{\trChainHomotopyEquivalent}{链同伦等价}%
\newcommand{\trIntersection}{相交}%
\newcommand{\trCoefficientField}{系数域}%
\newcommand{\trAlgebraicTopology}{代数拓扑}%
\newcommand{\trFiltration}{滤过}%
\newcommand{\trDirectedSet}{有向集}%
\newcommand{\trDirectedLimit}{直极限}%
\newcommand{\trPseudoConvex}{伪凸}%
\newcommand{\trPolydisc}{多圆盘}%
\newcommand{\trPacking}{填装}%
\newcommand{\trFlat}{平坦}%
\newcommand{\trBifurcationDiagram}{分岔图}%
\newcommand{\trTautologicalForm}{重言形式}%
\newcommand{\trClosedForm}{闭形式}%
\newcommand{\trBrakeOrbit}{刹车轨道}%
\newcommand{\trKineticEnergy}{动能}%
\newcommand{\trPotentialEnergy}{势能}%
\newcommand{\trConservativeSystem}{保守系统}%
\newcommand{\trDegreesOfFreedom}{自由度}%
\newcommand{\trAmplitude}{振幅}%
\newcommand{\trTensorProduct}{张量积}%
\newcommand{\trSymplecticSplitting}{辛分裂}%
\newcommand{\trHomeomorphic}{同胚}%
\newcommand{\trCell}{胞腔}%
\newcommand{\trCelestialMechanics}{天体力学}%
\newcommand{\trTwoBodyProblem}{二体问题}%
\newcommand{\trThreeBodyProblem}{三体问题}%
\newcommand{\trRestrictedThreeBodyProblem}{限制性三体问题}%
\newcommand{\trCollision}{碰撞}%
\newcommand{\trSingularity}{奇点}%
\newcommand{\trClassicalMechanics}{经典力学}%
\newcommand{\trConleyIndex}{\trConley{}指数}%
\newcommand{\trIsolatedInvariantSet}{孤立不变集}%
\newcommand{\trContinuation}{延拓}%
\newcommand{\trPerturbation}{扰动}%
\newcommand{\trBifurcation}{分岔}%
\newcommand{\trOscillation}{振荡}%
\newcommand{\trForcedOscillation}{受迫振荡}%
\newcommand{\trRestPoint}{驻点}%
\newcommand{\trIntegrableSystem}{可积系统}%
\newcommand{\trObstruction}{阻碍}%
\newcommand{\trInitialValueProblem}{初值问题}%
\newcommand{\trDifferentialEquation}{微分方程}%
\newcommand{\trPartialDifferentialEquation}{偏微分方程}%
\newcommand{\trTimeDependent}{时间依赖}%
\newcommand{\trTimeOneMap}{时间-1映射}%
\newcommand{\trAdmissible}{容许}%
\newcommand{\trTransformationLaw}{变换法则}%
\newcommand{\trTransformation}{变换}%
\newcommand{\trElementary}{初等}%
\newcommand{\trPermutation}{置换}%
\newcommand{\trScalarFunction}{数值函数}%
\newcommand{\trLattice}{格}%
\newcommand{\trIsotropyGroup}{迷向群}%
\newcommand{\trFractional}{分数阶}%
\newcommand{\trUnimodular}{幺模}%
\newcommand{\trPhaseShift}{相移}%
\newcommand{\trPartitionOfUnity}{单位分解}%
\newcommand{\trAtlas}{坐标图册}%
\newcommand{\trGreenFunction}{格林函数}%
\newcommand{\trSupNorm}{上确界范数}%
\newcommand{\trDifferenceQuotient}{差商}%
\newcommand{\trDirectionalDerivative}{方向导数}%
\newcommand{\trCone}{锥}%
\newcommand{\trDistribution}{广义函数}%
\newcommand{\trRemovableSingularity}{可去奇点}%
\newcommand{\trTotallyReal}{全实}%
\newcommand{\trHalfBall}{半球}%
\newcommand{\trEssentiallyBounded}{本质有界}%
\newcommand{\trHomomorphism}{同态}%
\newcommand{\trContravariant}{反变}%
\newcommand{\trFunctor}{函子}%
\newcommand{\trCategoryTheory}{范畴}%
\newcommand{\trMorphism}{态射}%
\newcommand{\trAlgebraicCategory}{代数范畴}%
\newcommand{\trAlgebraicObject}{代数对象}%
\newcommand{\trFixedPointIndex}{不动点指数}%

\newcommand{\trVariational}{变分}%
\newcommand{\trVariationalPrinciple}{变分原理}%
\newcommand{\trCalculusOfVariations}{变分法}%
\newcommand{\trDirectMethod}{直接方法}%
\newcommand{\trMinimizingSequence}{极小化序列}%
\newcommand{\trMinimax}{极小极大}%
\newcommand{\trCritical}{临界}%
\newcommand{\trCriticalPoint}{临界点}%
\newcommand{\trCriticalValue}{临界值}%
\newcommand{\trSaddlePoint}{鞍点}%
\newcommand{\trGradient}{梯度}%
\newcommand{\trGradientLike}{类梯度}%
\newcommand{\trPseudoGradient}{伪梯度}%
\newcommand{\trMorseTheory}{\trMorse{}理论}%
\newcommand{\trMorseIndex}{\trMorse{}指数}%
\newcommand{\trMorseInequalities}{\trMorse{}不等式}%
\newcommand{\trCategory}{畴数}%
\newcommand{\trSmoothCategory}{光滑范畴}%
\newcommand{\trGenus}{亏格}%
\newcommand{\trSuspension}{悬垂}%
\newcommand{\trFunctional}{泛函}%
\newcommand{\trTestFunction}{检验函数}%
\newcommand{\trCoercive}{强制}%
\newcommand{\trBoundedBelow}{下方有界}%
\newcommand{\trBoundedAbove}{上方有界}%
\newcommand{\trPalaisSmale}{\trPalais-\trSmale}%
\newcommand{\trMountainPassLemma}{山路引理}%

\newcommand{\trHomology}{同调}%
\newcommand{\trCohomology}{上同调}%
\newcommand{\trHomologous}{同调的}%
\newcommand{\trRelativeHomology}{相对同调}%
\newcommand{\trSingularHomology}{奇异同调}%
\newcommand{\trExact}{正合}%
\newcommand{\trExactSequence}{正合序列}%
\newcommand{\trChainComplex}{链复形}%
\newcommand{\trBoundaryOperator}{边缘算子}%
\newcommand{\trCycle}{闭链}%
\newcommand{\trCocycle}{上闭链}%
\newcommand{\trCohomologyClass}{上同调类}%
\newcommand{\trCupProduct}{杯积}%
\newcommand{\trBettiNumber}{贝蒂数}%
\newcommand{\trEulerCharacteristic}{欧拉示性数}%
\newcommand{\trFundamentalGroup}{基本群}%
\newcommand{\trHomotopy}{同伦}%
\newcommand{\trHomotopic}{同伦的}%
\newcommand{\trHomotopyInvariant}{同伦不变量}%
\newcommand{\trInvariance}{不变性}%
\newcommand{\trIsotopy}{同痕}%
\newcommand{\trIsotopic}{同痕的}%
\newcommand{\trContractible}{可缩}%
\newcommand{\trCoveringSpace}{覆盖空间}%
\newcommand{\trUniversalCover}{泛覆盖}%
\newcommand{\trDeformation}{形变}%
\newcommand{\trDeformationRetract}{形变收缩}%
\newcommand{\trQuotient}{商}%
\newcommand{\trDegree}{度}%
\newcommand{\trMonotonicity}{单调性}%
\newcommand{\trExcision}{切除}%
\newcommand{\trAdditivity}{可加性}%
\newcommand{\trNormalization}{规范性}%
\newcommand{\trNaturality}{自然性}%
\newcommand{\trConformality}{共形性}%
\newcommand{\trSolutionProperty}{解性质}%
\newcommand{\trHomotopyInvarianceAxiom}{同伦不变性公理}%

\newcommand{\trManifold}{流形}%
\newcommand{\trSubmanifold}{子流形}%
\newcommand{\trHypersurface}{超曲面}%
\newcommand{\trHyperplane}{超平面}%
\newcommand{\trSphere}{球面}%
\newcommand{\trSurface}{曲面}%
\newcommand{\trVectorField}{向量场}%
\newcommand{\trDiffeomorphism}{微分同胚}%
\newcommand{\trDiffeomorphic}{微分同胚的}%
\newcommand{\trHomeomorphism}{同胚}%
\newcommand{\trEmbedding}{嵌入}%
\newcommand{\trImmersion}{浸入}%
\newcommand{\trDifferentialForm}{微分形式}%
\newcommand{\trOneForm}{1-形式}%
\newcommand{\trTwoForm}{2-形式}%
\newcommand{\trExteriorDerivative}{外微分}%
\newcommand{\trCartanCalculus}{\trCartan{}演算}%
\newcommand{\trPullback}{拉回}%
\newcommand{\trPushforward}{推前}%
\newcommand{\trNondegenerate}{非退化}%
\newcommand{\trDegenerate}{退化}%
\newcommand{\trSkewSymmetric}{反对称}%
\let\trAntisymmetric\trSkewSymmetric%
\newcommand{\trSymmetric}{对称}%
\newcommand{\trDiagonal}{对角}%
\newcommand{\trBilinear}{双线性}%
\newcommand{\trOrthogonal}{正交}%
\newcommand{\trOrthogonality}{正交性}%
\newcommand{\trOrthogonalComplement}{正交补}%
\newcommand{\trOrthogonalSum}{正交和}%
\newcommand{\trProjection}{投影}%
\newcommand{\trEuclidean}{欧几里得}%
\newcommand{\trInnerProduct}{内积}%
\newcommand{\trDual}{对偶}%
\newcommand{\trSubspace}{子空间}%
\newcommand{\trProductSpace}{乘积空间}%
\newcommand{\trInjective}{单射}%
\newcommand{\trSurjective}{满射}%
\newcommand{\trNonsingular}{非奇异}%
\newcommand{\trBundle}{丛}%
\newcommand{\trVectorBundle}{向量丛}%
\newcommand{\trLineBundle}{线丛}%
\newcommand{\trSection}{截面}%
\newcommand{\trFiber}{纤维}%
\newcommand{\trTangentBundle}{切丛}%
\newcommand{\trCotangentBundle}{余切丛}%
\newcommand{\trNormalBundle}{法丛}%
\newcommand{\trTangentSpace}{切空间}%
\newcommand{\trCotangentSpace}{余切空间}%
\newcommand{\trPhaseSpace}{相空间}%
\newcommand{\trConfigurationSpace}{位形空间}%
\newcommand{\trLoop}{闭路}%
\newcommand{\trLoopSpace}{闭路空间}%
\newcommand{\trGeodesic}{测地线}%
\newcommand{\trPath}{路径}%
\newcommand{\trPathConnected}{道路连通}%
\newcommand{\trParameterizedHypersurfaceFamily}{参数化超曲面族}%
\newcommand{\trIsometry}{等距}%
\newcommand{\trClosed}{闭}%
\newcommand{\trCurve}{曲线}%
\newcommand{\trSmooth}{光滑}%
\newcommand{\trDerivative}{导数}%
\newcommand{\trMetric}{度量}%
\newcommand{\trFubiniStudyMetric}{富比尼-施图迪度量}%
\newcommand{\trBiInvariantMetric}{双不变度量}%
\newcommand{\trDistance}{距离}%
\newcommand{\trTopology}{拓扑}%
\newcommand{\trCompact}{紧致}%
\newcommand{\trCompactness}{紧性}%
\newcommand{\trCompactSupport}{紧支撑}%
\newcommand{\trSupport}{支撑}%
\newcommand{\trConnected}{连通}%
\newcommand{\trConnectedComponent}{连通分支}%
\newcommand{\trSeparate}{分离}%
\newcommand{\trSimplyConnected}{单连通}%
\newcommand{\trOriented}{定向}%
\newcommand{\trOrientable}{可定向}%
\newcommand{\trOrientation}{定向}%
\newcommand{\trConvex}{凸}%
\newcommand{\trStrictlyConvex}{严格凸}%
\newcommand{\trConvexity}{凸性}%
\newcommand{\trConcave}{凹}%
\newcommand{\trStarShaped}{星形}%
\newcommand{\trDomain}{区域}%
\newcommand{\trBoundary}{边界}%
\newcommand{\trBoundaryCondition}{边界条件}%
\newcommand{\trBoundaryValueProblem}{边值问题}%
\newcommand{\trInterior}{内部}%
\newcommand{\trClosure}{闭包}%
\newcommand{\trNeighborhood}{邻域}%
\newcommand{\trHolomorphic}{全纯}%
\newcommand{\trPseudoholomorphic}{伪全纯}%
\newcommand{\trConformal}{共形}%
\newcommand{\trHarmonic}{调和}%
\newcommand{\trAnalytic}{解析}%
\newcommand{\trStandard}{标准}%
\newcommand{\trExotic}{怪异}%

\newcommand{\trOperator}{算子}%
\newcommand{\trLinear}{线性}%
\newcommand{\trLinearOperator}{线性算子}%
\newcommand{\trLinearlyIndependent}{线性无关}%
\newcommand{\trRationallyIndependent}{有理无关}%
\newcommand{\trLinearTransformation}{线性变换}%
\newcommand{\trChainRule}{链式法则}%
\newcommand{\trLinearization}{线性化}%
\newcommand{\trNonlinear}{非线性}%
\newcommand{\trHomogeneous}{齐次}%
\newcommand{\trHomogeneity}{齐次性}%
\newcommand{\trSelfAdjoint}{自伴}%
\newcommand{\trAdjoint}{伴随}%
\newcommand{\trNormalOperator}{正规算子}%
\newcommand{\trBounded}{有界}%
\newcommand{\trCountable}{可数}%
\newcommand{\trAlmostEvery}{几乎所有}%
\newcommand{\trAlmostEverywhere}{几乎处处}%
\newcommand{\trUnbounded}{无界}%
\newcommand{\trCompactOperator}{紧算子}%
\newcommand{\trKernel}{核}%
\newcommand{\trRank}{秩}%
\newcommand{\trNorm}{范数}%
\newcommand{\trHessian}{黑塞矩阵}%
\newcommand{\trCodimension}{余维数}%
\newcommand{\trLevelSet}{水平集}%
\newcommand{\trParametrization}{参数化}%
\newcommand{\trVariationOfConstantsFormula}{常数变易公式}%
\newcommand{\trReparametrization}{重新参数化}%
\newcommand{\trOpenSet}{开集}%
\newcommand{\trRange}{值域}%
\newcommand{\trEigenvalue}{特征值}%
\newcommand{\trExponentialMap}{指数映射}%
\newcommand{\trEigenvector}{特征向量}%
\newcommand{\trCharacteristicPolynomial}{特征多项式}%
\newcommand{\trReferenceSolution}{参考解}%
\newcommand{\trCounterexample}{反例}%
\newcommand{\trOvertwisted}{过扭}%
\newcommand{\trTight}{紧}%
\newcommand{\trPlaneField}{平面场}%
\newcommand{\trHandlebody}{柄体}%
\newcommand{\trContractionMap}{收缩映射}%
\newcommand{\trDilatation}{伸缩}%
\newcommand{\trEigenfunction}{特征函数}%
\newcommand{\trEigenspace}{特征空间}%
\newcommand{\trMultiplicity}{重数}%
\newcommand{\trSpectrum}{谱}%
\newcommand{\trSpectral}{谱}%
\newcommand{\trElliptic}{椭圆}%
\newcommand{\trEllipticEstimate}{椭圆估计}%
\newcommand{\trEllipticRegularity}{椭圆正则性}%
\newcommand{\trEstimate}{估计}%
\newcommand{\trWeakSolution}{弱解}%
\newcommand{\trClassicalSolution}{古典解}%
\newcommand{\trRegularity}{正则性}%
\newcommand{\trRegularization}{正则化}%
\newcommand{\trFormalPowerSeries}{形式幂级数}%
\newcommand{\trSmallDenominator}{小分母}%
\newcommand{\trSmoothness}{光滑性}%
\newcommand{\trDense}{稠密}%
\newcommand{\trNowhereDense}{无处稠密}%
\newcommand{\trApproximation}{逼近}%
\newcommand{\trImplicitFunctionTheorem}{隐函数定理}%
\newcommand{\trGeneric}{通有}%
\newcommand{\trResidual}{剩余}%
\newcommand{\trTransversality}{横截性}%
\newcommand{\trTransversal}{横截}%
\newcommand{\trTransversalSection}{横截面}%
\newcommand{\trSimilarityPrinciple}{相似原理}%
\newcommand{\trBubblingOffAnalysis}{冒泡分析}%
\newcommand{\trContinuity}{连续性}%
\newcommand{\trCauchyRiemann}{\trCauchy-\trRiemann}%
\newcommand{\trConverge}{收敛}%
\newcommand{\trAsymptotic}{渐近}%
\newcommand{\trDecay}{衰减}%
\newcommand{\trIsomorphism}{同构}%
\newcommand{\trIsomorphic}{同构的}%
\newcommand{\trCutoff}{截断}%
\newcommand{\trHausdorff}{豪斯多夫}%
\newcommand{\trParacompact}{仿紧}%
\newcommand{\trMetrizable}{可度量化}%
\newcommand{\trNeighborhoodRetract}{邻域收缩核}%
\newcommand{\trTaut}{贴合}%


\newtheoremstyle{hzplain}{\topsep}{\topsep}{\itshape}{}
  {\bfseries}{}{.5em}{#1~#2\ifstrempty{#3}{}{（#3）}}
\newtheoremstyle{hzdefinition}{\topsep}{\topsep}{}{}
  {\bfseries}{}{.5em}{#1~#2\ifstrempty{#3}{}{（#3）}}
\newtheoremstyle{hzremark}{\topsep}{\topsep}{\upshape}{}
  {\itshape}{}{.5em}{#1~#2\ifstrempty{#3}{}{（#3）}}

\theoremstyle{hzplain}
\newtheorem{theorem}{\trhTheorem}
\newtheorem{proposition}{\trhProposition}
\newtheorem{lemma}{\trhLemma}
\newtheorem{corollary}{\trhCorollary}
\theoremstyle{hzdefinition}
\newtheorem{definition}{\trhDefinition}
\newtheorem{example}{\trhExample}
\newtheorem{axiom}{\trhAxiom}
\theoremstyle{hzremark}
\newtheorem{remark}{\trhRemark}

\counterwithin*{theorem}{chapter}
\counterwithin*{proposition}{chapter}
\counterwithin*{lemma}{chapter}
\counterwithin*{corollary}{chapter}
\counterwithin*{definition}{chapter}
\counterwithin*{example}{chapter}
\counterwithin*{axiom}{chapter}
\counterwithin*{remark}{chapter}

\newcommand{\hzsync}[2]{\setcounter{#1}{\numexpr(#2)-1\relax}}

\renewcommand{\proofname}{\trhProof}
\renewcommand{\qedsymbol}{$\blacksquare$}

\numberwithin{equation}{chapter}
\counterwithin{figure}{chapter}

\newcommand{\epigraph}[2]{%
  \begin{center}\begin{minipage}{0.86\textwidth}
    {\kaishu #1}\par\vspace{2pt}
    \hfill{\kaishu\small #2}\par
  \end{minipage}\end{center}\vspace{6pt}
}

\newcommand{\R}{\mathbb{R}}
\newcommand{\N}{\mathbb{N}}
\newcommand{\Z}{\mathbb{Z}}
\newcommand{\C}{\mathbb{C}}
\newcommand{\Q}{\mathbb{Q}}
\newcommand{\Sp}{\operatorname{Sp}}     %
\newcommand{\GL}{\operatorname{GL}}
\newcommand{\SO}{\operatorname{SO}}
\newcommand{\Diff}{\operatorname{Diff}}
\newcommand{\id}{\mathrm{id}}
\newcommand{\Hh}[2]{H^{#1/#2}}          %
\newcommand{\osc}{\operatorname{osc}}   %

\usepackage{xcolor}
\usepackage{hyperref}
\hypersetup{
  colorlinks=true,
  linkcolor=blue!55!black,
  citecolor=teal,
  urlcolor=violet,
  bookmarksnumbered=true,
  pdflang={zh-CN},
  pdftitle={辛不变量与哈密顿动力学（中文译本）},
  pdfauthor={Helmut Hofer, Eduard Zehnder（原著）},
}

\usepackage{imakeidx}
\makeindex[title={译本术语索引},intoc,noautomatic]
\pdfstringdefDisableCommands{\def\index#1{}}
\makeatletter
\let\SIHD@tableofcontents\tableofcontents
\def\tableofcontents{\begingroup\def\index##1{}\SIHD@tableofcontents\endgroup}
\makeatother

\usepackage[capitalise,noabbrev]{cleveref}
\crefname{theorem}    {\trhTheorem}    {\trhTheorem}
\crefname{proposition}{\trhProposition}{\trhProposition}
\crefname{lemma}      {\trhLemma}      {\trhLemma}
\crefname{corollary}  {\trhCorollary}  {\trhCorollary}
\crefname{definition} {\trhDefinition} {\trhDefinition}
\crefname{example}    {\trhExample}    {\trhExample}
\crefname{axiom}      {\trhAxiom}      {\trhAxiom}
\crefname{remark}     {\trhRemark}     {\trhRemark}
\crefname{figure}     {图}{图}
\crefname{equation}   {}{}  %
\crefformat{equation}{#2(#1)#3}

\usepackage{booktabs,longtable,multicol}
\newcommand{\BookGlossEntry}[2]{%
  \par\noindent\hangindent=1em\hangafter=1 #1\quad #2\par}
\newcommand{\BookSourceIndexEntry}[2]{%
  \par\begingroup\leftskip=#1em\relax\hangindent=1em\relax
  \noindent #2\par\endgroup}
\makeatletter
\newif\ifBookIndexing
\BookIndexingfalse
\newcommand{\BookBoldPage}[1]{\textbf{\hyperpage{#1}}}
\DeclareRobustCommand{\BookRecordIndex}[1]{%
  \ifBookIndexing
    \expandafter\index\expandafter{\csname BookIndex@#1\endcsname|hyperpage}%
  \fi}
\DeclareRobustCommand{\BookIndexDefinition}[1]{%
  \ifBookIndexing
    \expandafter\index\expandafter{\csname BookIndex@#1\endcsname|BookBoldPage}%
  \fi}
\DeclareRobustCommand{\BookTerm}[2]{%
  \BookRecordIndex{#1}%
  \begingroup\let\BookRecordIndex\@gobble #2\endgroup}
\AddToHook{cmd/@outputpage/before}{\let\BookRecordIndex\@gobble}
\pdfstringdefDisableCommands{%
  \def\BookTerm#1#2{#2}%
  \let\BookRecordIndex\@gobble
  \let\BookIndexDefinition\@gobble}
\expandafter\def\csname BookIndex@T0001\endcsname{xin@辛（symplectic）}
\expandafter\let\csname BookOriginal@trSymplectic\endcsname\trSymplectic
\renewcommand{\trSymplectic}{\BookTerm{T0001}{\csname BookOriginal@trSymplectic\endcsname}}
\expandafter\def\csname BookIndex@T0002\endcsname{fanxin@反辛（antisymplectic）}
\expandafter\let\csname BookOriginal@trAntisymplectic\endcsname\trAntisymplectic
\renewcommand{\trAntisymplectic}{\BookTerm{T0002}{\csname BookOriginal@trAntisymplectic\endcsname}}
\expandafter\def\csname BookIndex@T0003\endcsname{xinjihe@辛几何（symplectic geometry）}
\expandafter\let\csname BookOriginal@trSymplecticGeometry\endcsname\trSymplecticGeometry
\renewcommand{\trSymplecticGeometry}{\BookTerm{T0003}{\csname BookOriginal@trSymplecticGeometry\endcsname}}
\expandafter\def\csname BookIndex@T0004\endcsname{xintapu@辛拓扑（symplectic topology）}
\expandafter\let\csname BookOriginal@trSymplecticTopology\endcsname\trSymplecticTopology
\renewcommand{\trSymplecticTopology}{\BookTerm{T0004}{\csname BookOriginal@trSymplecticTopology\endcsname}}
\expandafter\def\csname BookIndex@T0005\endcsname{xinliuxing@辛流形（symplectic manifold）}
\expandafter\let\csname BookOriginal@trSymplecticManifold\endcsname\trSymplecticManifold
\renewcommand{\trSymplecticManifold}{\BookTerm{T0005}{\csname BookOriginal@trSymplecticManifold\endcsname}}
\expandafter\def\csname BookIndex@T0006\endcsname{xinxiangliangkongjian@辛向量空间（symplectic vector space）}
\expandafter\let\csname BookOriginal@trSymplecticVectorSpace\endcsname\trSymplecticVectorSpace
\renewcommand{\trSymplecticVectorSpace}{\BookTerm{T0006}{\csname BookOriginal@trSymplecticVectorSpace\endcsname}}
\expandafter\def\csname BookIndex@T0007\endcsname{xinxingshi@辛形式（symplectic form）}
\expandafter\let\csname BookOriginal@trSymplecticForm\endcsname\trSymplecticForm
\renewcommand{\trSymplecticForm}{\BookTerm{T0007}{\csname BookOriginal@trSymplecticForm\endcsname}}
\expandafter\def\csname BookIndex@T0008\endcsname{xinjiegou@辛结构（symplectic structure）}
\expandafter\let\csname BookOriginal@trSymplecticStructure\endcsname\trSymplecticStructure
\renewcommand{\trSymplecticStructure}{\BookTerm{T0008}{\csname BookOriginal@trSymplecticStructure\endcsname}}
\expandafter\def\csname BookIndex@T0009\endcsname{xinjuzhen@辛矩阵（symplectic matrix）}
\expandafter\let\csname BookOriginal@trSymplecticMatrix\endcsname\trSymplecticMatrix
\renewcommand{\trSymplecticMatrix}{\BookTerm{T0009}{\csname BookOriginal@trSymplecticMatrix\endcsname}}
\expandafter\def\csname BookIndex@T0010\endcsname{qunzuoyong@群作用（group action）}
\expandafter\let\csname BookOriginal@trGroupAction\endcsname\trGroupAction
\renewcommand{\trGroupAction}{\BookTerm{T0010}{\csname BookOriginal@trGroupAction\endcsname}}
\expandafter\def\csname BookIndex@T0011\endcsname{jiaohuan@交换（abelian）}
\expandafter\let\csname BookOriginal@trAbelian\endcsname\trAbelian
\renewcommand{\trAbelian}{\BookTerm{T0011}{\csname BookOriginal@trAbelian\endcsname}}
\expandafter\def\csname BookIndex@T0012\endcsname{yingshe@映射（mapping）}
\expandafter\let\csname BookOriginal@trMapping\endcsname\trMapping
\renewcommand{\trMapping}{\BookTerm{T0012}{\csname BookOriginal@trMapping\endcsname}}
\expandafter\def\csname BookIndex@T0013\endcsname{hengtongyingshe@恒同映射（identity map）}
\expandafter\let\csname BookOriginal@trIdentityMap\endcsname\trIdentityMap
\renewcommand{\trIdentityMap}{\BookTerm{T0013}{\csname BookOriginal@trIdentityMap\endcsname}}
\expandafter\def\csname BookIndex@T0014\endcsname{xinyingshe@辛映射（symplectic mapping）}
\expandafter\let\csname BookOriginal@trSymplecticMapping\endcsname\trSymplecticMapping
\renewcommand{\trSymplecticMapping}{\BookTerm{T0014}{\csname BookOriginal@trSymplecticMapping\endcsname}}
\expandafter\def\csname BookIndex@T0015\endcsname{xinweifentongpei@辛微分同胚（symplectomorphism）}
\expandafter\let\csname BookOriginal@trSymplecticDiffeomorphism\endcsname\trSymplecticDiffeomorphism
\renewcommand{\trSymplecticDiffeomorphism}{\BookTerm{T0015}{\csname BookOriginal@trSymplecticDiffeomorphism\endcsname}}
\expandafter\def\csname BookIndex@T0016\endcsname{xinweifentongpei@辛微分同胚（symplectomorphism）}
\expandafter\let\csname BookOriginal@trSymplectomorphism\endcsname\trSymplectomorphism
\renewcommand{\trSymplectomorphism}{\BookTerm{T0016}{\csname BookOriginal@trSymplectomorphism\endcsname}}
\expandafter\def\csname BookIndex@T0017\endcsname{xintongpei@辛同胚（symplectic homeomorphism）}
\expandafter\let\csname BookOriginal@trSymplecticHomeomorphism\endcsname\trSymplecticHomeomorphism
\renewcommand{\trSymplecticHomeomorphism}{\BookTerm{T0017}{\csname BookOriginal@trSymplecticHomeomorphism\endcsname}}
\expandafter\def\csname BookIndex@T0018\endcsname{xinqianru@辛嵌入（symplectic embedding）}
\expandafter\let\csname BookOriginal@trSymplecticEmbedding\endcsname\trSymplecticEmbedding
\renewcommand{\trSymplecticEmbedding}{\BookTerm{T0018}{\csname BookOriginal@trSymplecticEmbedding\endcsname}}
\expandafter\def\csname BookIndex@T0019\endcsname{xinrongliang@辛容量（symplectic capacity）}
\expandafter\let\csname BookOriginal@trSymplecticCapacity\endcsname\trSymplecticCapacity
\renewcommand{\trSymplecticCapacity}{\BookTerm{T0019}{\csname BookOriginal@trSymplecticCapacity\endcsname}}
\expandafter\def\csname BookIndex@T0020\endcsname{xinbubianliang@辛不变量（symplectic invariant）}
\expandafter\let\csname BookOriginal@trSymplecticInvariant\endcsname\trSymplecticInvariant
\renewcommand{\trSymplecticInvariant}{\BookTerm{T0020}{\csname BookOriginal@trSymplecticInvariant\endcsname}}
\expandafter\def\csname BookIndex@T0021\endcsname{xinji@辛基（symplectic basis）}
\expandafter\let\csname BookOriginal@trSymplecticBasis\endcsname\trSymplecticBasis
\renewcommand{\trSymplecticBasis}{\BookTerm{T0021}{\csname BookOriginal@trSymplecticBasis\endcsname}}
\expandafter\def\csname BookIndex@T0022\endcsname{xinzuobiao@辛坐标（symplectic coordinates）}
\expandafter\let\csname BookOriginal@trSymplecticCoordinates\endcsname\trSymplecticCoordinates
\renewcommand{\trSymplecticCoordinates}{\BookTerm{T0022}{\csname BookOriginal@trSymplecticCoordinates\endcsname}}
\expandafter\def\csname BookIndex@T0023\endcsname{xinbu@辛补（symplectic complement）}
\expandafter\let\csname BookOriginal@trSymplecticComplement\endcsname\trSymplecticComplement
\renewcommand{\trSymplecticComplement}{\BookTerm{T0023}{\csname BookOriginal@trSymplecticComplement\endcsname}}
\expandafter\def\csname BookIndex@T0024\endcsname{xintonggou@辛同构（symplectic isomorphism）}
\expandafter\let\csname BookOriginal@trSymplecticIsomorphism\endcsname\trSymplecticIsomorphism
\renewcommand{\trSymplecticIsomorphism}{\BookTerm{T0024}{\csname BookOriginal@trSymplecticIsomorphism\endcsname}}
\expandafter\def\csname BookIndex@T0025\endcsname{rongliang@容量（capacity）}
\expandafter\let\csname BookOriginal@trCapacity\endcsname\trCapacity
\renewcommand{\trCapacity}{\BookTerm{T0025}{\csname BookOriginal@trCapacity\endcsname}}
\expandafter\def\csname BookIndex@T0026\endcsname{neirongliang@内容量（inner capacity）}
\expandafter\let\csname BookOriginal@trInnerCapacity\endcsname\trInnerCapacity
\renewcommand{\trInnerCapacity}{\BookTerm{T0026}{\csname BookOriginal@trInnerCapacity\endcsname}}
\expandafter\def\csname BookIndex@T0027\endcsname{wairongliang@外容量（outer capacity）}
\expandafter\let\csname BookOriginal@trOuterCapacity\endcsname\trOuterCapacity
\renewcommand{\trOuterCapacity}{\BookTerm{T0027}{\csname BookOriginal@trOuterCapacity\endcsname}}
\expandafter\def\csname BookIndex@T0028\endcsname{neizhengzexing@内正则性（inner regularity）}
\expandafter\let\csname BookOriginal@trInnerRegularity\endcsname\trInnerRegularity
\renewcommand{\trInnerRegularity}{\BookTerm{T0028}{\csname BookOriginal@trInnerRegularity\endcsname}}
\expandafter\def\csname BookIndex@T0029\endcsname{waizhengzexing@外正则性（outer regularity）}
\expandafter\let\csname BookOriginal@trOuterRegularity\endcsname\trOuterRegularity
\renewcommand{\trOuterRegularity}{\BookTerm{T0029}{\csname BookOriginal@trOuterRegularity\endcsname}}
\expandafter\def\csname BookIndex@T0030\endcsname{feipingfanxing@非平凡性（nontriviality）}
\expandafter\let\csname BookOriginal@trNontriviality\endcsname\trNontriviality
\renewcommand{\trNontriviality}{\BookTerm{T0030}{\csname BookOriginal@trNontriviality\endcsname}}
\expandafter\def\csname BookIndex@T0031\endcsname{baohanyingshe@包含映射（inclusion mapping）}
\expandafter\let\csname BookOriginal@trInclusionMapping\endcsname\trInclusionMapping
\renewcommand{\trInclusionMapping}{\BookTerm{T0031}{\csname BookOriginal@trInclusionMapping\endcsname}}
\expandafter\def\csname BookIndex@T0032\endcsname{danzhihua@单值化（uniformization）}
\expandafter\let\csname BookOriginal@trUniformization\endcsname\trUniformization
\renewcommand{\trUniformization}{\BookTerm{T0032}{\csname BookOriginal@trUniformization\endcsname}}
\expandafter\def\csname BookIndex@T0033\endcsname{geluomofukuandu@格罗莫夫宽度（Gromov width）}
\expandafter\let\csname BookOriginal@trGromovWidth\endcsname\trGromovWidth
\renewcommand{\trGromovWidth}{\BookTerm{T0033}{\csname BookOriginal@trGromovWidth\endcsname}}
\expandafter\def\csname BookIndex@T0034\endcsname{jiyadingli@挤压定理（squeezing theorem）}
\expandafter\let\csname BookOriginal@trSqueezingTheorem\endcsname\trSqueezingTheorem
\renewcommand{\trSqueezingTheorem}{\BookTerm{T0034}{\csname BookOriginal@trSqueezingTheorem\endcsname}}
\expandafter\def\csname BookIndex@T0035\endcsname{hamidun@哈密顿（Hamiltonian）}
\expandafter\let\csname BookOriginal@trHamiltonian\endcsname\trHamiltonian
\renewcommand{\trHamiltonian}{\BookTerm{T0035}{\csname BookOriginal@trHamiltonian\endcsname}}
\expandafter\def\csname BookIndex@T0036\endcsname{hamidunxiangliangchang@哈密顿向量场（Hamiltonian vector field）}
\expandafter\let\csname BookOriginal@trHamiltonianVectorField\endcsname\trHamiltonianVectorField
\renewcommand{\trHamiltonianVectorField}{\BookTerm{T0036}{\csname BookOriginal@trHamiltonianVectorField\endcsname}}
\expandafter\def\csname BookIndex@T0037\endcsname{hamidunxitong@哈密顿系统（Hamiltonian system）}
\expandafter\let\csname BookOriginal@trHamiltonianSystem\endcsname\trHamiltonianSystem
\renewcommand{\trHamiltonianSystem}{\BookTerm{T0037}{\csname BookOriginal@trHamiltonianSystem\endcsname}}
\expandafter\def\csname BookIndex@T0038\endcsname{hamidunfangcheng@哈密顿方程（Hamiltonian equation）}
\expandafter\let\csname BookOriginal@trHamiltonianEquation\endcsname\trHamiltonianEquation
\renewcommand{\trHamiltonianEquation}{\BookTerm{T0038}{\csname BookOriginal@trHamiltonianEquation\endcsname}}
\expandafter\def\csname BookIndex@T0039\endcsname{hamidunhanshu@哈密顿函数（Hamiltonian function）}
\expandafter\let\csname BookOriginal@trHamiltonianFunction\endcsname\trHamiltonianFunction
\renewcommand{\trHamiltonianFunction}{\BookTerm{T0039}{\csname BookOriginal@trHamiltonianFunction\endcsname}}
\expandafter\def\csname BookIndex@T0040\endcsname{hamidundonglixue@哈密顿动力学（Hamiltonian dynamics）}
\expandafter\let\csname BookOriginal@trHamiltonianDynamics\endcsname\trHamiltonianDynamics
\renewcommand{\trHamiltonianDynamics}{\BookTerm{T0040}{\csname BookOriginal@trHamiltonianDynamics\endcsname}}
\expandafter\def\csname BookIndex@T0041\endcsname{hamidunyingshe@哈密顿映射（Hamiltonian map）}
\expandafter\let\csname BookOriginal@trHamiltonianMap\endcsname\trHamiltonianMap
\renewcommand{\trHamiltonianMap}{\BookTerm{T0041}{\csname BookOriginal@trHamiltonianMap\endcsname}}
\expandafter\def\csname BookIndex@T0042\endcsname{hamidunweifentongpei@哈密顿微分同胚（Hamiltonian diffeomorphism）}
\expandafter\let\csname BookOriginal@trHamiltonianDiffeomorphism\endcsname\trHamiltonianDiffeomorphism
\renewcommand{\trHamiltonianDiffeomorphism}{\BookTerm{T0042}{\csname BookOriginal@trHamiltonianDiffeomorphism\endcsname}}
\expandafter\def\csname BookIndex@T0043\endcsname{hamidunliu@哈密顿流（Hamiltonian flow）}
\expandafter\let\csname BookOriginal@trHamiltonianFlow\endcsname\trHamiltonianFlow
\renewcommand{\trHamiltonianFlow}{\BookTerm{T0043}{\csname BookOriginal@trHamiltonianFlow\endcsname}}
\expandafter\def\csname BookIndex@T0044\endcsname{hamidunxingshitixi@哈密顿形式体系（Hamiltonian formalism）}
\expandafter\let\csname BookOriginal@trHamiltonianFormalism\endcsname\trHamiltonianFormalism
\renewcommand{\trHamiltonianFormalism}{\BookTerm{T0044}{\csname BookOriginal@trHamiltonianFormalism\endcsname}}
\expandafter\def\csname BookIndex@T0045\endcsname{jinfujiegou@近复结构（almost complex structure）}
\expandafter\let\csname BookOriginal@trAlmostComplexStructure\endcsname\trAlmostComplexStructure
\renewcommand{\trAlmostComplexStructure}{\BookTerm{T0045}{\csname BookOriginal@trAlmostComplexStructure\endcsname}}
\expandafter\def\csname BookIndex@T0046\endcsname{fujiegou@复结构（complex structure）}
\expandafter\let\csname BookOriginal@trComplexStructure\endcsname\trComplexStructure
\renewcommand{\trComplexStructure}{\BookTerm{T0046}{\csname BookOriginal@trComplexStructure\endcsname}}
\expandafter\def\csname BookIndex@T0047\endcsname{xiangrong@相容（compatible）}
\expandafter\let\csname BookOriginal@trCompatible\endcsname\trCompatible
\renewcommand{\trCompatible}{\BookTerm{T0047}{\csname BookOriginal@trCompatible\endcsname}}
\expandafter\def\csname BookIndex@T0048\endcsname{qiechu@切触（contact）}
\expandafter\let\csname BookOriginal@trContact\endcsname\trContact
\renewcommand{\trContact}{\BookTerm{T0048}{\csname BookOriginal@trContact\endcsname}}
\expandafter\def\csname BookIndex@T0049\endcsname{qiechujiegou@切触结构（contact structure）}
\expandafter\let\csname BookOriginal@trContactStructure\endcsname\trContactStructure
\renewcommand{\trContactStructure}{\BookTerm{T0049}{\csname BookOriginal@trContactStructure\endcsname}}
\expandafter\def\csname BookIndex@T0050\endcsname{mixiang@迷向（isotropic）}
\expandafter\let\csname BookOriginal@trIsotropic\endcsname\trIsotropic
\renewcommand{\trIsotropic}{\BookTerm{T0050}{\csname BookOriginal@trIsotropic\endcsname}}
\expandafter\def\csname BookIndex@T0051\endcsname{yezhuangjiegou@叶状结构（foliation）}
\expandafter\let\csname BookOriginal@trFoliation\endcsname\trFoliation
\renewcommand{\trFoliation}{\BookTerm{T0051}{\csname BookOriginal@trFoliation\endcsname}}
\expandafter\def\csname BookIndex@T0052\endcsname{ye@叶（leaf）}
\expandafter\let\csname BookOriginal@trLeaf\endcsname\trLeaf
\renewcommand{\trLeaf}{\BookTerm{T0052}{\csname BookOriginal@trLeaf\endcsname}}
\expandafter\def\csname BookIndex@T0053\endcsname{gonge@共轭（conjugate）}
\expandafter\let\csname BookOriginal@trConjugate\endcsname\trConjugate
\renewcommand{\trConjugate}{\BookTerm{T0053}{\csname BookOriginal@trConjugate\endcsname}}
\expandafter\def\csname BookIndex@T0054\endcsname{qiechuxingshi@切触形式（contact form）}
\expandafter\let\csname BookOriginal@trContactForm\endcsname\trContactForm
\renewcommand{\trContactForm}{\BookTerm{T0054}{\csname BookOriginal@trContactForm\endcsname}}
\expandafter\def\csname BookIndex@T0055\endcsname{qiechuxing@切触型（contact type）}
\expandafter\let\csname BookOriginal@trContactType\endcsname\trContactType
\renewcommand{\trContactType}{\BookTerm{T0055}{\csname BookOriginal@trContactType\endcsname}}
\expandafter\def\csname BookIndex@T0056\endcsname{tezheng@特征（characteristic）}
\expandafter\let\csname BookOriginal@trCharacteristic\endcsname\trCharacteristic
\renewcommand{\trCharacteristic}{\BookTerm{T0056}{\csname BookOriginal@trCharacteristic\endcsname}}
\expandafter\def\csname BookIndex@T0057\endcsname{bitezheng@闭特征（closed characteristic）}
\expandafter\let\csname BookOriginal@trClosedCharacteristic\endcsname\trClosedCharacteristic
\renewcommand{\trClosedCharacteristic}{\BookTerm{T0057}{\csname BookOriginal@trClosedCharacteristic\endcsname}}
\expandafter\def\csname BookIndex@T0058\endcsname{tezhengxiancong@特征线丛（characteristic line bundle）}
\expandafter\let\csname BookOriginal@trCharacteristicLineBundle\endcsname\trCharacteristicLineBundle
\renewcommand{\trCharacteristicLineBundle}{\BookTerm{T0058}{\csname BookOriginal@trCharacteristicLineBundle\endcsname}}
\expandafter\def\csname BookIndex@T0059\endcsname{nengliang@能量（energy）}
\expandafter\let\csname BookOriginal@trEnergy\endcsname\trEnergy
\renewcommand{\trEnergy}{\BookTerm{T0059}{\csname BookOriginal@trEnergy\endcsname}}
\expandafter\def\csname BookIndex@T0060\endcsname{nengliangmian@能量面（energy surface）}
\expandafter\let\csname BookOriginal@trEnergySurface\endcsname\trEnergySurface
\renewcommand{\trEnergySurface}{\BookTerm{T0060}{\csname BookOriginal@trEnergySurface\endcsname}}
\expandafter\def\csname BookIndex@T0061\endcsname{zhengze@正则（canonical）}
\expandafter\let\csname BookOriginal@trCanonical\endcsname\trCanonical
\renewcommand{\trCanonical}{\BookTerm{T0061}{\csname BookOriginal@trCanonical\endcsname}}
\expandafter\def\csname BookIndex@T0062\endcsname{zhengze@正则（regular）}
\expandafter\let\csname BookOriginal@trRegular\endcsname\trRegular
\renewcommand{\trRegular}{\BookTerm{T0062}{\csname BookOriginal@trRegular\endcsname}}
\expandafter\def\csname BookIndex@T0063\endcsname{jizuobiao@极坐标（polar coordinates）}
\expandafter\let\csname BookOriginal@trPolarCoordinates\endcsname\trPolarCoordinates
\renewcommand{\trPolarCoordinates}{\BookTerm{T0063}{\csname BookOriginal@trPolarCoordinates\endcsname}}
\expandafter\def\csname BookIndex@T0064\endcsname{zuobiaoka@坐标卡（chart）}
\expandafter\let\csname BookOriginal@trChart\endcsname\trChart
\renewcommand{\trChart}{\BookTerm{T0064}{\csname BookOriginal@trChart\endcsname}}
\expandafter\def\csname BookIndex@T0065\endcsname{lifangti@立方体（cube）}
\expandafter\let\csname BookOriginal@trCube\endcsname\trCube
\renewcommand{\trCube}{\BookTerm{T0065}{\csname BookOriginal@trCube\endcsname}}
\expandafter\def\csname BookIndex@T0066\endcsname{zuoyongliangjiao@作用量-角（action-angle）}
\expandafter\let\csname BookOriginal@trActionAngle\endcsname\trActionAngle
\renewcommand{\trActionAngle}{\BookTerm{T0066}{\csname BookOriginal@trActionAngle\endcsname}}
\expandafter\def\csname BookIndex@T0067\endcsname{shengchenghanshu@生成函数（generating function）}
\expandafter\let\csname BookOriginal@trGeneratingFunction\endcsname\trGeneratingFunction
\renewcommand{\trGeneratingFunction}{\BookTerm{T0067}{\csname BookOriginal@trGeneratingFunction\endcsname}}
\expandafter\def\csname BookIndex@T0068\endcsname{zuoyongliang@作用量（action）}
\expandafter\let\csname BookOriginal@trAction\endcsname\trAction
\renewcommand{\trAction}{\BookTerm{T0068}{\csname BookOriginal@trAction\endcsname}}
\expandafter\def\csname BookIndex@T0069\endcsname{zuoyongliangfanhan@作用量泛函（action functional）}
\expandafter\let\csname BookOriginal@trActionFunctional\endcsname\trActionFunctional
\renewcommand{\trActionFunctional}{\BookTerm{T0069}{\csname BookOriginal@trActionFunctional\endcsname}}
\expandafter\def\csname BookIndex@T0070\endcsname{yuehuazuoyongliang@约化作用量（reduced action）}
\expandafter\let\csname BookOriginal@trReducedAction\endcsname\trReducedAction
\renewcommand{\trReducedAction}{\BookTerm{T0070}{\csname BookOriginal@trReducedAction\endcsname}}
\expandafter\def\csname BookIndex@T0071\endcsname{zuoyongliangpu@作用量谱（action spectrum）}
\expandafter\let\csname BookOriginal@trActionSpectrum\endcsname\trActionSpectrum
\renewcommand{\trActionSpectrum}{\BookTerm{T0071}{\csname BookOriginal@trActionSpectrum\endcsname}}
\expandafter\def\csname BookIndex@T0072\endcsname{zuixiaozuoyongliangyuanli@最小作用量原理（least action principle）}
\expandafter\let\csname BookOriginal@trLeastActionPrinciple\endcsname\trLeastActionPrinciple
\renewcommand{\trLeastActionPrinciple}{\BookTerm{T0072}{\csname BookOriginal@trLeastActionPrinciple\endcsname}}
\expandafter\def\csname BookIndex@T0073\endcsname{weiyinengliang@位移能量（displacement energy）}
\expandafter\let\csname BookOriginal@trDisplacementEnergy\endcsname\trDisplacementEnergy
\renewcommand{\trDisplacementEnergy}{\BookTerm{T0073}{\csname BookOriginal@trDisplacementEnergy\endcsname}}
\expandafter\def\csname BookIndex@T0074\endcsname{fanshi@范式（normal form）}
\expandafter\let\csname BookOriginal@trNormalForm\endcsname\trNormalForm
\renewcommand{\trNormalForm}{\BookTerm{T0074}{\csname BookOriginal@trNormalForm\endcsname}}
\expandafter\def\csname BookIndex@T0075\endcsname{bokehuofububianliang@伯克霍夫不变量（Birkhoff invariant）}
\expandafter\let\csname BookOriginal@trBirkhoffInvariant\endcsname\trBirkhoffInvariant
\renewcommand{\trBirkhoffInvariant}{\BookTerm{T0075}{\csname BookOriginal@trBirkhoffInvariant\endcsname}}
\expandafter\def\csname BookIndex@T0076\endcsname{feigongzhen@非共振（nonresonance）}
\expandafter\let\csname BookOriginal@trNonresonance\endcsname\trNonresonance
\renewcommand{\trNonresonance}{\BookTerm{T0076}{\csname BookOriginal@trNonresonance\endcsname}}
\expandafter\def\csname BookIndex@T0077\endcsname{gongzhen@共振（resonance）}
\expandafter\let\csname BookOriginal@trResonance\endcsname\trResonance
\renewcommand{\trResonance}{\BookTerm{T0077}{\csname BookOriginal@trResonance\endcsname}}
\expandafter\def\csname BookIndex@T0078\endcsname{ercixing@二次型（quadratic form）}
\expandafter\let\csname BookOriginal@trQuadraticForm\endcsname\trQuadraticForm
\renewcommand{\trQuadraticForm}{\BookTerm{T0078}{\csname BookOriginal@trQuadraticForm\endcsname}}
\expandafter\def\csname BookIndex@T0079\endcsname{zhengding@正定（positive definite）}
\expandafter\let\csname BookOriginal@trPositiveDefinite\endcsname\trPositiveDefinite
\renewcommand{\trPositiveDefinite}{\BookTerm{T0079}{\csname BookOriginal@trPositiveDefinite\endcsname}}
\expandafter\def\csname BookIndex@T0080\endcsname{tuoqiu@椭球（ellipsoid）}
\expandafter\let\csname BookOriginal@trEllipsoid\endcsname\trEllipsoid
\renewcommand{\trEllipsoid}{\BookTerm{T0080}{\csname BookOriginal@trEllipsoid\endcsname}}
\expandafter\def\csname BookIndex@T0081\endcsname{zhuti@柱体（cylinder）}
\expandafter\let\csname BookOriginal@trCylinder\endcsname\trCylinder
\renewcommand{\trCylinder}{\BookTerm{T0081}{\csname BookOriginal@trCylinder\endcsname}}
\expandafter\def\csname BookIndex@T0082\endcsname{huanmian@环面（torus）}
\expandafter\let\csname BookOriginal@trTorus\endcsname\trTorus
\renewcommand{\trTorus}{\BookTerm{T0082}{\csname BookOriginal@trTorus\endcsname}}
\expandafter\def\csname BookIndex@T0083\endcsname{yuanzhou@圆周（circle）}
\expandafter\let\csname BookOriginal@trCircle\endcsname\trCircle
\renewcommand{\trCircle}{\BookTerm{T0083}{\csname BookOriginal@trCircle\endcsname}}
\expandafter\def\csname BookIndex@T0084\endcsname{huanxingquyu@环形区域（annulus）}
\expandafter\let\csname BookOriginal@trAnnulus\endcsname\trAnnulus
\renewcommand{\trAnnulus}{\BookTerm{T0084}{\csname BookOriginal@trAnnulus\endcsname}}
\expandafter\def\csname BookIndex@T0085\endcsname{qiu@球（ball）}
\expandafter\let\csname BookOriginal@trBall\endcsname\trBall
\renewcommand{\trBall}{\BookTerm{T0085}{\csname BookOriginal@trBall\endcsname}}
\expandafter\def\csname BookIndex@T0086\endcsname{danweiqiu@单位球（unit ball）}
\expandafter\let\csname BookOriginal@trUnitBall\endcsname\trUnitBall
\renewcommand{\trUnitBall}{\BookTerm{T0086}{\csname BookOriginal@trUnitBall\endcsname}}
\expandafter\def\csname BookIndex@T0087\endcsname{yuanpan@圆盘（disc）}
\expandafter\let\csname BookOriginal@trDisc\endcsname\trDisc
\renewcommand{\trDisc}{\BookTerm{T0087}{\csname BookOriginal@trDisc\endcsname}}
\expandafter\def\csname BookIndex@T0088\endcsname{tiji@体积（volume）}
\expandafter\let\csname BookOriginal@trVolume\endcsname\trVolume
\renewcommand{\trVolume}{\BookTerm{T0088}{\csname BookOriginal@trVolume\endcsname}}
\expandafter\def\csname BookIndex@T0089\endcsname{tijixingshi@体积形式（volume form）}
\expandafter\let\csname BookOriginal@trVolumeForm\endcsname\trVolumeForm
\renewcommand{\trVolumeForm}{\BookTerm{T0089}{\csname BookOriginal@trVolumeForm\endcsname}}
\expandafter\def\csname BookIndex@T0090\endcsname{baotiji@保体积（volume preserving）}
\expandafter\let\csname BookOriginal@trVolumePreserving\endcsname\trVolumePreserving
\renewcommand{\trVolumePreserving}{\BookTerm{T0090}{\csname BookOriginal@trVolumePreserving\endcsname}}
\expandafter\def\csname BookIndex@T0091\endcsname{gangxing@刚性（rigidity）}
\expandafter\let\csname BookOriginal@trRigidity\endcsname\trRigidity
\renewcommand{\trRigidity}{\BookTerm{T0091}{\csname BookOriginal@trRigidity\endcsname}}
\expandafter\def\csname BookIndex@T0092\endcsname{donglixitong@动力系统（dynamical system）}
\expandafter\let\csname BookOriginal@trDynamicalSystem\endcsname\trDynamicalSystem
\renewcommand{\trDynamicalSystem}{\BookTerm{T0092}{\csname BookOriginal@trDynamicalSystem\endcsname}}
\expandafter\def\csname BookIndex@T0093\endcsname{liu@流（flow）}
\expandafter\let\csname BookOriginal@trFlow\endcsname\trFlow
\renewcommand{\trFlow}{\BookTerm{T0093}{\csname BookOriginal@trFlow\endcsname}}
\expandafter\def\csname BookIndex@T0094\endcsname{posongkuohao@泊松括号（Poisson bracket）}
\expandafter\let\csname BookOriginal@trPoissonBracket\endcsname\trPoissonBracket
\renewcommand{\trPoissonBracket}{\BookTerm{T0094}{\csname BookOriginal@trPoissonBracket\endcsname}}
\expandafter\def\csname BookIndex@T0095\endcsname{duihe@对合（involution）}
\expandafter\let\csname BookOriginal@trInvolution\endcsname\trInvolution
\renewcommand{\trInvolution}{\BookTerm{T0095}{\csname BookOriginal@trInvolution\endcsname}}
\expandafter\def\csname BookIndex@T0096\endcsname{jifen@积分（integral）}
\expandafter\let\csname BookOriginal@trIntegral\endcsname\trIntegral
\renewcommand{\trIntegral}{\BookTerm{T0096}{\csname BookOriginal@trIntegral\endcsname}}
\expandafter\def\csname BookIndex@T0097\endcsname{diyijifen@第一积分（first integral）}
\expandafter\let\csname BookOriginal@trFirstIntegral\endcsname\trFirstIntegral
\renewcommand{\trFirstIntegral}{\BookTerm{T0097}{\csname BookOriginal@trFirstIntegral\endcsname}}
\expandafter\def\csname BookIndex@T0098\endcsname{pinluxiangliang@频率向量（frequency vector）}
\expandafter\let\csname BookOriginal@trFrequencyVector\endcsname\trFrequencyVector
\renewcommand{\trFrequencyVector}{\BookTerm{T0098}{\csname BookOriginal@trFrequencyVector\endcsname}}
\expandafter\def\csname BookIndex@T0099\endcsname{judian@聚点（accumulation point）}
\expandafter\let\csname BookOriginal@trAccumulationPoint\endcsname\trAccumulationPoint
\renewcommand{\trAccumulationPoint}{\BookTerm{T0099}{\csname BookOriginal@trAccumulationPoint\endcsname}}
\expandafter\def\csname BookIndex@T0100\endcsname{shuangqu@双曲（hyperbolic）}
\expandafter\let\csname BookOriginal@trHyperbolic\endcsname\trHyperbolic
\renewcommand{\trHyperbolic}{\BookTerm{T0100}{\csname BookOriginal@trHyperbolic\endcsname}}
\expandafter\def\csname BookIndex@T0101\endcsname{fasanxing@发散性（divergence）}
\expandafter\let\csname BookOriginal@trDivergence\endcsname\trDivergence
\renewcommand{\trDivergence}{\BookTerm{T0101}{\csname BookOriginal@trDivergence\endcsname}}
\expandafter\def\csname BookIndex@T0102\endcsname{guidao@轨道（orbit）}
\expandafter\let\csname BookOriginal@trOrbit\endcsname\trOrbit
\renewcommand{\trOrbit}{\BookTerm{T0102}{\csname BookOriginal@trOrbit\endcsname}}
\expandafter\def\csname BookIndex@T0103\endcsname{guixian@轨线（trajectory）}
\expandafter\let\csname BookOriginal@trTrajectory\endcsname\trTrajectory
\renewcommand{\trTrajectory}{\BookTerm{T0103}{\csname BookOriginal@trTrajectory\endcsname}}
\expandafter\def\csname BookIndex@T0104\endcsname{zhouqi@周期（periodic）}
\expandafter\let\csname BookOriginal@trPeriodic\endcsname\trPeriodic
\renewcommand{\trPeriodic}{\BookTerm{T0104}{\csname BookOriginal@trPeriodic\endcsname}}
\expandafter\def\csname BookIndex@T0105\endcsname{zhouqiguidao@周期轨道（periodic orbit）}
\expandafter\let\csname BookOriginal@trPeriodicOrbit\endcsname\trPeriodicOrbit
\renewcommand{\trPeriodicOrbit}{\BookTerm{T0105}{\csname BookOriginal@trPeriodicOrbit\endcsname}}
\expandafter\def\csname BookIndex@T0106\endcsname{nizhouqi@拟周期（quasi periodic）}
\expandafter\let\csname BookOriginal@trQuasiPeriodic\endcsname\trQuasiPeriodic
\renewcommand{\trQuasiPeriodic}{\BookTerm{T0106}{\csname BookOriginal@trQuasiPeriodic\endcsname}}
\expandafter\def\csname BookIndex@T0107\endcsname{zhouqijie@周期解（periodic solution）}
\expandafter\let\csname BookOriginal@trPeriodicSolution\endcsname\trPeriodicSolution
\renewcommand{\trPeriodicSolution}{\BookTerm{T0107}{\csname BookOriginal@trPeriodicSolution\endcsname}}
\expandafter\def\csname BookIndex@T0108\endcsname{budongdian@不动点（fixed point）}
\expandafter\let\csname BookOriginal@trFixedPoint\endcsname\trFixedPoint
\renewcommand{\trFixedPoint}{\BookTerm{T0108}{\csname BookOriginal@trFixedPoint\endcsname}}
\expandafter\def\csname BookIndex@T0109\endcsname{pinghengdian@平衡点（equilibrium）}
\expandafter\let\csname BookOriginal@trEquilibrium\endcsname\trEquilibrium
\renewcommand{\trEquilibrium}{\BookTerm{T0109}{\csname BookOriginal@trEquilibrium\endcsname}}
\expandafter\def\csname BookIndex@T0110\endcsname{huigui@回归（recurrence）}
\expandafter\let\csname BookOriginal@trRecurrence\endcsname\trRecurrence
\renewcommand{\trRecurrence}{\BookTerm{T0110}{\csname BookOriginal@trRecurrence\endcsname}}
\expandafter\def\csname BookIndex@T0111\endcsname{bubian@不变（invariant）}
\expandafter\let\csname BookOriginal@trInvariant\endcsname\trInvariant
\renewcommand{\trInvariant}{\BookTerm{T0111}{\csname BookOriginal@trInvariant\endcsname}}
\expandafter\def\csname BookIndex@T0112\endcsname{bubiancedu@不变测度（invariant measure）}
\expandafter\let\csname BookOriginal@trInvariantMeasure\endcsname\trInvariantMeasure
\renewcommand{\trInvariantMeasure}{\BookTerm{T0112}{\csname BookOriginal@trInvariantMeasure\endcsname}}
\expandafter\def\csname BookIndex@T0113\endcsname{cedu@测度（measure）}
\expandafter\let\csname BookOriginal@trMeasure\endcsname\trMeasure
\renewcommand{\trMeasure}{\BookTerm{T0113}{\csname BookOriginal@trMeasure\endcsname}}
\expandafter\def\csname BookIndex@T0114\endcsname{baocedu@保测度（measure preserving）}
\expandafter\let\csname BookOriginal@trMeasurePreserving\endcsname\trMeasurePreserving
\renewcommand{\trMeasurePreserving}{\BookTerm{T0114}{\csname BookOriginal@trMeasurePreserving\endcsname}}
\expandafter\def\csname BookIndex@T0115\endcsname{lingce@零测（measure zero）}
\expandafter\let\csname BookOriginal@trMeasureZero\endcsname\trMeasureZero
\renewcommand{\trMeasureZero}{\BookTerm{T0115}{\csname BookOriginal@trMeasureZero\endcsname}}
\expandafter\def\csname BookIndex@T0116\endcsname{wending@稳定（stable）}
\expandafter\let\csname BookOriginal@trStable\endcsname\trStable
\renewcommand{\trStable}{\BookTerm{T0116}{\csname BookOriginal@trStable\endcsname}}
\expandafter\def\csname BookIndex@T0117\endcsname{wendingliuxing@稳定流形（stable manifold）}
\expandafter\let\csname BookOriginal@trStableManifold\endcsname\trStableManifold
\renewcommand{\trStableManifold}{\BookTerm{T0117}{\csname BookOriginal@trStableManifold\endcsname}}
\expandafter\def\csname BookIndex@T0118\endcsname{buwendingliuxing@不稳定流形（unstable manifold）}
\expandafter\let\csname BookOriginal@trUnstableManifold\endcsname\trUnstableManifold
\renewcommand{\trUnstableManifold}{\BookTerm{T0118}{\csname BookOriginal@trUnstableManifold\endcsname}}
\expandafter\def\csname BookIndex@T0119\endcsname{lianjieguidao@连接轨道（connecting orbit）}
\expandafter\let\csname BookOriginal@trConnectingOrbit\endcsname\trConnectingOrbit
\renewcommand{\trConnectingOrbit}{\BookTerm{T0119}{\csname BookOriginal@trConnectingOrbit\endcsname}}
\expandafter\def\csname BookIndex@T0120\endcsname{tongsu@同宿（homoclinic）}
\expandafter\let\csname BookOriginal@trHomoclinic\endcsname\trHomoclinic
\renewcommand{\trHomoclinic}{\BookTerm{T0120}{\csname BookOriginal@trHomoclinic\endcsname}}
\expandafter\def\csname BookIndex@T0121\endcsname{bihe@闭合（closing）}
\expandafter\let\csname BookOriginal@trClosing\endcsname\trClosing
\renewcommand{\trClosing}{\BookTerm{T0121}{\csname BookOriginal@trClosing\endcsname}}
\expandafter\def\csname BookIndex@T0122\endcsname{biheyinli@闭合引理（closing lemma）}
\expandafter\let\csname BookOriginal@trClosingLemma\endcsname\trClosingLemma
\renewcommand{\trClosingLemma}{\BookTerm{T0122}{\csname BookOriginal@trClosingLemma\endcsname}}
\expandafter\def\csname BookIndex@T0123\endcsname{feibiheyinli@非闭合引理（non-closing lemma）}
\expandafter\let\csname BookOriginal@trNonClosingLemma\endcsname\trNonClosingLemma
\renewcommand{\trNonClosingLemma}{\BookTerm{T0123}{\csname BookOriginal@trNonClosingLemma\endcsname}}
\expandafter\def\csname BookIndex@T0124\endcsname{bihewenti@闭合问题（closing problem）}
\expandafter\let\csname BookOriginal@trClosingProblem\endcsname\trClosingProblem
\renewcommand{\trClosingProblem}{\BookTerm{T0124}{\csname BookOriginal@trClosingProblem\endcsname}}
\expandafter\def\csname BookIndex@T0125\endcsname{diufantu@丢番图（diophantine）}
\expandafter\let\csname BookOriginal@trDiophantine\endcsname\trDiophantine
\renewcommand{\trDiophantine}{\BookTerm{T0125}{\csname BookOriginal@trDiophantine\endcsname}}
\expandafter\def\csname BookIndex@T0126\endcsname{wuli@无理（irrational）}
\expandafter\let\csname BookOriginal@trIrrational\endcsname\trIrrational
\renewcommand{\trIrrational}{\BookTerm{T0126}{\csname BookOriginal@trIrrational\endcsname}}
\expandafter\def\csname BookIndex@T0127\endcsname{pingyi@平移（translation）}
\expandafter\let\csname BookOriginal@trTranslation\endcsname\trTranslation
\renewcommand{\trTranslation}{\BookTerm{T0127}{\csname BookOriginal@trTranslation\endcsname}}
\expandafter\def\csname BookIndex@T0128\endcsname{donglixue@动力学（dynamics）}
\expandafter\let\csname BookOriginal@trDynamics\endcsname\trDynamics
\renewcommand{\trDynamics}{\BookTerm{T0128}{\csname BookOriginal@trDynamics\endcsname}}
\expandafter\def\csname BookIndex@T0129\endcsname{wanbeihua@完备化（completion）}
\expandafter\let\csname BookOriginal@trCompletion\endcsname\trCompletion
\renewcommand{\trCompletion}{\BookTerm{T0129}{\csname BookOriginal@trCompletion\endcsname}}
\expandafter\def\csname BookIndex@T0130\endcsname{jiaohuanzi@交换子（commutator）}
\expandafter\let\csname BookOriginal@trCommutator\endcsname\trCommutator
\renewcommand{\trCommutator}{\BookTerm{T0130}{\csname BookOriginal@trCommutator\endcsname}}
\expandafter\def\csname BookIndex@T0131\endcsname{weiduliang@伪度量（pseudo-metric）}
\expandafter\let\csname BookOriginal@trPseudoMetric\endcsname\trPseudoMetric
\renewcommand{\trPseudoMetric}{\BookTerm{T0131}{\csname BookOriginal@trPseudoMetric\endcsname}}
\expandafter\def\csname BookIndex@T0132\endcsname{tapuqun@拓扑群（topological group）}
\expandafter\let\csname BookOriginal@trTopologicalGroup\endcsname\trTopologicalGroup
\renewcommand{\trTopologicalGroup}{\BookTerm{T0132}{\csname BookOriginal@trTopologicalGroup\endcsname}}
\expandafter\def\csname BookIndex@T0133\endcsname{xujiegou@序结构（order structure）}
\expandafter\let\csname BookOriginal@trOrderStructure\endcsname\trOrderStructure
\renewcommand{\trOrderStructure}{\BookTerm{T0133}{\csname BookOriginal@trOrderStructure\endcsname}}
\expandafter\def\csname BookIndex@T0134\endcsname{pianxu@偏序（partial order）}
\expandafter\let\csname BookOriginal@trPartialOrder\endcsname\trPartialOrder
\renewcommand{\trPartialOrder}{\BookTerm{T0134}{\csname BookOriginal@trPartialOrder\endcsname}}
\expandafter\def\csname BookIndex@T0135\endcsname{shuangshe@双射（bijection）}
\expandafter\let\csname BookOriginal@trBijection\endcsname\trBijection
\renewcommand{\trBijection}{\BookTerm{T0135}{\csname BookOriginal@trBijection\endcsname}}
\expandafter\def\csname BookIndex@T0136\endcsname{yuanxingshi@原形式（primitive）}
\expandafter\let\csname BookOriginal@trPrimitive\endcsname\trPrimitive
\renewcommand{\trPrimitive}{\BookTerm{T0136}{\csname BookOriginal@trPrimitive\endcsname}}
\expandafter\def\csname BookIndex@T0137\endcsname{fenbujifen@分部积分（integration by parts）}
\expandafter\let\csname BookOriginal@trIntegrationByParts\endcsname\trIntegrationByParts
\renewcommand{\trIntegrationByParts}{\BookTerm{T0137}{\csname BookOriginal@trIntegrationByParts\endcsname}}
\expandafter\def\csname BookIndex@T0138\endcsname{longqihanshu@隆起函数（bump function）}
\expandafter\let\csname BookOriginal@trBumpFunction\endcsname\trBumpFunction
\renewcommand{\trBumpFunction}{\BookTerm{T0138}{\csname BookOriginal@trBumpFunction\endcsname}}
\expandafter\def\csname BookIndex@T0139\endcsname{dandiao@单调（monotone）}
\expandafter\let\csname BookOriginal@trMonotone\endcsname\trMonotone
\renewcommand{\trMonotone}{\BookTerm{T0139}{\csname BookOriginal@trMonotone\endcsname}}
\expandafter\def\csname BookIndex@T0140\endcsname{wanyou@万有（universal）}
\expandafter\let\csname BookOriginal@trUniversal\endcsname\trUniversal
\renewcommand{\trUniversal}{\BookTerm{T0140}{\csname BookOriginal@trUniversal\endcsname}}
\expandafter\def\csname BookIndex@T0141\endcsname{teshurongliang@特殊容量（special capacity）}
\expandafter\let\csname BookOriginal@trSpecialCapacity\endcsname\trSpecialCapacity
\renewcommand{\trSpecialCapacity}{\BookTerm{T0141}{\csname BookOriginal@trSpecialCapacity\endcsname}}
\expandafter\def\csname BookIndex@T0142\endcsname{buxiangjiao@不相交（disjoint）}
\expandafter\let\csname BookOriginal@trDisjoint\endcsname\trDisjoint
\renewcommand{\trDisjoint}{\BookTerm{T0142}{\csname BookOriginal@trDisjoint\endcsname}}
\expandafter\def\csname BookIndex@T0143\endcsname{zhengchangfenli@正常分离（properly separated）}
\expandafter\let\csname BookOriginal@trProperlySeparated\endcsname\trProperlySeparated
\renewcommand{\trProperlySeparated}{\BookTerm{T0143}{\csname BookOriginal@trProperlySeparated\endcsname}}
\expandafter\def\csname BookIndex@T0144\endcsname{tubao@凸包（convex hull）}
\expandafter\let\csname BookOriginal@trConvexHull\endcsname\trConvexHull
\renewcommand{\trConvexHull}{\BookTerm{T0144}{\csname BookOriginal@trConvexHull\endcsname}}
\expandafter\def\csname BookIndex@T0145\endcsname{nizizhi@拟自治（quasi autonomous）}
\expandafter\let\csname BookOriginal@trQuasiAutonomous\endcsname\trQuasiAutonomous
\renewcommand{\trQuasiAutonomous}{\BookTerm{T0145}{\csname BookOriginal@trQuasiAutonomous\endcsname}}
\expandafter\def\csname BookIndex@T0146\endcsname{jixiao@极小（minimal）}
\expandafter\let\csname BookOriginal@trMinimal\endcsname\trMinimal
\renewcommand{\trMinimal}{\BookTerm{T0146}{\csname BookOriginal@trMinimal\endcsname}}
\expandafter\def\csname BookIndex@T0147\endcsname{jiandan@简单（simple）}
\expandafter\let\csname BookOriginal@trSimple\endcsname\trSimple
\renewcommand{\trSimple}{\BookTerm{T0147}{\csname BookOriginal@trSimple\endcsname}}
\expandafter\def\csname BookIndex@T0148\endcsname{jiandanxing@简单性（simplicity）}
\expandafter\let\csname BookOriginal@trSimplicity\endcsname\trSimplicity
\renewcommand{\trSimplicity}{\BookTerm{T0148}{\csname BookOriginal@trSimplicity\endcsname}}
\expandafter\def\csname BookIndex@T0149\endcsname{beizhang@杯长（cup-length）}
\expandafter\let\csname BookOriginal@trCupLength\endcsname\trCupLength
\renewcommand{\trCupLength}{\BookTerm{T0149}{\csname BookOriginal@trCupLength\endcsname}}
\expandafter\def\csname BookIndex@T0150\endcsname{houan@猴鞍（monkey saddle）}
\expandafter\let\csname BookOriginal@trMonkeySaddle\endcsname\trMonkeySaddle
\renewcommand{\trMonkeySaddle}{\BookTerm{T0150}{\csname BookOriginal@trMonkeySaddle\endcsname}}
\expandafter\def\csname BookIndex@T0151\endcsname{zhixin@质心（center of mass）}
\expandafter\let\csname BookOriginal@trCenterOfMass\endcsname\trCenterOfMass
\renewcommand{\trCenterOfMass}{\BookTerm{T0151}{\csname BookOriginal@trCenterOfMass\endcsname}}
\expandafter\def\csname BookIndex@T0152\endcsname{jiequlu@截曲率（sectional curvature）}
\expandafter\let\csname BookOriginal@trSectionalCurvature\endcsname\trSectionalCurvature
\renewcommand{\trSectionalCurvature}{\BookTerm{T0152}{\csname BookOriginal@trSectionalCurvature\endcsname}}
\expandafter\def\csname BookIndex@T0153\endcsname{maoji@帽积（cap product）}
\expandafter\let\csname BookOriginal@trCapProduct\endcsname\trCapProduct
\renewcommand{\trCapProduct}{\BookTerm{T0153}{\csname BookOriginal@trCapProduct\endcsname}}
\expandafter\def\csname BookIndex@T0154\endcsname{qiadangyingshe@恰当映射（proper map）}
\expandafter\let\csname BookOriginal@trProperMap\endcsname\trProperMap
\renewcommand{\trProperMap}{\BookTerm{T0154}{\csname BookOriginal@trProperMap\endcsname}}
\expandafter\def\csname BookIndex@T0155\endcsname{lingjiemian@零截面（zero section）}
\expandafter\let\csname BookOriginal@trZeroSection\endcsname\trZeroSection
\renewcommand{\trZeroSection}{\BookTerm{T0155}{\csname BookOriginal@trZeroSection\endcsname}}
\expandafter\def\csname BookIndex@T0156\endcsname{jibenlei@基本类（fundamental class）}
\expandafter\let\csname BookOriginal@trFundamentalClass\endcsname\trFundamentalClass
\renewcommand{\trFundamentalClass}{\BookTerm{T0156}{\csname BookOriginal@trFundamentalClass\endcsname}}
\expandafter\def\csname BookIndex@T0157\endcsname{yuehua@约化（reduction）}
\expandafter\let\csname BookOriginal@trReduction\endcsname\trReduction
\renewcommand{\trReduction}{\BookTerm{T0157}{\csname BookOriginal@trReduction\endcsname}}
\expandafter\def\csname BookIndex@T0158\endcsname{xiangweihanshu@相位函数（phase function）}
\expandafter\let\csname BookOriginal@trPhaseFunction\endcsname\trPhaseFunction
\renewcommand{\trPhaseFunction}{\BookTerm{T0158}{\csname BookOriginal@trPhaseFunction\endcsname}}
\expandafter\def\csname BookIndex@T0159\endcsname{fuliyejifensuanzi@傅里叶积分算子（Fourier integral operator）}
\expandafter\let\csname BookOriginal@trFourierIntegralOperator\endcsname\trFourierIntegralOperator
\renewcommand{\trFourierIntegralOperator}{\BookTerm{T0159}{\csname BookOriginal@trFourierIntegralOperator\endcsname}}
\expandafter\def\csname BookIndex@T0160\endcsname{fenlie@分裂（splitting）}
\expandafter\let\csname BookOriginal@trSplitting\endcsname\trSplitting
\renewcommand{\trSplitting}{\BookTerm{T0160}{\csname BookOriginal@trSplitting\endcsname}}
\expandafter\def\csname BookIndex@T0161\endcsname{dengdulianxu@等度连续（equi-continuous）}
\expandafter\let\csname BookOriginal@trEquicontinuous\endcsname\trEquicontinuous
\renewcommand{\trEquicontinuous}{\BookTerm{T0161}{\csname BookOriginal@trEquicontinuous\endcsname}}
\expandafter\def\csname BookIndex@T0162\endcsname{yujin@预紧（precompact）}
\expandafter\let\csname BookOriginal@trPrecompact\endcsname\trPrecompact
\renewcommand{\trPrecompact}{\BookTerm{T0162}{\csname BookOriginal@trPrecompact\endcsname}}
\expandafter\def\csname BookIndex@T0163\endcsname{jixianji@极限集（limit set）}
\expandafter\let\csname BookOriginal@trLimitSet\endcsname\trLimitSet
\renewcommand{\trLimitSet}{\BookTerm{T0163}{\csname BookOriginal@trLimitSet\endcsname}}
\expandafter\def\csname BookIndex@T0164\endcsname{banliu@半流（semiflow）}
\expandafter\let\csname BookOriginal@trSemiflow\endcsname\trSemiflow
\renewcommand{\trSemiflow}{\BookTerm{T0164}{\csname BookOriginal@trSemiflow\endcsname}}
\expandafter\def\csname BookIndex@T0165\endcsname{lianxutong@连续统（continuum）}
\expandafter\let\csname BookOriginal@trContinuum\endcsname\trContinuum
\renewcommand{\trContinuum}{\BookTerm{T0165}{\csname BookOriginal@trContinuum\endcsname}}
\expandafter\def\csname BookIndex@T0166\endcsname{zhouqiqun@周期群（period group）}
\expandafter\let\csname BookOriginal@trPeriodGroup\endcsname\trPeriodGroup
\renewcommand{\trPeriodGroup}{\BookTerm{T0166}{\csname BookOriginal@trPeriodGroup\endcsname}}
\expandafter\def\csname BookIndex@T0167\endcsname{maopian@帽片（cap）}
\expandafter\let\csname BookOriginal@trCapDisk\endcsname\trCapDisk
\renewcommand{\trCapDisk}{\BookTerm{T0167}{\csname BookOriginal@trCapDisk\endcsname}}
\expandafter\def\csname BookIndex@T0168\endcsname{zhanhefangfa@粘合方法（gluing method）}
\expandafter\let\csname BookOriginal@trGluingMethod\endcsname\trGluingMethod
\renewcommand{\trGluingMethod}{\BookTerm{T0168}{\csname BookOriginal@trGluingMethod\endcsname}}
\expandafter\def\csname BookIndex@T0169\endcsname{polieguiji@破裂轨迹（broken trajectory）}
\expandafter\let\csname BookOriginal@trBrokenTrajectory\endcsname\trBrokenTrajectory
\renewcommand{\trBrokenTrajectory}{\BookTerm{T0169}{\csname BookOriginal@trBrokenTrajectory\endcsname}}
\expandafter\def\csname BookIndex@T0170\endcsname{liantongdanxi@连通单系（connected simple system）}
\expandafter\let\csname BookOriginal@trConnectedSimpleSystem\endcsname\trConnectedSimpleSystem
\renewcommand{\trConnectedSimpleSystem}{\BookTerm{T0170}{\csname BookOriginal@trConnectedSimpleSystem\endcsname}}
\expandafter\def\csname BookIndex@T0171\endcsname{mokongjian@模空间（moduli space）}
\expandafter\let\csname BookOriginal@trModuliSpace\endcsname\trModuliSpace
\renewcommand{\trModuliSpace}{\BookTerm{T0171}{\csname BookOriginal@trModuliSpace\endcsname}}
\expandafter\def\csname BookIndex@T0172\endcsname{pingfanhua@平凡化（trivialization）}
\expandafter\let\csname BookOriginal@trTrivialization\endcsname\trTrivialization
\renewcommand{\trTrivialization}{\BookTerm{T0172}{\csname BookOriginal@trTrivialization\endcsname}}
\expandafter\def\csname BookIndex@T0173\endcsname{shunzitongdiao@瞬子同调（instanton homology）}
\expandafter\let\csname BookOriginal@trInstantonHomology\endcsname\trInstantonHomology
\renewcommand{\trInstantonHomology}{\BookTerm{T0173}{\csname BookOriginal@trInstantonHomology\endcsname}}
\expandafter\def\csname BookIndex@T0174\endcsname{lianluo@联络（connection）}
\expandafter\let\csname BookOriginal@trConnection\endcsname\trConnection
\renewcommand{\trConnection}{\BookTerm{T0174}{\csname BookOriginal@trConnection\endcsname}}
\expandafter\def\csname BookIndex@T0175\endcsname{huanraoshu@环绕数（winding number）}
\expandafter\let\csname BookOriginal@trWindingNumber\endcsname\trWindingNumber
\renewcommand{\trWindingNumber}{\BookTerm{T0175}{\csname BookOriginal@trWindingNumber\endcsname}}
\expandafter\def\csname BookIndex@T0176\endcsname{liantongtai@链同态（chain map）}
\expandafter\let\csname BookOriginal@trChainMap\endcsname\trChainMap
\renewcommand{\trChainMap}{\BookTerm{T0176}{\csname BookOriginal@trChainMap\endcsname}}
\expandafter\def\csname BookIndex@T0177\endcsname{xiangjiao@相交（intersection）}
\expandafter\let\csname BookOriginal@trIntersection\endcsname\trIntersection
\renewcommand{\trIntersection}{\BookTerm{T0177}{\csname BookOriginal@trIntersection\endcsname}}
\expandafter\def\csname BookIndex@T0178\endcsname{xishuyu@系数域（coefficient field）}
\expandafter\let\csname BookOriginal@trCoefficientField\endcsname\trCoefficientField
\renewcommand{\trCoefficientField}{\BookTerm{T0178}{\csname BookOriginal@trCoefficientField\endcsname}}
\expandafter\def\csname BookIndex@T0179\endcsname{daishutapu@代数拓扑（algebraic topology）}
\expandafter\let\csname BookOriginal@trAlgebraicTopology\endcsname\trAlgebraicTopology
\renewcommand{\trAlgebraicTopology}{\BookTerm{T0179}{\csname BookOriginal@trAlgebraicTopology\endcsname}}
\expandafter\def\csname BookIndex@T0180\endcsname{luguo@滤过（filtration）}
\expandafter\let\csname BookOriginal@trFiltration\endcsname\trFiltration
\renewcommand{\trFiltration}{\BookTerm{T0180}{\csname BookOriginal@trFiltration\endcsname}}
\expandafter\def\csname BookIndex@T0181\endcsname{youxiangji@有向集（directed set）}
\expandafter\let\csname BookOriginal@trDirectedSet\endcsname\trDirectedSet
\renewcommand{\trDirectedSet}{\BookTerm{T0181}{\csname BookOriginal@trDirectedSet\endcsname}}
\expandafter\def\csname BookIndex@T0182\endcsname{zhijixian@直极限（direct limit）}
\expandafter\let\csname BookOriginal@trDirectedLimit\endcsname\trDirectedLimit
\renewcommand{\trDirectedLimit}{\BookTerm{T0182}{\csname BookOriginal@trDirectedLimit\endcsname}}
\expandafter\def\csname BookIndex@T0183\endcsname{weitu@伪凸（pseudo convex）}
\expandafter\let\csname BookOriginal@trPseudoConvex\endcsname\trPseudoConvex
\renewcommand{\trPseudoConvex}{\BookTerm{T0183}{\csname BookOriginal@trPseudoConvex\endcsname}}
\expandafter\def\csname BookIndex@T0184\endcsname{duoyuanpan@多圆盘（polydisc）}
\expandafter\let\csname BookOriginal@trPolydisc\endcsname\trPolydisc
\renewcommand{\trPolydisc}{\BookTerm{T0184}{\csname BookOriginal@trPolydisc\endcsname}}
\expandafter\def\csname BookIndex@T0185\endcsname{tianzhuang@填装（packing）}
\expandafter\let\csname BookOriginal@trPacking\endcsname\trPacking
\renewcommand{\trPacking}{\BookTerm{T0185}{\csname BookOriginal@trPacking\endcsname}}
\expandafter\def\csname BookIndex@T0186\endcsname{pingtan@平坦（flat）}
\expandafter\let\csname BookOriginal@trFlat\endcsname\trFlat
\renewcommand{\trFlat}{\BookTerm{T0186}{\csname BookOriginal@trFlat\endcsname}}
\expandafter\def\csname BookIndex@T0187\endcsname{fenchatu@分岔图（bifurcation diagram）}
\expandafter\let\csname BookOriginal@trBifurcationDiagram\endcsname\trBifurcationDiagram
\renewcommand{\trBifurcationDiagram}{\BookTerm{T0187}{\csname BookOriginal@trBifurcationDiagram\endcsname}}
\expandafter\def\csname BookIndex@T0188\endcsname{zhongyanxingshi@重言形式（tautological form）}
\expandafter\let\csname BookOriginal@trTautologicalForm\endcsname\trTautologicalForm
\renewcommand{\trTautologicalForm}{\BookTerm{T0188}{\csname BookOriginal@trTautologicalForm\endcsname}}
\expandafter\def\csname BookIndex@T0189\endcsname{bixingshi@闭形式（closed form）}
\expandafter\let\csname BookOriginal@trClosedForm\endcsname\trClosedForm
\renewcommand{\trClosedForm}{\BookTerm{T0189}{\csname BookOriginal@trClosedForm\endcsname}}
\expandafter\def\csname BookIndex@T0190\endcsname{shacheguidao@刹车轨道（brake orbit）}
\expandafter\let\csname BookOriginal@trBrakeOrbit\endcsname\trBrakeOrbit
\renewcommand{\trBrakeOrbit}{\BookTerm{T0190}{\csname BookOriginal@trBrakeOrbit\endcsname}}
\expandafter\def\csname BookIndex@T0191\endcsname{dongneng@动能（kinetic energy）}
\expandafter\let\csname BookOriginal@trKineticEnergy\endcsname\trKineticEnergy
\renewcommand{\trKineticEnergy}{\BookTerm{T0191}{\csname BookOriginal@trKineticEnergy\endcsname}}
\expandafter\def\csname BookIndex@T0192\endcsname{shineng@势能（potential energy）}
\expandafter\let\csname BookOriginal@trPotentialEnergy\endcsname\trPotentialEnergy
\renewcommand{\trPotentialEnergy}{\BookTerm{T0192}{\csname BookOriginal@trPotentialEnergy\endcsname}}
\expandafter\def\csname BookIndex@T0193\endcsname{baoshouxitong@保守系统（conservative system）}
\expandafter\let\csname BookOriginal@trConservativeSystem\endcsname\trConservativeSystem
\renewcommand{\trConservativeSystem}{\BookTerm{T0193}{\csname BookOriginal@trConservativeSystem\endcsname}}
\expandafter\def\csname BookIndex@T0194\endcsname{ziyoudu@自由度（degrees of freedom）}
\expandafter\let\csname BookOriginal@trDegreesOfFreedom\endcsname\trDegreesOfFreedom
\renewcommand{\trDegreesOfFreedom}{\BookTerm{T0194}{\csname BookOriginal@trDegreesOfFreedom\endcsname}}
\expandafter\def\csname BookIndex@T0195\endcsname{zhenfu@振幅（amplitude）}
\expandafter\let\csname BookOriginal@trAmplitude\endcsname\trAmplitude
\renewcommand{\trAmplitude}{\BookTerm{T0195}{\csname BookOriginal@trAmplitude\endcsname}}
\expandafter\def\csname BookIndex@T0196\endcsname{xinfenlie@辛分裂（symplectic splitting）}
\expandafter\let\csname BookOriginal@trSymplecticSplitting\endcsname\trSymplecticSplitting
\renewcommand{\trSymplecticSplitting}{\BookTerm{T0196}{\csname BookOriginal@trSymplecticSplitting\endcsname}}
\expandafter\def\csname BookIndex@T0197\endcsname{tongpei@同胚（homeomorphic）}
\expandafter\let\csname BookOriginal@trHomeomorphic\endcsname\trHomeomorphic
\renewcommand{\trHomeomorphic}{\BookTerm{T0197}{\csname BookOriginal@trHomeomorphic\endcsname}}
\expandafter\def\csname BookIndex@T0198\endcsname{baoqiang@胞腔（cell）}
\expandafter\let\csname BookOriginal@trCell\endcsname\trCell
\renewcommand{\trCell}{\BookTerm{T0198}{\csname BookOriginal@trCell\endcsname}}
\expandafter\def\csname BookIndex@T0199\endcsname{tiantilixue@天体力学（celestial mechanics）}
\expandafter\let\csname BookOriginal@trCelestialMechanics\endcsname\trCelestialMechanics
\renewcommand{\trCelestialMechanics}{\BookTerm{T0199}{\csname BookOriginal@trCelestialMechanics\endcsname}}
\expandafter\def\csname BookIndex@T0200\endcsname{ertiwenti@二体问题（two-body problem）}
\expandafter\let\csname BookOriginal@trTwoBodyProblem\endcsname\trTwoBodyProblem
\renewcommand{\trTwoBodyProblem}{\BookTerm{T0200}{\csname BookOriginal@trTwoBodyProblem\endcsname}}
\expandafter\def\csname BookIndex@T0201\endcsname{xianzhixingsantiwenti@限制性三体问题（restricted three-body problem）}
\expandafter\let\csname BookOriginal@trRestrictedThreeBodyProblem\endcsname\trRestrictedThreeBodyProblem
\renewcommand{\trRestrictedThreeBodyProblem}{\BookTerm{T0201}{\csname BookOriginal@trRestrictedThreeBodyProblem\endcsname}}
\expandafter\def\csname BookIndex@T0202\endcsname{pengzhuang@碰撞（collision）}
\expandafter\let\csname BookOriginal@trCollision\endcsname\trCollision
\renewcommand{\trCollision}{\BookTerm{T0202}{\csname BookOriginal@trCollision\endcsname}}
\expandafter\def\csname BookIndex@T0203\endcsname{qidian@奇点（singularity）}
\expandafter\let\csname BookOriginal@trSingularity\endcsname\trSingularity
\renewcommand{\trSingularity}{\BookTerm{T0203}{\csname BookOriginal@trSingularity\endcsname}}
\expandafter\def\csname BookIndex@T0204\endcsname{jingdianlixue@经典力学（classical mechanics）}
\expandafter\let\csname BookOriginal@trClassicalMechanics\endcsname\trClassicalMechanics
\renewcommand{\trClassicalMechanics}{\BookTerm{T0204}{\csname BookOriginal@trClassicalMechanics\endcsname}}
\expandafter\def\csname BookIndex@T0205\endcsname{kanglizhishu@康利指数（Conley index）}
\expandafter\let\csname BookOriginal@trConleyIndex\endcsname\trConleyIndex
\renewcommand{\trConleyIndex}{\BookTerm{T0205}{\csname BookOriginal@trConleyIndex\endcsname}}
\expandafter\def\csname BookIndex@T0206\endcsname{gulibubianji@孤立不变集（isolated invariant set）}
\expandafter\let\csname BookOriginal@trIsolatedInvariantSet\endcsname\trIsolatedInvariantSet
\renewcommand{\trIsolatedInvariantSet}{\BookTerm{T0206}{\csname BookOriginal@trIsolatedInvariantSet\endcsname}}
\expandafter\def\csname BookIndex@T0207\endcsname{yanta@延拓（continuation）}
\expandafter\let\csname BookOriginal@trContinuation\endcsname\trContinuation
\renewcommand{\trContinuation}{\BookTerm{T0207}{\csname BookOriginal@trContinuation\endcsname}}
\expandafter\def\csname BookIndex@T0208\endcsname{raodong@扰动（perturbation）}
\expandafter\let\csname BookOriginal@trPerturbation\endcsname\trPerturbation
\renewcommand{\trPerturbation}{\BookTerm{T0208}{\csname BookOriginal@trPerturbation\endcsname}}
\expandafter\def\csname BookIndex@T0209\endcsname{fencha@分岔（bifurcation）}
\expandafter\let\csname BookOriginal@trBifurcation\endcsname\trBifurcation
\renewcommand{\trBifurcation}{\BookTerm{T0209}{\csname BookOriginal@trBifurcation\endcsname}}
\expandafter\def\csname BookIndex@T0210\endcsname{zhendang@振荡（oscillation）}
\expandafter\let\csname BookOriginal@trOscillation\endcsname\trOscillation
\renewcommand{\trOscillation}{\BookTerm{T0210}{\csname BookOriginal@trOscillation\endcsname}}
\expandafter\def\csname BookIndex@T0211\endcsname{shoupozhendang@受迫振荡（forced oscillation）}
\expandafter\let\csname BookOriginal@trForcedOscillation\endcsname\trForcedOscillation
\renewcommand{\trForcedOscillation}{\BookTerm{T0211}{\csname BookOriginal@trForcedOscillation\endcsname}}
\expandafter\def\csname BookIndex@T0212\endcsname{zhudian@驻点（rest point）}
\expandafter\let\csname BookOriginal@trRestPoint\endcsname\trRestPoint
\renewcommand{\trRestPoint}{\BookTerm{T0212}{\csname BookOriginal@trRestPoint\endcsname}}
\expandafter\def\csname BookIndex@T0213\endcsname{kejixitong@可积系统（integrable system）}
\expandafter\let\csname BookOriginal@trIntegrableSystem\endcsname\trIntegrableSystem
\renewcommand{\trIntegrableSystem}{\BookTerm{T0213}{\csname BookOriginal@trIntegrableSystem\endcsname}}
\expandafter\def\csname BookIndex@T0214\endcsname{zuai@阻碍（obstruction）}
\expandafter\let\csname BookOriginal@trObstruction\endcsname\trObstruction
\renewcommand{\trObstruction}{\BookTerm{T0214}{\csname BookOriginal@trObstruction\endcsname}}
\expandafter\def\csname BookIndex@T0215\endcsname{chuzhiwenti@初值问题（initial value problem）}
\expandafter\let\csname BookOriginal@trInitialValueProblem\endcsname\trInitialValueProblem
\renewcommand{\trInitialValueProblem}{\BookTerm{T0215}{\csname BookOriginal@trInitialValueProblem\endcsname}}
\expandafter\def\csname BookIndex@T0216\endcsname{weifenfangcheng@微分方程（differential equation）}
\expandafter\let\csname BookOriginal@trDifferentialEquation\endcsname\trDifferentialEquation
\renewcommand{\trDifferentialEquation}{\BookTerm{T0216}{\csname BookOriginal@trDifferentialEquation\endcsname}}
\expandafter\def\csname BookIndex@T0217\endcsname{pianweifenfangcheng@偏微分方程（partial differential equation）}
\expandafter\let\csname BookOriginal@trPartialDifferentialEquation\endcsname\trPartialDifferentialEquation
\renewcommand{\trPartialDifferentialEquation}{\BookTerm{T0217}{\csname BookOriginal@trPartialDifferentialEquation\endcsname}}
\expandafter\def\csname BookIndex@T0218\endcsname{shijianyilai@时间依赖（time-dependent）}
\expandafter\let\csname BookOriginal@trTimeDependent\endcsname\trTimeDependent
\renewcommand{\trTimeDependent}{\BookTerm{T0218}{\csname BookOriginal@trTimeDependent\endcsname}}
\expandafter\def\csname BookIndex@T0219\endcsname{shijian1yingshe@时间-1映射（time-1 map）}
\expandafter\let\csname BookOriginal@trTimeOneMap\endcsname\trTimeOneMap
\renewcommand{\trTimeOneMap}{\BookTerm{T0219}{\csname BookOriginal@trTimeOneMap\endcsname}}
\expandafter\def\csname BookIndex@T0220\endcsname{rongxu@容许（admissible）}
\expandafter\let\csname BookOriginal@trAdmissible\endcsname\trAdmissible
\renewcommand{\trAdmissible}{\BookTerm{T0220}{\csname BookOriginal@trAdmissible\endcsname}}
\expandafter\def\csname BookIndex@T0221\endcsname{bianhuanfaze@变换法则（transformation law）}
\expandafter\let\csname BookOriginal@trTransformationLaw\endcsname\trTransformationLaw
\renewcommand{\trTransformationLaw}{\BookTerm{T0221}{\csname BookOriginal@trTransformationLaw\endcsname}}
\expandafter\def\csname BookIndex@T0222\endcsname{bianhuan@变换（transformation）}
\expandafter\let\csname BookOriginal@trTransformation\endcsname\trTransformation
\renewcommand{\trTransformation}{\BookTerm{T0222}{\csname BookOriginal@trTransformation\endcsname}}
\expandafter\def\csname BookIndex@T0223\endcsname{chudeng@初等（elementary）}
\expandafter\let\csname BookOriginal@trElementary\endcsname\trElementary
\renewcommand{\trElementary}{\BookTerm{T0223}{\csname BookOriginal@trElementary\endcsname}}
\expandafter\def\csname BookIndex@T0224\endcsname{zhihuan@置换（permutation）}
\expandafter\let\csname BookOriginal@trPermutation\endcsname\trPermutation
\renewcommand{\trPermutation}{\BookTerm{T0224}{\csname BookOriginal@trPermutation\endcsname}}
\expandafter\def\csname BookIndex@T0225\endcsname{shuzhihanshu@数值函数（scalar function）}
\expandafter\let\csname BookOriginal@trScalarFunction\endcsname\trScalarFunction
\renewcommand{\trScalarFunction}{\BookTerm{T0225}{\csname BookOriginal@trScalarFunction\endcsname}}
\expandafter\def\csname BookIndex@T0226\endcsname{ge@格（lattice）}
\expandafter\let\csname BookOriginal@trLattice\endcsname\trLattice
\renewcommand{\trLattice}{\BookTerm{T0226}{\csname BookOriginal@trLattice\endcsname}}
\expandafter\def\csname BookIndex@T0227\endcsname{mixiangqun@迷向群（isotropy group）}
\expandafter\let\csname BookOriginal@trIsotropyGroup\endcsname\trIsotropyGroup
\renewcommand{\trIsotropyGroup}{\BookTerm{T0227}{\csname BookOriginal@trIsotropyGroup\endcsname}}
\expandafter\def\csname BookIndex@T0228\endcsname{fenshujie@分数阶（fractional）}
\expandafter\let\csname BookOriginal@trFractional\endcsname\trFractional
\renewcommand{\trFractional}{\BookTerm{T0228}{\csname BookOriginal@trFractional\endcsname}}
\expandafter\def\csname BookIndex@T0229\endcsname{yaomo@幺模（unimodular）}
\expandafter\let\csname BookOriginal@trUnimodular\endcsname\trUnimodular
\renewcommand{\trUnimodular}{\BookTerm{T0229}{\csname BookOriginal@trUnimodular\endcsname}}
\expandafter\def\csname BookIndex@T0230\endcsname{xiangyi@相移（phase shift）}
\expandafter\let\csname BookOriginal@trPhaseShift\endcsname\trPhaseShift
\renewcommand{\trPhaseShift}{\BookTerm{T0230}{\csname BookOriginal@trPhaseShift\endcsname}}
\expandafter\def\csname BookIndex@T0231\endcsname{danweifenjie@单位分解（partition of unity）}
\expandafter\let\csname BookOriginal@trPartitionOfUnity\endcsname\trPartitionOfUnity
\renewcommand{\trPartitionOfUnity}{\BookTerm{T0231}{\csname BookOriginal@trPartitionOfUnity\endcsname}}
\expandafter\def\csname BookIndex@T0232\endcsname{zuobiaotuce@坐标图册（atlas）}
\expandafter\let\csname BookOriginal@trAtlas\endcsname\trAtlas
\renewcommand{\trAtlas}{\BookTerm{T0232}{\csname BookOriginal@trAtlas\endcsname}}
\expandafter\def\csname BookIndex@T0233\endcsname{gelinhanshu@格林函数（Green's function）}
\expandafter\let\csname BookOriginal@trGreenFunction\endcsname\trGreenFunction
\renewcommand{\trGreenFunction}{\BookTerm{T0233}{\csname BookOriginal@trGreenFunction\endcsname}}
\expandafter\def\csname BookIndex@T0234\endcsname{shangquejiefanshu@上确界范数（sup-norm）}
\expandafter\let\csname BookOriginal@trSupNorm\endcsname\trSupNorm
\renewcommand{\trSupNorm}{\BookTerm{T0234}{\csname BookOriginal@trSupNorm\endcsname}}
\expandafter\def\csname BookIndex@T0235\endcsname{chashang@差商（difference quotient）}
\expandafter\let\csname BookOriginal@trDifferenceQuotient\endcsname\trDifferenceQuotient
\renewcommand{\trDifferenceQuotient}{\BookTerm{T0235}{\csname BookOriginal@trDifferenceQuotient\endcsname}}
\expandafter\def\csname BookIndex@T0236\endcsname{fangxiangdaoshu@方向导数（directional derivative）}
\expandafter\let\csname BookOriginal@trDirectionalDerivative\endcsname\trDirectionalDerivative
\renewcommand{\trDirectionalDerivative}{\BookTerm{T0236}{\csname BookOriginal@trDirectionalDerivative\endcsname}}
\expandafter\def\csname BookIndex@T0237\endcsname{zhui@锥（cone）}
\expandafter\let\csname BookOriginal@trCone\endcsname\trCone
\renewcommand{\trCone}{\BookTerm{T0237}{\csname BookOriginal@trCone\endcsname}}
\expandafter\def\csname BookIndex@T0238\endcsname{guangyihanshu@广义函数（distribution）}
\expandafter\let\csname BookOriginal@trDistribution\endcsname\trDistribution
\renewcommand{\trDistribution}{\BookTerm{T0238}{\csname BookOriginal@trDistribution\endcsname}}
\expandafter\def\csname BookIndex@T0239\endcsname{kequqidian@可去奇点（removable singularity）}
\expandafter\let\csname BookOriginal@trRemovableSingularity\endcsname\trRemovableSingularity
\renewcommand{\trRemovableSingularity}{\BookTerm{T0239}{\csname BookOriginal@trRemovableSingularity\endcsname}}
\expandafter\def\csname BookIndex@T0240\endcsname{quanshi@全实（totally real）}
\expandafter\let\csname BookOriginal@trTotallyReal\endcsname\trTotallyReal
\renewcommand{\trTotallyReal}{\BookTerm{T0240}{\csname BookOriginal@trTotallyReal\endcsname}}
\expandafter\def\csname BookIndex@T0241\endcsname{banqiu@半球（half ball）}
\expandafter\let\csname BookOriginal@trHalfBall\endcsname\trHalfBall
\renewcommand{\trHalfBall}{\BookTerm{T0241}{\csname BookOriginal@trHalfBall\endcsname}}
\expandafter\def\csname BookIndex@T0242\endcsname{benzhiyoujie@本质有界（essentially bounded）}
\expandafter\let\csname BookOriginal@trEssentiallyBounded\endcsname\trEssentiallyBounded
\renewcommand{\trEssentiallyBounded}{\BookTerm{T0242}{\csname BookOriginal@trEssentiallyBounded\endcsname}}
\expandafter\def\csname BookIndex@T0243\endcsname{tongtai@同态（homomorphism）}
\expandafter\let\csname BookOriginal@trHomomorphism\endcsname\trHomomorphism
\renewcommand{\trHomomorphism}{\BookTerm{T0243}{\csname BookOriginal@trHomomorphism\endcsname}}
\expandafter\def\csname BookIndex@T0244\endcsname{fanbian@反变（contravariant）}
\expandafter\let\csname BookOriginal@trContravariant\endcsname\trContravariant
\renewcommand{\trContravariant}{\BookTerm{T0244}{\csname BookOriginal@trContravariant\endcsname}}
\expandafter\def\csname BookIndex@T0245\endcsname{hanzi@函子（functor）}
\expandafter\let\csname BookOriginal@trFunctor\endcsname\trFunctor
\renewcommand{\trFunctor}{\BookTerm{T0245}{\csname BookOriginal@trFunctor\endcsname}}
\expandafter\def\csname BookIndex@T0246\endcsname{fanchou@范畴（category）}
\expandafter\let\csname BookOriginal@trCategoryTheory\endcsname\trCategoryTheory
\renewcommand{\trCategoryTheory}{\BookTerm{T0246}{\csname BookOriginal@trCategoryTheory\endcsname}}
\expandafter\def\csname BookIndex@T0247\endcsname{taishe@态射（morphism）}
\expandafter\let\csname BookOriginal@trMorphism\endcsname\trMorphism
\renewcommand{\trMorphism}{\BookTerm{T0247}{\csname BookOriginal@trMorphism\endcsname}}
\expandafter\def\csname BookIndex@T0248\endcsname{daishufanchou@代数范畴（algebraic category）}
\expandafter\let\csname BookOriginal@trAlgebraicCategory\endcsname\trAlgebraicCategory
\renewcommand{\trAlgebraicCategory}{\BookTerm{T0248}{\csname BookOriginal@trAlgebraicCategory\endcsname}}
\expandafter\def\csname BookIndex@T0249\endcsname{daishuduixiang@代数对象（algebraic object）}
\expandafter\let\csname BookOriginal@trAlgebraicObject\endcsname\trAlgebraicObject
\renewcommand{\trAlgebraicObject}{\BookTerm{T0249}{\csname BookOriginal@trAlgebraicObject\endcsname}}
\expandafter\def\csname BookIndex@T0250\endcsname{budongdianzhishu@不动点指数（fixed point index）}
\expandafter\let\csname BookOriginal@trFixedPointIndex\endcsname\trFixedPointIndex
\renewcommand{\trFixedPointIndex}{\BookTerm{T0250}{\csname BookOriginal@trFixedPointIndex\endcsname}}
\expandafter\def\csname BookIndex@T0251\endcsname{bianfen@变分（variational）}
\expandafter\let\csname BookOriginal@trVariational\endcsname\trVariational
\renewcommand{\trVariational}{\BookTerm{T0251}{\csname BookOriginal@trVariational\endcsname}}
\expandafter\def\csname BookIndex@T0252\endcsname{bianfenyuanli@变分原理（variational principle）}
\expandafter\let\csname BookOriginal@trVariationalPrinciple\endcsname\trVariationalPrinciple
\renewcommand{\trVariationalPrinciple}{\BookTerm{T0252}{\csname BookOriginal@trVariationalPrinciple\endcsname}}
\expandafter\def\csname BookIndex@T0253\endcsname{bianfenfa@变分法（calculus of variations）}
\expandafter\let\csname BookOriginal@trCalculusOfVariations\endcsname\trCalculusOfVariations
\renewcommand{\trCalculusOfVariations}{\BookTerm{T0253}{\csname BookOriginal@trCalculusOfVariations\endcsname}}
\expandafter\def\csname BookIndex@T0254\endcsname{zhijiefangfa@直接方法（direct method）}
\expandafter\let\csname BookOriginal@trDirectMethod\endcsname\trDirectMethod
\renewcommand{\trDirectMethod}{\BookTerm{T0254}{\csname BookOriginal@trDirectMethod\endcsname}}
\expandafter\def\csname BookIndex@T0255\endcsname{jixiaohuaxulie@极小化序列（minimizing sequence）}
\expandafter\let\csname BookOriginal@trMinimizingSequence\endcsname\trMinimizingSequence
\renewcommand{\trMinimizingSequence}{\BookTerm{T0255}{\csname BookOriginal@trMinimizingSequence\endcsname}}
\expandafter\def\csname BookIndex@T0256\endcsname{jixiaojida@极小极大（minimax）}
\expandafter\let\csname BookOriginal@trMinimax\endcsname\trMinimax
\renewcommand{\trMinimax}{\BookTerm{T0256}{\csname BookOriginal@trMinimax\endcsname}}
\expandafter\def\csname BookIndex@T0257\endcsname{linjie@临界（critical）}
\expandafter\let\csname BookOriginal@trCritical\endcsname\trCritical
\renewcommand{\trCritical}{\BookTerm{T0257}{\csname BookOriginal@trCritical\endcsname}}
\expandafter\def\csname BookIndex@T0258\endcsname{linjiedian@临界点（critical point）}
\expandafter\let\csname BookOriginal@trCriticalPoint\endcsname\trCriticalPoint
\renewcommand{\trCriticalPoint}{\BookTerm{T0258}{\csname BookOriginal@trCriticalPoint\endcsname}}
\expandafter\def\csname BookIndex@T0259\endcsname{linjiezhi@临界值（critical value）}
\expandafter\let\csname BookOriginal@trCriticalValue\endcsname\trCriticalValue
\renewcommand{\trCriticalValue}{\BookTerm{T0259}{\csname BookOriginal@trCriticalValue\endcsname}}
\expandafter\def\csname BookIndex@T0260\endcsname{andian@鞍点（saddle point）}
\expandafter\let\csname BookOriginal@trSaddlePoint\endcsname\trSaddlePoint
\renewcommand{\trSaddlePoint}{\BookTerm{T0260}{\csname BookOriginal@trSaddlePoint\endcsname}}
\expandafter\def\csname BookIndex@T0261\endcsname{tidu@梯度（gradient）}
\expandafter\let\csname BookOriginal@trGradient\endcsname\trGradient
\renewcommand{\trGradient}{\BookTerm{T0261}{\csname BookOriginal@trGradient\endcsname}}
\expandafter\def\csname BookIndex@T0262\endcsname{leitidu@类梯度（gradient-like）}
\expandafter\let\csname BookOriginal@trGradientLike\endcsname\trGradientLike
\renewcommand{\trGradientLike}{\BookTerm{T0262}{\csname BookOriginal@trGradientLike\endcsname}}
\expandafter\def\csname BookIndex@T0263\endcsname{moersililun@莫尔斯理论（Morse theory）}
\expandafter\let\csname BookOriginal@trMorseTheory\endcsname\trMorseTheory
\renewcommand{\trMorseTheory}{\BookTerm{T0263}{\csname BookOriginal@trMorseTheory\endcsname}}
\expandafter\def\csname BookIndex@T0264\endcsname{moersizhishu@莫尔斯指数（Morse index）}
\expandafter\let\csname BookOriginal@trMorseIndex\endcsname\trMorseIndex
\renewcommand{\trMorseIndex}{\BookTerm{T0264}{\csname BookOriginal@trMorseIndex\endcsname}}
\expandafter\def\csname BookIndex@T0265\endcsname{guanghuafanchou@光滑范畴（smooth category）}
\expandafter\let\csname BookOriginal@trSmoothCategory\endcsname\trSmoothCategory
\renewcommand{\trSmoothCategory}{\BookTerm{T0265}{\csname BookOriginal@trSmoothCategory\endcsname}}
\expandafter\def\csname BookIndex@T0266\endcsname{kuige@亏格（genus）}
\expandafter\let\csname BookOriginal@trGenus\endcsname\trGenus
\renewcommand{\trGenus}{\BookTerm{T0266}{\csname BookOriginal@trGenus\endcsname}}
\expandafter\def\csname BookIndex@T0267\endcsname{fanhan@泛函（functional）}
\expandafter\let\csname BookOriginal@trFunctional\endcsname\trFunctional
\renewcommand{\trFunctional}{\BookTerm{T0267}{\csname BookOriginal@trFunctional\endcsname}}
\expandafter\def\csname BookIndex@T0268\endcsname{jianyanhanshu@检验函数（test function）}
\expandafter\let\csname BookOriginal@trTestFunction\endcsname\trTestFunction
\renewcommand{\trTestFunction}{\BookTerm{T0268}{\csname BookOriginal@trTestFunction\endcsname}}
\expandafter\def\csname BookIndex@T0269\endcsname{xiafangyoujie@下方有界（bounded below）}
\expandafter\let\csname BookOriginal@trBoundedBelow\endcsname\trBoundedBelow
\renewcommand{\trBoundedBelow}{\BookTerm{T0269}{\csname BookOriginal@trBoundedBelow\endcsname}}
\expandafter\def\csname BookIndex@T0270\endcsname{shangfangyoujie@上方有界（bounded above）}
\expandafter\let\csname BookOriginal@trBoundedAbove\endcsname\trBoundedAbove
\renewcommand{\trBoundedAbove}{\BookTerm{T0270}{\csname BookOriginal@trBoundedAbove\endcsname}}
\expandafter\def\csname BookIndex@T0271\endcsname{palaisimeier@帕莱-斯梅尔（Palais-Smale）}
\expandafter\let\csname BookOriginal@trPalaisSmale\endcsname\trPalaisSmale
\renewcommand{\trPalaisSmale}{\BookTerm{T0271}{\csname BookOriginal@trPalaisSmale\endcsname}}
\expandafter\def\csname BookIndex@T0272\endcsname{shanluyinli@山路引理（mountain pass lemma）}
\expandafter\let\csname BookOriginal@trMountainPassLemma\endcsname\trMountainPassLemma
\renewcommand{\trMountainPassLemma}{\BookTerm{T0272}{\csname BookOriginal@trMountainPassLemma\endcsname}}
\expandafter\def\csname BookIndex@T0273\endcsname{tongdiao@同调（homology）}
\expandafter\let\csname BookOriginal@trHomology\endcsname\trHomology
\renewcommand{\trHomology}{\BookTerm{T0273}{\csname BookOriginal@trHomology\endcsname}}
\expandafter\def\csname BookIndex@T0274\endcsname{shangtongdiao@上同调（cohomology）}
\expandafter\let\csname BookOriginal@trCohomology\endcsname\trCohomology
\renewcommand{\trCohomology}{\BookTerm{T0274}{\csname BookOriginal@trCohomology\endcsname}}
\expandafter\def\csname BookIndex@T0275\endcsname{tongdiaode@同调的（homologous）}
\expandafter\let\csname BookOriginal@trHomologous\endcsname\trHomologous
\renewcommand{\trHomologous}{\BookTerm{T0275}{\csname BookOriginal@trHomologous\endcsname}}
\expandafter\def\csname BookIndex@T0276\endcsname{qiyitongdiao@奇异同调（singular homology）}
\expandafter\let\csname BookOriginal@trSingularHomology\endcsname\trSingularHomology
\renewcommand{\trSingularHomology}{\BookTerm{T0276}{\csname BookOriginal@trSingularHomology\endcsname}}
\expandafter\def\csname BookIndex@T0277\endcsname{zhenghe@正合（exact）}
\expandafter\let\csname BookOriginal@trExact\endcsname\trExact
\renewcommand{\trExact}{\BookTerm{T0277}{\csname BookOriginal@trExact\endcsname}}
\expandafter\def\csname BookIndex@T0278\endcsname{zhenghexulie@正合序列（exact sequence）}
\expandafter\let\csname BookOriginal@trExactSequence\endcsname\trExactSequence
\renewcommand{\trExactSequence}{\BookTerm{T0278}{\csname BookOriginal@trExactSequence\endcsname}}
\expandafter\def\csname BookIndex@T0279\endcsname{lianfuxing@链复形（chain complex）}
\expandafter\let\csname BookOriginal@trChainComplex\endcsname\trChainComplex
\renewcommand{\trChainComplex}{\BookTerm{T0279}{\csname BookOriginal@trChainComplex\endcsname}}
\expandafter\def\csname BookIndex@T0280\endcsname{bianyuansuanzi@边缘算子（boundary operator）}
\expandafter\let\csname BookOriginal@trBoundaryOperator\endcsname\trBoundaryOperator
\renewcommand{\trBoundaryOperator}{\BookTerm{T0280}{\csname BookOriginal@trBoundaryOperator\endcsname}}
\expandafter\def\csname BookIndex@T0281\endcsname{beiji@杯积（cup product）}
\expandafter\let\csname BookOriginal@trCupProduct\endcsname\trCupProduct
\renewcommand{\trCupProduct}{\BookTerm{T0281}{\csname BookOriginal@trCupProduct\endcsname}}
\expandafter\def\csname BookIndex@T0282\endcsname{beidishu@贝蒂数（Betti number）}
\expandafter\let\csname BookOriginal@trBettiNumber\endcsname\trBettiNumber
\renewcommand{\trBettiNumber}{\BookTerm{T0282}{\csname BookOriginal@trBettiNumber\endcsname}}
\expandafter\def\csname BookIndex@T0283\endcsname{oulashixingshu@欧拉示性数（Euler characteristic）}
\expandafter\let\csname BookOriginal@trEulerCharacteristic\endcsname\trEulerCharacteristic
\renewcommand{\trEulerCharacteristic}{\BookTerm{T0283}{\csname BookOriginal@trEulerCharacteristic\endcsname}}
\expandafter\def\csname BookIndex@T0284\endcsname{jibenqun@基本群（fundamental group）}
\expandafter\let\csname BookOriginal@trFundamentalGroup\endcsname\trFundamentalGroup
\renewcommand{\trFundamentalGroup}{\BookTerm{T0284}{\csname BookOriginal@trFundamentalGroup\endcsname}}
\expandafter\def\csname BookIndex@T0285\endcsname{tonglun@同伦（homotopy）}
\expandafter\let\csname BookOriginal@trHomotopy\endcsname\trHomotopy
\renewcommand{\trHomotopy}{\BookTerm{T0285}{\csname BookOriginal@trHomotopy\endcsname}}
\expandafter\def\csname BookIndex@T0286\endcsname{tonglunde@同伦的（homotopic）}
\expandafter\let\csname BookOriginal@trHomotopic\endcsname\trHomotopic
\renewcommand{\trHomotopic}{\BookTerm{T0286}{\csname BookOriginal@trHomotopic\endcsname}}
\expandafter\def\csname BookIndex@T0287\endcsname{tonglunbubianliang@同伦不变量（homotopy invariant）}
\expandafter\let\csname BookOriginal@trHomotopyInvariant\endcsname\trHomotopyInvariant
\renewcommand{\trHomotopyInvariant}{\BookTerm{T0287}{\csname BookOriginal@trHomotopyInvariant\endcsname}}
\expandafter\def\csname BookIndex@T0288\endcsname{bubianxing@不变性（invariance）}
\expandafter\let\csname BookOriginal@trInvariance\endcsname\trInvariance
\renewcommand{\trInvariance}{\BookTerm{T0288}{\csname BookOriginal@trInvariance\endcsname}}
\expandafter\def\csname BookIndex@T0289\endcsname{tonghen@同痕（isotopy）}
\expandafter\let\csname BookOriginal@trIsotopy\endcsname\trIsotopy
\renewcommand{\trIsotopy}{\BookTerm{T0289}{\csname BookOriginal@trIsotopy\endcsname}}
\expandafter\def\csname BookIndex@T0290\endcsname{tonghende@同痕的（isotopic）}
\expandafter\let\csname BookOriginal@trIsotopic\endcsname\trIsotopic
\renewcommand{\trIsotopic}{\BookTerm{T0290}{\csname BookOriginal@trIsotopic\endcsname}}
\expandafter\def\csname BookIndex@T0291\endcsname{kesuo@可缩（contractible）}
\expandafter\let\csname BookOriginal@trContractible\endcsname\trContractible
\renewcommand{\trContractible}{\BookTerm{T0291}{\csname BookOriginal@trContractible\endcsname}}
\expandafter\def\csname BookIndex@T0292\endcsname{xingbian@形变（deformation）}
\expandafter\let\csname BookOriginal@trDeformation\endcsname\trDeformation
\renewcommand{\trDeformation}{\BookTerm{T0292}{\csname BookOriginal@trDeformation\endcsname}}
\expandafter\def\csname BookIndex@T0293\endcsname{shang@商（quotient）}
\expandafter\let\csname BookOriginal@trQuotient\endcsname\trQuotient
\renewcommand{\trQuotient}{\BookTerm{T0293}{\csname BookOriginal@trQuotient\endcsname}}
\expandafter\def\csname BookIndex@T0294\endcsname{du@度（degree）}
\expandafter\let\csname BookOriginal@trDegree\endcsname\trDegree
\renewcommand{\trDegree}{\BookTerm{T0294}{\csname BookOriginal@trDegree\endcsname}}
\expandafter\def\csname BookIndex@T0295\endcsname{dandiaoxing@单调性（monotonicity）}
\expandafter\let\csname BookOriginal@trMonotonicity\endcsname\trMonotonicity
\renewcommand{\trMonotonicity}{\BookTerm{T0295}{\csname BookOriginal@trMonotonicity\endcsname}}
\expandafter\def\csname BookIndex@T0296\endcsname{qiechu@切除（excision）}
\expandafter\let\csname BookOriginal@trExcision\endcsname\trExcision
\renewcommand{\trExcision}{\BookTerm{T0296}{\csname BookOriginal@trExcision\endcsname}}
\expandafter\def\csname BookIndex@T0297\endcsname{kejiaxing@可加性（additivity）}
\expandafter\let\csname BookOriginal@trAdditivity\endcsname\trAdditivity
\renewcommand{\trAdditivity}{\BookTerm{T0297}{\csname BookOriginal@trAdditivity\endcsname}}
\expandafter\def\csname BookIndex@T0298\endcsname{guifanxing@规范性（normalization）}
\expandafter\let\csname BookOriginal@trNormalization\endcsname\trNormalization
\renewcommand{\trNormalization}{\BookTerm{T0298}{\csname BookOriginal@trNormalization\endcsname}}
\expandafter\def\csname BookIndex@T0299\endcsname{ziranxing@自然性（naturality）}
\expandafter\let\csname BookOriginal@trNaturality\endcsname\trNaturality
\renewcommand{\trNaturality}{\BookTerm{T0299}{\csname BookOriginal@trNaturality\endcsname}}
\expandafter\def\csname BookIndex@T0300\endcsname{gongxingxing@共形性（conformality）}
\expandafter\let\csname BookOriginal@trConformality\endcsname\trConformality
\renewcommand{\trConformality}{\BookTerm{T0300}{\csname BookOriginal@trConformality\endcsname}}
\expandafter\def\csname BookIndex@T0301\endcsname{jiexingzhi@解性质（solution property）}
\expandafter\let\csname BookOriginal@trSolutionProperty\endcsname\trSolutionProperty
\renewcommand{\trSolutionProperty}{\BookTerm{T0301}{\csname BookOriginal@trSolutionProperty\endcsname}}
\expandafter\def\csname BookIndex@T0302\endcsname{liuxing@流形（manifold）}
\expandafter\let\csname BookOriginal@trManifold\endcsname\trManifold
\renewcommand{\trManifold}{\BookTerm{T0302}{\csname BookOriginal@trManifold\endcsname}}
\expandafter\def\csname BookIndex@T0303\endcsname{ziliuxing@子流形（submanifold）}
\expandafter\let\csname BookOriginal@trSubmanifold\endcsname\trSubmanifold
\renewcommand{\trSubmanifold}{\BookTerm{T0303}{\csname BookOriginal@trSubmanifold\endcsname}}
\expandafter\def\csname BookIndex@T0304\endcsname{chaoqumian@超曲面（hypersurface）}
\expandafter\let\csname BookOriginal@trHypersurface\endcsname\trHypersurface
\renewcommand{\trHypersurface}{\BookTerm{T0304}{\csname BookOriginal@trHypersurface\endcsname}}
\expandafter\def\csname BookIndex@T0305\endcsname{chaopingmian@超平面（hyperplane）}
\expandafter\let\csname BookOriginal@trHyperplane\endcsname\trHyperplane
\renewcommand{\trHyperplane}{\BookTerm{T0305}{\csname BookOriginal@trHyperplane\endcsname}}
\expandafter\def\csname BookIndex@T0306\endcsname{qiumian@球面（sphere）}
\expandafter\let\csname BookOriginal@trSphere\endcsname\trSphere
\renewcommand{\trSphere}{\BookTerm{T0306}{\csname BookOriginal@trSphere\endcsname}}
\expandafter\def\csname BookIndex@T0307\endcsname{qumian@曲面（surface）}
\expandafter\let\csname BookOriginal@trSurface\endcsname\trSurface
\renewcommand{\trSurface}{\BookTerm{T0307}{\csname BookOriginal@trSurface\endcsname}}
\expandafter\def\csname BookIndex@T0308\endcsname{xiangliangchang@向量场（vector field）}
\expandafter\let\csname BookOriginal@trVectorField\endcsname\trVectorField
\renewcommand{\trVectorField}{\BookTerm{T0308}{\csname BookOriginal@trVectorField\endcsname}}
\expandafter\def\csname BookIndex@T0309\endcsname{weifentongpei@微分同胚（diffeomorphism）}
\expandafter\let\csname BookOriginal@trDiffeomorphism\endcsname\trDiffeomorphism
\renewcommand{\trDiffeomorphism}{\BookTerm{T0309}{\csname BookOriginal@trDiffeomorphism\endcsname}}
\expandafter\def\csname BookIndex@T0310\endcsname{weifentongpeide@微分同胚的（diffeomorphic）}
\expandafter\let\csname BookOriginal@trDiffeomorphic\endcsname\trDiffeomorphic
\renewcommand{\trDiffeomorphic}{\BookTerm{T0310}{\csname BookOriginal@trDiffeomorphic\endcsname}}
\expandafter\def\csname BookIndex@T0311\endcsname{tongpei@同胚（homeomorphism）}
\expandafter\let\csname BookOriginal@trHomeomorphism\endcsname\trHomeomorphism
\renewcommand{\trHomeomorphism}{\BookTerm{T0311}{\csname BookOriginal@trHomeomorphism\endcsname}}
\expandafter\def\csname BookIndex@T0312\endcsname{qianru@嵌入（embedding）}
\expandafter\let\csname BookOriginal@trEmbedding\endcsname\trEmbedding
\renewcommand{\trEmbedding}{\BookTerm{T0312}{\csname BookOriginal@trEmbedding\endcsname}}
\expandafter\def\csname BookIndex@T0313\endcsname{jinru@浸入（immersion）}
\expandafter\let\csname BookOriginal@trImmersion\endcsname\trImmersion
\renewcommand{\trImmersion}{\BookTerm{T0313}{\csname BookOriginal@trImmersion\endcsname}}
\expandafter\def\csname BookIndex@T0314\endcsname{weifenxingshi@微分形式（differential form）}
\expandafter\let\csname BookOriginal@trDifferentialForm\endcsname\trDifferentialForm
\renewcommand{\trDifferentialForm}{\BookTerm{T0314}{\csname BookOriginal@trDifferentialForm\endcsname}}
\expandafter\def\csname BookIndex@T0315\endcsname{1xingshi@1-形式（one-form）}
\expandafter\let\csname BookOriginal@trOneForm\endcsname\trOneForm
\renewcommand{\trOneForm}{\BookTerm{T0315}{\csname BookOriginal@trOneForm\endcsname}}
\expandafter\def\csname BookIndex@T0316\endcsname{2xingshi@2-形式（two-form）}
\expandafter\let\csname BookOriginal@trTwoForm\endcsname\trTwoForm
\renewcommand{\trTwoForm}{\BookTerm{T0316}{\csname BookOriginal@trTwoForm\endcsname}}
\expandafter\def\csname BookIndex@T0317\endcsname{waiweifen@外微分（exterior derivative）}
\expandafter\let\csname BookOriginal@trExteriorDerivative\endcsname\trExteriorDerivative
\renewcommand{\trExteriorDerivative}{\BookTerm{T0317}{\csname BookOriginal@trExteriorDerivative\endcsname}}
\expandafter\def\csname BookIndex@T0318\endcsname{jiadangyansuan@嘉当演算（Cartan's calculus）}
\expandafter\let\csname BookOriginal@trCartanCalculus\endcsname\trCartanCalculus
\renewcommand{\trCartanCalculus}{\BookTerm{T0318}{\csname BookOriginal@trCartanCalculus\endcsname}}
\expandafter\def\csname BookIndex@T0319\endcsname{lahui@拉回（pullback）}
\expandafter\let\csname BookOriginal@trPullback\endcsname\trPullback
\renewcommand{\trPullback}{\BookTerm{T0319}{\csname BookOriginal@trPullback\endcsname}}
\expandafter\def\csname BookIndex@T0320\endcsname{tuiqian@推前（pushforward）}
\expandafter\let\csname BookOriginal@trPushforward\endcsname\trPushforward
\renewcommand{\trPushforward}{\BookTerm{T0320}{\csname BookOriginal@trPushforward\endcsname}}
\expandafter\def\csname BookIndex@T0321\endcsname{feituihua@非退化（nondegenerate）}
\expandafter\let\csname BookOriginal@trNondegenerate\endcsname\trNondegenerate
\renewcommand{\trNondegenerate}{\BookTerm{T0321}{\csname BookOriginal@trNondegenerate\endcsname}}
\expandafter\def\csname BookIndex@T0322\endcsname{tuihua@退化（degenerate）}
\expandafter\let\csname BookOriginal@trDegenerate\endcsname\trDegenerate
\renewcommand{\trDegenerate}{\BookTerm{T0322}{\csname BookOriginal@trDegenerate\endcsname}}
\expandafter\def\csname BookIndex@T0323\endcsname{fanduicheng@反对称（skew symmetric）}
\expandafter\let\csname BookOriginal@trSkewSymmetric\endcsname\trSkewSymmetric
\renewcommand{\trSkewSymmetric}{\BookTerm{T0323}{\csname BookOriginal@trSkewSymmetric\endcsname}}
\expandafter\def\csname BookIndex@T0324\endcsname{duicheng@对称（symmetric）}
\expandafter\let\csname BookOriginal@trSymmetric\endcsname\trSymmetric
\renewcommand{\trSymmetric}{\BookTerm{T0324}{\csname BookOriginal@trSymmetric\endcsname}}
\expandafter\def\csname BookIndex@T0325\endcsname{duijiao@对角（diagonal）}
\expandafter\let\csname BookOriginal@trDiagonal\endcsname\trDiagonal
\renewcommand{\trDiagonal}{\BookTerm{T0325}{\csname BookOriginal@trDiagonal\endcsname}}
\expandafter\def\csname BookIndex@T0326\endcsname{shuangxianxing@双线性（bilinear）}
\expandafter\let\csname BookOriginal@trBilinear\endcsname\trBilinear
\renewcommand{\trBilinear}{\BookTerm{T0326}{\csname BookOriginal@trBilinear\endcsname}}
\expandafter\def\csname BookIndex@T0327\endcsname{zhengjiao@正交（orthogonal）}
\expandafter\let\csname BookOriginal@trOrthogonal\endcsname\trOrthogonal
\renewcommand{\trOrthogonal}{\BookTerm{T0327}{\csname BookOriginal@trOrthogonal\endcsname}}
\expandafter\def\csname BookIndex@T0328\endcsname{zhengjiaoxing@正交性（orthogonality）}
\expandafter\let\csname BookOriginal@trOrthogonality\endcsname\trOrthogonality
\renewcommand{\trOrthogonality}{\BookTerm{T0328}{\csname BookOriginal@trOrthogonality\endcsname}}
\expandafter\def\csname BookIndex@T0329\endcsname{zhengjiaobu@正交补（orthogonal complement）}
\expandafter\let\csname BookOriginal@trOrthogonalComplement\endcsname\trOrthogonalComplement
\renewcommand{\trOrthogonalComplement}{\BookTerm{T0329}{\csname BookOriginal@trOrthogonalComplement\endcsname}}
\expandafter\def\csname BookIndex@T0330\endcsname{zhengjiaohe@正交和（orthogonal sum）}
\expandafter\let\csname BookOriginal@trOrthogonalSum\endcsname\trOrthogonalSum
\renewcommand{\trOrthogonalSum}{\BookTerm{T0330}{\csname BookOriginal@trOrthogonalSum\endcsname}}
\expandafter\def\csname BookIndex@T0331\endcsname{touying@投影（projection）}
\expandafter\let\csname BookOriginal@trProjection\endcsname\trProjection
\renewcommand{\trProjection}{\BookTerm{T0331}{\csname BookOriginal@trProjection\endcsname}}
\expandafter\def\csname BookIndex@T0332\endcsname{oujilide@欧几里得（Euclidean）}
\expandafter\let\csname BookOriginal@trEuclidean\endcsname\trEuclidean
\renewcommand{\trEuclidean}{\BookTerm{T0332}{\csname BookOriginal@trEuclidean\endcsname}}
\expandafter\def\csname BookIndex@T0333\endcsname{neiji@内积（inner product）}
\expandafter\let\csname BookOriginal@trInnerProduct\endcsname\trInnerProduct
\renewcommand{\trInnerProduct}{\BookTerm{T0333}{\csname BookOriginal@trInnerProduct\endcsname}}
\expandafter\def\csname BookIndex@T0334\endcsname{duiou@对偶（dual）}
\expandafter\let\csname BookOriginal@trDual\endcsname\trDual
\renewcommand{\trDual}{\BookTerm{T0334}{\csname BookOriginal@trDual\endcsname}}
\expandafter\def\csname BookIndex@T0335\endcsname{zikongjian@子空间（linear subspace）}
\expandafter\let\csname BookOriginal@trSubspace\endcsname\trSubspace
\renewcommand{\trSubspace}{\BookTerm{T0335}{\csname BookOriginal@trSubspace\endcsname}}
\expandafter\def\csname BookIndex@T0336\endcsname{chengjikongjian@乘积空间（product space）}
\expandafter\let\csname BookOriginal@trProductSpace\endcsname\trProductSpace
\renewcommand{\trProductSpace}{\BookTerm{T0336}{\csname BookOriginal@trProductSpace\endcsname}}
\expandafter\def\csname BookIndex@T0337\endcsname{danshe@单射（injective）}
\expandafter\let\csname BookOriginal@trInjective\endcsname\trInjective
\renewcommand{\trInjective}{\BookTerm{T0337}{\csname BookOriginal@trInjective\endcsname}}
\expandafter\def\csname BookIndex@T0338\endcsname{manshe@满射（surjective）}
\expandafter\let\csname BookOriginal@trSurjective\endcsname\trSurjective
\renewcommand{\trSurjective}{\BookTerm{T0338}{\csname BookOriginal@trSurjective\endcsname}}
\expandafter\def\csname BookIndex@T0339\endcsname{feiqiyi@非奇异（nonsingular）}
\expandafter\let\csname BookOriginal@trNonsingular\endcsname\trNonsingular
\renewcommand{\trNonsingular}{\BookTerm{T0339}{\csname BookOriginal@trNonsingular\endcsname}}
\expandafter\def\csname BookIndex@T0340\endcsname{cong@丛（bundle）}
\expandafter\let\csname BookOriginal@trBundle\endcsname\trBundle
\renewcommand{\trBundle}{\BookTerm{T0340}{\csname BookOriginal@trBundle\endcsname}}
\expandafter\def\csname BookIndex@T0341\endcsname{xiangliangcong@向量丛（vector bundle）}
\expandafter\let\csname BookOriginal@trVectorBundle\endcsname\trVectorBundle
\renewcommand{\trVectorBundle}{\BookTerm{T0341}{\csname BookOriginal@trVectorBundle\endcsname}}
\expandafter\def\csname BookIndex@T0342\endcsname{xiancong@线丛（line bundle）}
\expandafter\let\csname BookOriginal@trLineBundle\endcsname\trLineBundle
\renewcommand{\trLineBundle}{\BookTerm{T0342}{\csname BookOriginal@trLineBundle\endcsname}}
\expandafter\def\csname BookIndex@T0343\endcsname{jiemian@截面（section）}
\expandafter\let\csname BookOriginal@trSection\endcsname\trSection
\renewcommand{\trSection}{\BookTerm{T0343}{\csname BookOriginal@trSection\endcsname}}
\expandafter\def\csname BookIndex@T0344\endcsname{xianwei@纤维（fiber）}
\expandafter\let\csname BookOriginal@trFiber\endcsname\trFiber
\renewcommand{\trFiber}{\BookTerm{T0344}{\csname BookOriginal@trFiber\endcsname}}
\expandafter\def\csname BookIndex@T0345\endcsname{qiecong@切丛（tangent bundle）}
\expandafter\let\csname BookOriginal@trTangentBundle\endcsname\trTangentBundle
\renewcommand{\trTangentBundle}{\BookTerm{T0345}{\csname BookOriginal@trTangentBundle\endcsname}}
\expandafter\def\csname BookIndex@T0346\endcsname{yuqiecong@余切丛（cotangent bundle）}
\expandafter\let\csname BookOriginal@trCotangentBundle\endcsname\trCotangentBundle
\renewcommand{\trCotangentBundle}{\BookTerm{T0346}{\csname BookOriginal@trCotangentBundle\endcsname}}
\expandafter\def\csname BookIndex@T0347\endcsname{facong@法丛（normal bundle）}
\expandafter\let\csname BookOriginal@trNormalBundle\endcsname\trNormalBundle
\renewcommand{\trNormalBundle}{\BookTerm{T0347}{\csname BookOriginal@trNormalBundle\endcsname}}
\expandafter\def\csname BookIndex@T0348\endcsname{qiekongjian@切空间（tangent space）}
\expandafter\let\csname BookOriginal@trTangentSpace\endcsname\trTangentSpace
\renewcommand{\trTangentSpace}{\BookTerm{T0348}{\csname BookOriginal@trTangentSpace\endcsname}}
\expandafter\def\csname BookIndex@T0349\endcsname{yuqiekongjian@余切空间（cotangent space）}
\expandafter\let\csname BookOriginal@trCotangentSpace\endcsname\trCotangentSpace
\renewcommand{\trCotangentSpace}{\BookTerm{T0349}{\csname BookOriginal@trCotangentSpace\endcsname}}
\expandafter\def\csname BookIndex@T0350\endcsname{xiangkongjian@相空间（phase space）}
\expandafter\let\csname BookOriginal@trPhaseSpace\endcsname\trPhaseSpace
\renewcommand{\trPhaseSpace}{\BookTerm{T0350}{\csname BookOriginal@trPhaseSpace\endcsname}}
\expandafter\def\csname BookIndex@T0351\endcsname{weixingkongjian@位形空间（configuration space）}
\expandafter\let\csname BookOriginal@trConfigurationSpace\endcsname\trConfigurationSpace
\renewcommand{\trConfigurationSpace}{\BookTerm{T0351}{\csname BookOriginal@trConfigurationSpace\endcsname}}
\expandafter\def\csname BookIndex@T0352\endcsname{bilu@闭路（loop）}
\expandafter\let\csname BookOriginal@trLoop\endcsname\trLoop
\renewcommand{\trLoop}{\BookTerm{T0352}{\csname BookOriginal@trLoop\endcsname}}
\expandafter\def\csname BookIndex@T0353\endcsname{bilukongjian@闭路空间（loop space）}
\expandafter\let\csname BookOriginal@trLoopSpace\endcsname\trLoopSpace
\renewcommand{\trLoopSpace}{\BookTerm{T0353}{\csname BookOriginal@trLoopSpace\endcsname}}
\expandafter\def\csname BookIndex@T0354\endcsname{cedexian@测地线（geodesic）}
\expandafter\let\csname BookOriginal@trGeodesic\endcsname\trGeodesic
\renewcommand{\trGeodesic}{\BookTerm{T0354}{\csname BookOriginal@trGeodesic\endcsname}}
\expandafter\def\csname BookIndex@T0355\endcsname{lujing@路径（path）}
\expandafter\let\csname BookOriginal@trPath\endcsname\trPath
\renewcommand{\trPath}{\BookTerm{T0355}{\csname BookOriginal@trPath\endcsname}}
\expandafter\def\csname BookIndex@T0356\endcsname{daoluliantong@道路连通（path-connected）}
\expandafter\let\csname BookOriginal@trPathConnected\endcsname\trPathConnected
\renewcommand{\trPathConnected}{\BookTerm{T0356}{\csname BookOriginal@trPathConnected\endcsname}}
\expandafter\def\csname BookIndex@T0357\endcsname{canshuhuachaoqumianzu@参数化超曲面族（parameterized family of hypersurfaces）}
\expandafter\let\csname BookOriginal@trParameterizedHypersurfaceFamily\endcsname\trParameterizedHypersurfaceFamily
\renewcommand{\trParameterizedHypersurfaceFamily}{\BookTerm{T0357}{\csname BookOriginal@trParameterizedHypersurfaceFamily\endcsname}}
\expandafter\def\csname BookIndex@T0358\endcsname{dengju@等距（isometry）}
\expandafter\let\csname BookOriginal@trIsometry\endcsname\trIsometry
\renewcommand{\trIsometry}{\BookTerm{T0358}{\csname BookOriginal@trIsometry\endcsname}}
\expandafter\def\csname BookIndex@T0359\endcsname{bi@闭（closed）}
\expandafter\let\csname BookOriginal@trClosed\endcsname\trClosed
\renewcommand{\trClosed}{\BookTerm{T0359}{\csname BookOriginal@trClosed\endcsname}}
\expandafter\def\csname BookIndex@T0360\endcsname{quxian@曲线（curve）}
\expandafter\let\csname BookOriginal@trCurve\endcsname\trCurve
\renewcommand{\trCurve}{\BookTerm{T0360}{\csname BookOriginal@trCurve\endcsname}}
\expandafter\def\csname BookIndex@T0361\endcsname{guanghua@光滑（smooth）}
\expandafter\let\csname BookOriginal@trSmooth\endcsname\trSmooth
\renewcommand{\trSmooth}{\BookTerm{T0361}{\csname BookOriginal@trSmooth\endcsname}}
\expandafter\def\csname BookIndex@T0362\endcsname{daoshu@导数（derivative）}
\expandafter\let\csname BookOriginal@trDerivative\endcsname\trDerivative
\renewcommand{\trDerivative}{\BookTerm{T0362}{\csname BookOriginal@trDerivative\endcsname}}
\expandafter\def\csname BookIndex@T0363\endcsname{duliang@度量（metric）}
\expandafter\let\csname BookOriginal@trMetric\endcsname\trMetric
\renewcommand{\trMetric}{\BookTerm{T0363}{\csname BookOriginal@trMetric\endcsname}}
\expandafter\def\csname BookIndex@T0364\endcsname{fubinishitudiduliang@富比尼-施图迪度量（Fubini-Study metric）}
\expandafter\let\csname BookOriginal@trFubiniStudyMetric\endcsname\trFubiniStudyMetric
\renewcommand{\trFubiniStudyMetric}{\BookTerm{T0364}{\csname BookOriginal@trFubiniStudyMetric\endcsname}}
\expandafter\def\csname BookIndex@T0365\endcsname{shuangbubianduliang@双不变度量（bi-invariant metric）}
\expandafter\let\csname BookOriginal@trBiInvariantMetric\endcsname\trBiInvariantMetric
\renewcommand{\trBiInvariantMetric}{\BookTerm{T0365}{\csname BookOriginal@trBiInvariantMetric\endcsname}}
\expandafter\def\csname BookIndex@T0366\endcsname{juli@距离（distance）}
\expandafter\let\csname BookOriginal@trDistance\endcsname\trDistance
\renewcommand{\trDistance}{\BookTerm{T0366}{\csname BookOriginal@trDistance\endcsname}}
\expandafter\def\csname BookIndex@T0367\endcsname{tapu@拓扑（topology）}
\expandafter\let\csname BookOriginal@trTopology\endcsname\trTopology
\renewcommand{\trTopology}{\BookTerm{T0367}{\csname BookOriginal@trTopology\endcsname}}
\expandafter\def\csname BookIndex@T0368\endcsname{jinzhi@紧致（compact）}
\expandafter\let\csname BookOriginal@trCompact\endcsname\trCompact
\renewcommand{\trCompact}{\BookTerm{T0368}{\csname BookOriginal@trCompact\endcsname}}
\expandafter\def\csname BookIndex@T0369\endcsname{jinxing@紧性（compactness）}
\expandafter\let\csname BookOriginal@trCompactness\endcsname\trCompactness
\renewcommand{\trCompactness}{\BookTerm{T0369}{\csname BookOriginal@trCompactness\endcsname}}
\expandafter\def\csname BookIndex@T0370\endcsname{jinzhicheng@紧支撑（compact support）}
\expandafter\let\csname BookOriginal@trCompactSupport\endcsname\trCompactSupport
\renewcommand{\trCompactSupport}{\BookTerm{T0370}{\csname BookOriginal@trCompactSupport\endcsname}}
\expandafter\def\csname BookIndex@T0371\endcsname{zhicheng@支撑（support）}
\expandafter\let\csname BookOriginal@trSupport\endcsname\trSupport
\renewcommand{\trSupport}{\BookTerm{T0371}{\csname BookOriginal@trSupport\endcsname}}
\expandafter\def\csname BookIndex@T0372\endcsname{liantong@连通（connected）}
\expandafter\let\csname BookOriginal@trConnected\endcsname\trConnected
\renewcommand{\trConnected}{\BookTerm{T0372}{\csname BookOriginal@trConnected\endcsname}}
\expandafter\def\csname BookIndex@T0373\endcsname{liantongfenzhi@连通分支（connected component）}
\expandafter\let\csname BookOriginal@trConnectedComponent\endcsname\trConnectedComponent
\renewcommand{\trConnectedComponent}{\BookTerm{T0373}{\csname BookOriginal@trConnectedComponent\endcsname}}
\expandafter\def\csname BookIndex@T0374\endcsname{fenli@分离（separate）}
\expandafter\let\csname BookOriginal@trSeparate\endcsname\trSeparate
\renewcommand{\trSeparate}{\BookTerm{T0374}{\csname BookOriginal@trSeparate\endcsname}}
\expandafter\def\csname BookIndex@T0375\endcsname{danliantong@单连通（simply connected）}
\expandafter\let\csname BookOriginal@trSimplyConnected\endcsname\trSimplyConnected
\renewcommand{\trSimplyConnected}{\BookTerm{T0375}{\csname BookOriginal@trSimplyConnected\endcsname}}
\expandafter\def\csname BookIndex@T0376\endcsname{dingxiang@定向（oriented）}
\expandafter\let\csname BookOriginal@trOriented\endcsname\trOriented
\renewcommand{\trOriented}{\BookTerm{T0376}{\csname BookOriginal@trOriented\endcsname}}
\expandafter\def\csname BookIndex@T0377\endcsname{kedingxiang@可定向（orientable）}
\expandafter\let\csname BookOriginal@trOrientable\endcsname\trOrientable
\renewcommand{\trOrientable}{\BookTerm{T0377}{\csname BookOriginal@trOrientable\endcsname}}
\expandafter\def\csname BookIndex@T0378\endcsname{dingxiang@定向（orientation）}
\expandafter\let\csname BookOriginal@trOrientation\endcsname\trOrientation
\renewcommand{\trOrientation}{\BookTerm{T0378}{\csname BookOriginal@trOrientation\endcsname}}
\expandafter\def\csname BookIndex@T0379\endcsname{tu@凸（convex）}
\expandafter\let\csname BookOriginal@trConvex\endcsname\trConvex
\renewcommand{\trConvex}{\BookTerm{T0379}{\csname BookOriginal@trConvex\endcsname}}
\expandafter\def\csname BookIndex@T0380\endcsname{yangetu@严格凸（strictly convex）}
\expandafter\let\csname BookOriginal@trStrictlyConvex\endcsname\trStrictlyConvex
\renewcommand{\trStrictlyConvex}{\BookTerm{T0380}{\csname BookOriginal@trStrictlyConvex\endcsname}}
\expandafter\def\csname BookIndex@T0381\endcsname{tuxing@凸性（convexity）}
\expandafter\let\csname BookOriginal@trConvexity\endcsname\trConvexity
\renewcommand{\trConvexity}{\BookTerm{T0381}{\csname BookOriginal@trConvexity\endcsname}}
\expandafter\def\csname BookIndex@T0382\endcsname{xingxing@星形（star shaped）}
\expandafter\let\csname BookOriginal@trStarShaped\endcsname\trStarShaped
\renewcommand{\trStarShaped}{\BookTerm{T0382}{\csname BookOriginal@trStarShaped\endcsname}}
\expandafter\def\csname BookIndex@T0383\endcsname{quyu@区域（domain）}
\expandafter\let\csname BookOriginal@trDomain\endcsname\trDomain
\renewcommand{\trDomain}{\BookTerm{T0383}{\csname BookOriginal@trDomain\endcsname}}
\expandafter\def\csname BookIndex@T0384\endcsname{bianjie@边界（boundary）}
\expandafter\let\csname BookOriginal@trBoundary\endcsname\trBoundary
\renewcommand{\trBoundary}{\BookTerm{T0384}{\csname BookOriginal@trBoundary\endcsname}}
\expandafter\def\csname BookIndex@T0385\endcsname{bianjietiaojian@边界条件（boundary condition）}
\expandafter\let\csname BookOriginal@trBoundaryCondition\endcsname\trBoundaryCondition
\renewcommand{\trBoundaryCondition}{\BookTerm{T0385}{\csname BookOriginal@trBoundaryCondition\endcsname}}
\expandafter\def\csname BookIndex@T0386\endcsname{bianzhiwenti@边值问题（boundary value problem）}
\expandafter\let\csname BookOriginal@trBoundaryValueProblem\endcsname\trBoundaryValueProblem
\renewcommand{\trBoundaryValueProblem}{\BookTerm{T0386}{\csname BookOriginal@trBoundaryValueProblem\endcsname}}
\expandafter\def\csname BookIndex@T0387\endcsname{neibu@内部（interior）}
\expandafter\let\csname BookOriginal@trInterior\endcsname\trInterior
\renewcommand{\trInterior}{\BookTerm{T0387}{\csname BookOriginal@trInterior\endcsname}}
\expandafter\def\csname BookIndex@T0388\endcsname{bibao@闭包（closure）}
\expandafter\let\csname BookOriginal@trClosure\endcsname\trClosure
\renewcommand{\trClosure}{\BookTerm{T0388}{\csname BookOriginal@trClosure\endcsname}}
\expandafter\def\csname BookIndex@T0389\endcsname{linyu@邻域（neighborhood）}
\expandafter\let\csname BookOriginal@trNeighborhood\endcsname\trNeighborhood
\renewcommand{\trNeighborhood}{\BookTerm{T0389}{\csname BookOriginal@trNeighborhood\endcsname}}
\expandafter\def\csname BookIndex@T0390\endcsname{quanchun@全纯（holomorphic）}
\expandafter\let\csname BookOriginal@trHolomorphic\endcsname\trHolomorphic
\renewcommand{\trHolomorphic}{\BookTerm{T0390}{\csname BookOriginal@trHolomorphic\endcsname}}
\expandafter\def\csname BookIndex@T0391\endcsname{weiquanchun@伪全纯（pseudoholomorphic）}
\expandafter\let\csname BookOriginal@trPseudoholomorphic\endcsname\trPseudoholomorphic
\renewcommand{\trPseudoholomorphic}{\BookTerm{T0391}{\csname BookOriginal@trPseudoholomorphic\endcsname}}
\expandafter\def\csname BookIndex@T0392\endcsname{gongxing@共形（conformal）}
\expandafter\let\csname BookOriginal@trConformal\endcsname\trConformal
\renewcommand{\trConformal}{\BookTerm{T0392}{\csname BookOriginal@trConformal\endcsname}}
\expandafter\def\csname BookIndex@T0393\endcsname{diaohe@调和（harmonic）}
\expandafter\let\csname BookOriginal@trHarmonic\endcsname\trHarmonic
\renewcommand{\trHarmonic}{\BookTerm{T0393}{\csname BookOriginal@trHarmonic\endcsname}}
\expandafter\def\csname BookIndex@T0394\endcsname{jiexi@解析（analytic）}
\expandafter\let\csname BookOriginal@trAnalytic\endcsname\trAnalytic
\renewcommand{\trAnalytic}{\BookTerm{T0394}{\csname BookOriginal@trAnalytic\endcsname}}
\expandafter\def\csname BookIndex@T0395\endcsname{biaozhun@标准（standard）}
\expandafter\let\csname BookOriginal@trStandard\endcsname\trStandard
\renewcommand{\trStandard}{\BookTerm{T0395}{\csname BookOriginal@trStandard\endcsname}}
\expandafter\def\csname BookIndex@T0396\endcsname{guaiyi@怪异（exotic）}
\expandafter\let\csname BookOriginal@trExotic\endcsname\trExotic
\renewcommand{\trExotic}{\BookTerm{T0396}{\csname BookOriginal@trExotic\endcsname}}
\expandafter\def\csname BookIndex@T0397\endcsname{suanzi@算子（operator）}
\expandafter\let\csname BookOriginal@trOperator\endcsname\trOperator
\renewcommand{\trOperator}{\BookTerm{T0397}{\csname BookOriginal@trOperator\endcsname}}
\expandafter\def\csname BookIndex@T0398\endcsname{xianxing@线性（linear）}
\expandafter\let\csname BookOriginal@trLinear\endcsname\trLinear
\renewcommand{\trLinear}{\BookTerm{T0398}{\csname BookOriginal@trLinear\endcsname}}
\expandafter\def\csname BookIndex@T0399\endcsname{xianxingsuanzi@线性算子（linear operator）}
\expandafter\let\csname BookOriginal@trLinearOperator\endcsname\trLinearOperator
\renewcommand{\trLinearOperator}{\BookTerm{T0399}{\csname BookOriginal@trLinearOperator\endcsname}}
\expandafter\def\csname BookIndex@T0400\endcsname{xianxingwuguan@线性无关（linearly independent）}
\expandafter\let\csname BookOriginal@trLinearlyIndependent\endcsname\trLinearlyIndependent
\renewcommand{\trLinearlyIndependent}{\BookTerm{T0400}{\csname BookOriginal@trLinearlyIndependent\endcsname}}
\expandafter\def\csname BookIndex@T0401\endcsname{youliwuguan@有理无关（rationally independent）}
\expandafter\let\csname BookOriginal@trRationallyIndependent\endcsname\trRationallyIndependent
\renewcommand{\trRationallyIndependent}{\BookTerm{T0401}{\csname BookOriginal@trRationallyIndependent\endcsname}}
\expandafter\def\csname BookIndex@T0402\endcsname{xianxingbianhuan@线性变换（linear transformation）}
\expandafter\let\csname BookOriginal@trLinearTransformation\endcsname\trLinearTransformation
\renewcommand{\trLinearTransformation}{\BookTerm{T0402}{\csname BookOriginal@trLinearTransformation\endcsname}}
\expandafter\def\csname BookIndex@T0403\endcsname{lianshifaze@链式法则（chain rule）}
\expandafter\let\csname BookOriginal@trChainRule\endcsname\trChainRule
\renewcommand{\trChainRule}{\BookTerm{T0403}{\csname BookOriginal@trChainRule\endcsname}}
\expandafter\def\csname BookIndex@T0404\endcsname{xianxinghua@线性化（linearization）}
\expandafter\let\csname BookOriginal@trLinearization\endcsname\trLinearization
\renewcommand{\trLinearization}{\BookTerm{T0404}{\csname BookOriginal@trLinearization\endcsname}}
\expandafter\def\csname BookIndex@T0405\endcsname{feixianxing@非线性（nonlinear）}
\expandafter\let\csname BookOriginal@trNonlinear\endcsname\trNonlinear
\renewcommand{\trNonlinear}{\BookTerm{T0405}{\csname BookOriginal@trNonlinear\endcsname}}
\expandafter\def\csname BookIndex@T0406\endcsname{qici@齐次（homogeneous）}
\expandafter\let\csname BookOriginal@trHomogeneous\endcsname\trHomogeneous
\renewcommand{\trHomogeneous}{\BookTerm{T0406}{\csname BookOriginal@trHomogeneous\endcsname}}
\expandafter\def\csname BookIndex@T0407\endcsname{qicixing@齐次性（homogeneity）}
\expandafter\let\csname BookOriginal@trHomogeneity\endcsname\trHomogeneity
\renewcommand{\trHomogeneity}{\BookTerm{T0407}{\csname BookOriginal@trHomogeneity\endcsname}}
\expandafter\def\csname BookIndex@T0408\endcsname{ziban@自伴（self adjoint）}
\expandafter\let\csname BookOriginal@trSelfAdjoint\endcsname\trSelfAdjoint
\renewcommand{\trSelfAdjoint}{\BookTerm{T0408}{\csname BookOriginal@trSelfAdjoint\endcsname}}
\expandafter\def\csname BookIndex@T0409\endcsname{bansui@伴随（adjoint）}
\expandafter\let\csname BookOriginal@trAdjoint\endcsname\trAdjoint
\renewcommand{\trAdjoint}{\BookTerm{T0409}{\csname BookOriginal@trAdjoint\endcsname}}
\expandafter\def\csname BookIndex@T0410\endcsname{zhengguisuanzi@正规算子（normal operator）}
\expandafter\let\csname BookOriginal@trNormalOperator\endcsname\trNormalOperator
\renewcommand{\trNormalOperator}{\BookTerm{T0410}{\csname BookOriginal@trNormalOperator\endcsname}}
\expandafter\def\csname BookIndex@T0411\endcsname{youjie@有界（bounded）}
\expandafter\let\csname BookOriginal@trBounded\endcsname\trBounded
\renewcommand{\trBounded}{\BookTerm{T0411}{\csname BookOriginal@trBounded\endcsname}}
\expandafter\def\csname BookIndex@T0412\endcsname{keshu@可数（countable）}
\expandafter\let\csname BookOriginal@trCountable\endcsname\trCountable
\renewcommand{\trCountable}{\BookTerm{T0412}{\csname BookOriginal@trCountable\endcsname}}
\expandafter\def\csname BookIndex@T0413\endcsname{jihusuoyou@几乎所有（almost every）}
\expandafter\let\csname BookOriginal@trAlmostEvery\endcsname\trAlmostEvery
\renewcommand{\trAlmostEvery}{\BookTerm{T0413}{\csname BookOriginal@trAlmostEvery\endcsname}}
\expandafter\def\csname BookIndex@T0414\endcsname{jihuchuchu@几乎处处（almost everywhere）}
\expandafter\let\csname BookOriginal@trAlmostEverywhere\endcsname\trAlmostEverywhere
\renewcommand{\trAlmostEverywhere}{\BookTerm{T0414}{\csname BookOriginal@trAlmostEverywhere\endcsname}}
\expandafter\def\csname BookIndex@T0415\endcsname{jinsuanzi@紧算子（compact operator）}
\expandafter\let\csname BookOriginal@trCompactOperator\endcsname\trCompactOperator
\renewcommand{\trCompactOperator}{\BookTerm{T0415}{\csname BookOriginal@trCompactOperator\endcsname}}
\expandafter\def\csname BookIndex@T0416\endcsname{he@核（kernel）}
\expandafter\let\csname BookOriginal@trKernel\endcsname\trKernel
\renewcommand{\trKernel}{\BookTerm{T0416}{\csname BookOriginal@trKernel\endcsname}}
\expandafter\def\csname BookIndex@T0417\endcsname{zhi@秩（rank）}
\expandafter\let\csname BookOriginal@trRank\endcsname\trRank
\renewcommand{\trRank}{\BookTerm{T0417}{\csname BookOriginal@trRank\endcsname}}
\expandafter\def\csname BookIndex@T0418\endcsname{fanshu@范数（norm）}
\expandafter\let\csname BookOriginal@trNorm\endcsname\trNorm
\renewcommand{\trNorm}{\BookTerm{T0418}{\csname BookOriginal@trNorm\endcsname}}
\expandafter\def\csname BookIndex@T0419\endcsname{heisaijuzhen@黑塞矩阵（Hessian）}
\expandafter\let\csname BookOriginal@trHessian\endcsname\trHessian
\renewcommand{\trHessian}{\BookTerm{T0419}{\csname BookOriginal@trHessian\endcsname}}
\expandafter\def\csname BookIndex@T0420\endcsname{yuweishu@余维数（codimension）}
\expandafter\let\csname BookOriginal@trCodimension\endcsname\trCodimension
\renewcommand{\trCodimension}{\BookTerm{T0420}{\csname BookOriginal@trCodimension\endcsname}}
\expandafter\def\csname BookIndex@T0421\endcsname{shuipingji@水平集（level set）}
\expandafter\let\csname BookOriginal@trLevelSet\endcsname\trLevelSet
\renewcommand{\trLevelSet}{\BookTerm{T0421}{\csname BookOriginal@trLevelSet\endcsname}}
\expandafter\def\csname BookIndex@T0422\endcsname{canshuhua@参数化（parametrization）}
\expandafter\let\csname BookOriginal@trParametrization\endcsname\trParametrization
\renewcommand{\trParametrization}{\BookTerm{T0422}{\csname BookOriginal@trParametrization\endcsname}}
\expandafter\def\csname BookIndex@T0423\endcsname{changshubianyigongshi@常数变易公式（variation of constants formula）}
\expandafter\let\csname BookOriginal@trVariationOfConstantsFormula\endcsname\trVariationOfConstantsFormula
\renewcommand{\trVariationOfConstantsFormula}{\BookTerm{T0423}{\csname BookOriginal@trVariationOfConstantsFormula\endcsname}}
\expandafter\def\csname BookIndex@T0424\endcsname{zhongxincanshuhua@重新参数化（reparametrization）}
\expandafter\let\csname BookOriginal@trReparametrization\endcsname\trReparametrization
\renewcommand{\trReparametrization}{\BookTerm{T0424}{\csname BookOriginal@trReparametrization\endcsname}}
\expandafter\def\csname BookIndex@T0425\endcsname{kaiji@开集（open set）}
\expandafter\let\csname BookOriginal@trOpenSet\endcsname\trOpenSet
\renewcommand{\trOpenSet}{\BookTerm{T0425}{\csname BookOriginal@trOpenSet\endcsname}}
\expandafter\def\csname BookIndex@T0426\endcsname{zhiyu@值域（range）}
\expandafter\let\csname BookOriginal@trRange\endcsname\trRange
\renewcommand{\trRange}{\BookTerm{T0426}{\csname BookOriginal@trRange\endcsname}}
\expandafter\def\csname BookIndex@T0427\endcsname{tezhengzhi@特征值（eigenvalue）}
\expandafter\let\csname BookOriginal@trEigenvalue\endcsname\trEigenvalue
\renewcommand{\trEigenvalue}{\BookTerm{T0427}{\csname BookOriginal@trEigenvalue\endcsname}}
\expandafter\def\csname BookIndex@T0428\endcsname{zhishuyingshe@指数映射（exponential map）}
\expandafter\let\csname BookOriginal@trExponentialMap\endcsname\trExponentialMap
\renewcommand{\trExponentialMap}{\BookTerm{T0428}{\csname BookOriginal@trExponentialMap\endcsname}}
\expandafter\def\csname BookIndex@T0429\endcsname{tezhengxiangliang@特征向量（eigenvector）}
\expandafter\let\csname BookOriginal@trEigenvector\endcsname\trEigenvector
\renewcommand{\trEigenvector}{\BookTerm{T0429}{\csname BookOriginal@trEigenvector\endcsname}}
\expandafter\def\csname BookIndex@T0430\endcsname{tezhengduoxiangshi@特征多项式（characteristic polynomial）}
\expandafter\let\csname BookOriginal@trCharacteristicPolynomial\endcsname\trCharacteristicPolynomial
\renewcommand{\trCharacteristicPolynomial}{\BookTerm{T0430}{\csname BookOriginal@trCharacteristicPolynomial\endcsname}}
\expandafter\def\csname BookIndex@T0431\endcsname{cankaojie@参考解（reference solution）}
\expandafter\let\csname BookOriginal@trReferenceSolution\endcsname\trReferenceSolution
\renewcommand{\trReferenceSolution}{\BookTerm{T0431}{\csname BookOriginal@trReferenceSolution\endcsname}}
\expandafter\def\csname BookIndex@T0432\endcsname{fanli@反例（counterexample）}
\expandafter\let\csname BookOriginal@trCounterexample\endcsname\trCounterexample
\renewcommand{\trCounterexample}{\BookTerm{T0432}{\csname BookOriginal@trCounterexample\endcsname}}
\expandafter\def\csname BookIndex@T0433\endcsname{guoniu@过扭（overtwisted）}
\expandafter\let\csname BookOriginal@trOvertwisted\endcsname\trOvertwisted
\renewcommand{\trOvertwisted}{\BookTerm{T0433}{\csname BookOriginal@trOvertwisted\endcsname}}
\expandafter\def\csname BookIndex@T0434\endcsname{jin@紧（tight）}
\expandafter\let\csname BookOriginal@trTight\endcsname\trTight
\renewcommand{\trTight}{\BookTerm{T0434}{\csname BookOriginal@trTight\endcsname}}
\expandafter\def\csname BookIndex@T0435\endcsname{pingmianchang@平面场（plane field）}
\expandafter\let\csname BookOriginal@trPlaneField\endcsname\trPlaneField
\renewcommand{\trPlaneField}{\BookTerm{T0435}{\csname BookOriginal@trPlaneField\endcsname}}
\expandafter\def\csname BookIndex@T0436\endcsname{bingti@柄体（handlebody）}
\expandafter\let\csname BookOriginal@trHandlebody\endcsname\trHandlebody
\renewcommand{\trHandlebody}{\BookTerm{T0436}{\csname BookOriginal@trHandlebody\endcsname}}
\expandafter\def\csname BookIndex@T0437\endcsname{shousuoyingshe@收缩映射（contraction map）}
\expandafter\let\csname BookOriginal@trContractionMap\endcsname\trContractionMap
\renewcommand{\trContractionMap}{\BookTerm{T0437}{\csname BookOriginal@trContractionMap\endcsname}}
\expandafter\def\csname BookIndex@T0438\endcsname{shensuo@伸缩（dilatation）}
\expandafter\let\csname BookOriginal@trDilatation\endcsname\trDilatation
\renewcommand{\trDilatation}{\BookTerm{T0438}{\csname BookOriginal@trDilatation\endcsname}}
\expandafter\def\csname BookIndex@T0439\endcsname{tezhengkongjian@特征空间（eigenspace）}
\expandafter\let\csname BookOriginal@trEigenspace\endcsname\trEigenspace
\renewcommand{\trEigenspace}{\BookTerm{T0439}{\csname BookOriginal@trEigenspace\endcsname}}
\expandafter\def\csname BookIndex@T0440\endcsname{zhongshu@重数（multiplicity）}
\expandafter\let\csname BookOriginal@trMultiplicity\endcsname\trMultiplicity
\renewcommand{\trMultiplicity}{\BookTerm{T0440}{\csname BookOriginal@trMultiplicity\endcsname}}
\expandafter\def\csname BookIndex@T0441\endcsname{tuoyuan@椭圆（elliptic）}
\expandafter\let\csname BookOriginal@trElliptic\endcsname\trElliptic
\renewcommand{\trElliptic}{\BookTerm{T0441}{\csname BookOriginal@trElliptic\endcsname}}
\expandafter\def\csname BookIndex@T0442\endcsname{tuoyuanguji@椭圆估计（elliptic estimate）}
\expandafter\let\csname BookOriginal@trEllipticEstimate\endcsname\trEllipticEstimate
\renewcommand{\trEllipticEstimate}{\BookTerm{T0442}{\csname BookOriginal@trEllipticEstimate\endcsname}}
\expandafter\def\csname BookIndex@T0443\endcsname{tuoyuanzhengzexing@椭圆正则性（elliptic regularity）}
\expandafter\let\csname BookOriginal@trEllipticRegularity\endcsname\trEllipticRegularity
\renewcommand{\trEllipticRegularity}{\BookTerm{T0443}{\csname BookOriginal@trEllipticRegularity\endcsname}}
\expandafter\def\csname BookIndex@T0444\endcsname{guji@估计（estimate）}
\expandafter\let\csname BookOriginal@trEstimate\endcsname\trEstimate
\renewcommand{\trEstimate}{\BookTerm{T0444}{\csname BookOriginal@trEstimate\endcsname}}
\expandafter\def\csname BookIndex@T0445\endcsname{ruojie@弱解（weak solution）}
\expandafter\let\csname BookOriginal@trWeakSolution\endcsname\trWeakSolution
\renewcommand{\trWeakSolution}{\BookTerm{T0445}{\csname BookOriginal@trWeakSolution\endcsname}}
\expandafter\def\csname BookIndex@T0446\endcsname{gudianjie@古典解（classical solution）}
\expandafter\let\csname BookOriginal@trClassicalSolution\endcsname\trClassicalSolution
\renewcommand{\trClassicalSolution}{\BookTerm{T0446}{\csname BookOriginal@trClassicalSolution\endcsname}}
\expandafter\def\csname BookIndex@T0447\endcsname{zhengzexing@正则性（regularity）}
\expandafter\let\csname BookOriginal@trRegularity\endcsname\trRegularity
\renewcommand{\trRegularity}{\BookTerm{T0447}{\csname BookOriginal@trRegularity\endcsname}}
\expandafter\def\csname BookIndex@T0448\endcsname{zhengzehua@正则化（regularization）}
\expandafter\let\csname BookOriginal@trRegularization\endcsname\trRegularization
\renewcommand{\trRegularization}{\BookTerm{T0448}{\csname BookOriginal@trRegularization\endcsname}}
\expandafter\def\csname BookIndex@T0449\endcsname{xingshimijishu@形式幂级数（formal power series）}
\expandafter\let\csname BookOriginal@trFormalPowerSeries\endcsname\trFormalPowerSeries
\renewcommand{\trFormalPowerSeries}{\BookTerm{T0449}{\csname BookOriginal@trFormalPowerSeries\endcsname}}
\expandafter\def\csname BookIndex@T0450\endcsname{xiaofenmu@小分母（small divisor）}
\expandafter\let\csname BookOriginal@trSmallDenominator\endcsname\trSmallDenominator
\renewcommand{\trSmallDenominator}{\BookTerm{T0450}{\csname BookOriginal@trSmallDenominator\endcsname}}
\expandafter\def\csname BookIndex@T0451\endcsname{guanghuaxing@光滑性（smoothness）}
\expandafter\let\csname BookOriginal@trSmoothness\endcsname\trSmoothness
\renewcommand{\trSmoothness}{\BookTerm{T0451}{\csname BookOriginal@trSmoothness\endcsname}}
\expandafter\def\csname BookIndex@T0452\endcsname{choumi@稠密（dense）}
\expandafter\let\csname BookOriginal@trDense\endcsname\trDense
\renewcommand{\trDense}{\BookTerm{T0452}{\csname BookOriginal@trDense\endcsname}}
\expandafter\def\csname BookIndex@T0453\endcsname{wuchuchoumi@无处稠密（nowhere dense）}
\expandafter\let\csname BookOriginal@trNowhereDense\endcsname\trNowhereDense
\renewcommand{\trNowhereDense}{\BookTerm{T0453}{\csname BookOriginal@trNowhereDense\endcsname}}
\expandafter\def\csname BookIndex@T0454\endcsname{bijin@逼近（approximation）}
\expandafter\let\csname BookOriginal@trApproximation\endcsname\trApproximation
\renewcommand{\trApproximation}{\BookTerm{T0454}{\csname BookOriginal@trApproximation\endcsname}}
\expandafter\def\csname BookIndex@T0455\endcsname{yinhanshudingli@隐函数定理（implicit function theorem）}
\expandafter\let\csname BookOriginal@trImplicitFunctionTheorem\endcsname\trImplicitFunctionTheorem
\renewcommand{\trImplicitFunctionTheorem}{\BookTerm{T0455}{\csname BookOriginal@trImplicitFunctionTheorem\endcsname}}
\expandafter\def\csname BookIndex@T0456\endcsname{tongyou@通有（generic）}
\expandafter\let\csname BookOriginal@trGeneric\endcsname\trGeneric
\renewcommand{\trGeneric}{\BookTerm{T0456}{\csname BookOriginal@trGeneric\endcsname}}
\expandafter\def\csname BookIndex@T0457\endcsname{shengyu@剩余（residual）}
\expandafter\let\csname BookOriginal@trResidual\endcsname\trResidual
\renewcommand{\trResidual}{\BookTerm{T0457}{\csname BookOriginal@trResidual\endcsname}}
\expandafter\def\csname BookIndex@T0458\endcsname{hengjiexing@横截性（transversality）}
\expandafter\let\csname BookOriginal@trTransversality\endcsname\trTransversality
\renewcommand{\trTransversality}{\BookTerm{T0458}{\csname BookOriginal@trTransversality\endcsname}}
\expandafter\def\csname BookIndex@T0459\endcsname{hengjie@横截（transversal）}
\expandafter\let\csname BookOriginal@trTransversal\endcsname\trTransversal
\renewcommand{\trTransversal}{\BookTerm{T0459}{\csname BookOriginal@trTransversal\endcsname}}
\expandafter\def\csname BookIndex@T0460\endcsname{hengjiemian@横截面（transversal section）}
\expandafter\let\csname BookOriginal@trTransversalSection\endcsname\trTransversalSection
\renewcommand{\trTransversalSection}{\BookTerm{T0460}{\csname BookOriginal@trTransversalSection\endcsname}}
\expandafter\def\csname BookIndex@T0461\endcsname{xiangsiyuanli@相似原理（similarity principle）}
\expandafter\let\csname BookOriginal@trSimilarityPrinciple\endcsname\trSimilarityPrinciple
\renewcommand{\trSimilarityPrinciple}{\BookTerm{T0461}{\csname BookOriginal@trSimilarityPrinciple\endcsname}}
\expandafter\def\csname BookIndex@T0462\endcsname{maopaofenxi@冒泡分析（bubbling-off analysis）}
\expandafter\let\csname BookOriginal@trBubblingOffAnalysis\endcsname\trBubblingOffAnalysis
\renewcommand{\trBubblingOffAnalysis}{\BookTerm{T0462}{\csname BookOriginal@trBubblingOffAnalysis\endcsname}}
\expandafter\def\csname BookIndex@T0463\endcsname{lianxuxing@连续性（continuity）}
\expandafter\let\csname BookOriginal@trContinuity\endcsname\trContinuity
\renewcommand{\trContinuity}{\BookTerm{T0463}{\csname BookOriginal@trContinuity\endcsname}}
\expandafter\def\csname BookIndex@T0464\endcsname{kexiliman@柯西-黎曼（Cauchy-Riemann）}
\expandafter\let\csname BookOriginal@trCauchyRiemann\endcsname\trCauchyRiemann
\renewcommand{\trCauchyRiemann}{\BookTerm{T0464}{\csname BookOriginal@trCauchyRiemann\endcsname}}
\expandafter\def\csname BookIndex@T0465\endcsname{shoulian@收敛（converge）}
\expandafter\let\csname BookOriginal@trConverge\endcsname\trConverge
\renewcommand{\trConverge}{\BookTerm{T0465}{\csname BookOriginal@trConverge\endcsname}}
\expandafter\def\csname BookIndex@T0466\endcsname{jianjin@渐近（asymptotic）}
\expandafter\let\csname BookOriginal@trAsymptotic\endcsname\trAsymptotic
\renewcommand{\trAsymptotic}{\BookTerm{T0466}{\csname BookOriginal@trAsymptotic\endcsname}}
\expandafter\def\csname BookIndex@T0467\endcsname{tonggou@同构（isomorphism）}
\expandafter\let\csname BookOriginal@trIsomorphism\endcsname\trIsomorphism
\renewcommand{\trIsomorphism}{\BookTerm{T0467}{\csname BookOriginal@trIsomorphism\endcsname}}
\expandafter\def\csname BookIndex@T0468\endcsname{tonggoude@同构的（isomorphic）}
\expandafter\let\csname BookOriginal@trIsomorphic\endcsname\trIsomorphic
\renewcommand{\trIsomorphic}{\BookTerm{T0468}{\csname BookOriginal@trIsomorphic\endcsname}}
\expandafter\def\csname BookIndex@T0469\endcsname{jieduan@截断（cutoff）}
\expandafter\let\csname BookOriginal@trCutoff\endcsname\trCutoff
\renewcommand{\trCutoff}{\BookTerm{T0469}{\csname BookOriginal@trCutoff\endcsname}}
\expandafter\def\csname BookIndex@T0470\endcsname{fangjin@仿紧（paracompact）}
\expandafter\let\csname BookOriginal@trParacompact\endcsname\trParacompact
\renewcommand{\trParacompact}{\BookTerm{T0470}{\csname BookOriginal@trParacompact\endcsname}}
\expandafter\def\csname BookIndex@T0471\endcsname{kedulianghua@可度量化（metrizable）}
\expandafter\let\csname BookOriginal@trMetrizable\endcsname\trMetrizable
\renewcommand{\trMetrizable}{\BookTerm{T0471}{\csname BookOriginal@trMetrizable\endcsname}}
\expandafter\def\csname BookIndex@T0472\endcsname{linyushousuohe@邻域收缩核（neighborhood retract）}
\expandafter\let\csname BookOriginal@trNeighborhoodRetract\endcsname\trNeighborhoodRetract
\renewcommand{\trNeighborhoodRetract}{\BookTerm{T0472}{\csname BookOriginal@trNeighborhoodRetract\endcsname}}
\expandafter\def\csname BookIndex@T0473\endcsname{tiehe@贴合（taut）}
\expandafter\let\csname BookOriginal@trTaut\endcsname\trTaut
\renewcommand{\trTaut}{\BookTerm{T0473}{\csname BookOriginal@trTaut\endcsname}}
\expandafter\def\csname BookIndex@T0474\endcsname{fanduicheng@反对称（antisymmetric）}
\expandafter\let\csname BookOriginal@trAntisymmetric\endcsname\trAntisymmetric
\renewcommand{\trAntisymmetric}{\BookTerm{T0474}{\csname BookOriginal@trAntisymmetric\endcsname}}
\makeatother

